% Amenable groups with nearly exponential sofic profile, and quantum channels that need nearly linear memory.
% Compile: python build.py (three pdflatex passes).
\documentclass[11pt,a4paper]{article}

\usepackage[utf8]{inputenc}
\usepackage[T1]{fontenc}
\usepackage{lmodern}
\usepackage{amsmath,amssymb,amsthm}
\usepackage[margin=1.1in]{geometry}
\usepackage{booktabs}
\usepackage{enumitem}
\usepackage{microtype}
\setlength{\emergencystretch}{2em}
\usepackage{xcolor}
\usepackage[colorlinks=true,linkcolor=blue!50!black,citecolor=blue!50!black,urlcolor=blue!50!black]{hyperref}
\hypersetup{pdftitle={Amenable groups with nearly exponential sofic profile, and quantum channels that need nearly linear memory},
pdfauthor={Seth Douglas and Nidhal Mghirbi},pdfkeywords={sofic profile, amenable groups, Houghton groups, Thompson's groups,
wreath products, quantum channels, quantum memory}}

\theoremstyle{plain}
\newtheorem{theorem}{Theorem}[section]
\newtheorem{proposition}[theorem]{Proposition}
\newtheorem{corollary}[theorem]{Corollary}
\newtheorem{lemma}[theorem]{Lemma}
\theoremstyle{definition}
\newtheorem{definition}[theorem]{Definition}
\theoremstyle{remark}
\newtheorem{remark}[theorem]{Remark}
\newtheorem*{theoremA}{Theorem A}
\newtheorem*{theoremB}{Theorem B}

\DeclareMathOperator{\Tr}{Tr}
\DeclareMathOperator{\Sym}{Sym}
\DeclareMathOperator{\FSym}{FSym}
\DeclareMathOperator{\GL}{GL}
\DeclareMathOperator{\AGL}{AGL}
\DeclareMathOperator{\Area}{Area}
\DeclareMathOperator{\rank}{rank}
\newcommand{\ddia}{d_\diamond}
\newcommand{\cK}{\mathcal K}
\newcommand{\F}{\mathbb F}
\newcommand{\Z}{\mathbb Z}
\newcommand{\N}{\mathbb N}

\title{Amenable groups with nearly exponential sofic profile,\\ and quantum channels that need nearly linear memory}
\author{Seth Douglas \and Nidhal Mghirbi}
\date{September 2026}

\begin{document}
\maketitle

\begin{abstract}
How far from finite can an amenable group be? Cornulier's sofic profile measures the least number of points on which a finite piece of a group
can be modelled by permutations at accuracy $1/r$. In earlier work we showed, by a counting criterion that can never give more
than $\exp(O(r))$, that Houghton's group has a finite piece with profile $\exp(\Omega(r^{0.129}))$; to our knowledge it was the first amenable
group shown to have superpolynomial profile. We construct a finitely presented elementary amenable group
with a finite piece whose profile is $\exp(r^{1+o(1)})$, with a lower bound certified by $25$ explicit relations: at least $\exp(c\,r/(\log r)^{12.75})$ and at
most $\exp(O(r\log r))$. The idea is to place copies of $S_3$ not on points but on configurations: Houghton's group moves the points of three rays,
finitely supported affine maps act on their binary configurations, and one copy of $S_3$ sits on each configuration. Then $N$ points carry $2^N$
copies, while any two copies are brought together by relations that touch only three points at a time. We reach this group through a ladder of
transport mechanisms, including an elementary amenable group at exponent $1/3$ and a subgroup of Thompson's group $V$ whose models, if they
exist, need $\exp(c\,r/\log^2r)$ points. Our group embeds in Brin's higher-dimensional Thompson group $3V$, so $3V$ has a chunk with finite
profile $\exp(r^{1+o(1)})$. Finally, the group becomes an explicit quantum channel on $\mathbb C^{873}$ whose repeated use is expensive.
A device that serves $n$ uses, releasing each output before the next input arrives, needs memory about $n/B$ when it may spend purity $B$, and
exact devices attain this trade-off up to polylogarithmic factors; with logarithmic purity it needs memory $n^{1-o(1)}$, and at the optimal
exchange rate its memory is $n^{1/2+o(1)}$, attained with purity of the same order. The engine is a one-round theorem that reads an approximate representation of the
group directly off the memory state of any such device. As a consequence, a local lattice process in $\nu$ dimensions whose bath starts
maximally mixed needs time about $\delta^{-1/\nu}$ to produce this channel to accuracy $\delta$. Autonomous dynamics can still produce such
channels: a bounded time-independent Hamiltonian, not local, coupled to a tracial bath, produces in the long run a channel with the same
memory-purity trade-off.
\end{abstract}


{\small\tableofcontents}

%======================================================================
\section{Introduction}\label{sec:intro}
%======================================================================

A group is sofic when its finite pieces can be modelled by permutations of finite sets, up to a small error. How large must those models be?
For a finite piece $E$ of a group, a \emph{chunk}, Cornulier's \emph{sofic profile} $\sigma_E(r)$ is the least number of points on which $E$ can be
modelled with multiplication correct on all but a fraction $1/r$ of the points and distinct elements separated almost everywhere~\cite{Cornulier2013}.
Amenable groups are sofic, so their profiles are finite. Cornulier asked whether a sofic group can have profile that is unbounded and not linear
in $r$~\cite[Problem~3.18]{Cornulier2013}. In~\cite{DouglasMghirbiA} we showed that Houghton's group $H_3$, which is finitely
presented~\cite{Brown1987}, with an explicit presentation due to Johnson~\cite{Johnson1999}, and elementary amenable, has a chunk whose profile is at least $\exp(\Omega(r^{c}))$, $c=1/(1+\log107)\approx0.129$. The proof rests on a counting criterion
(Theorem~\ref{thm:criterion}): if finitely many relations force many displaced copies of $S_3$ to commute cheaply, every model must be large,
because approximately commuting copies of $S_3$ carry entropy. That criterion can never give more than $\exp(O(r))$.

This paper shows that the criterion nearly reaches its ceiling in an amenable group, and that the group then becomes a quantum channel whose
repeated use needs nearly linear memory.

\begin{theoremA}
There is a finitely presented elementary amenable group $\Gamma$, with five generators, and a chunk $E$ of $\Gamma$ such that, for all $r\ge2$,
\[
\exp\Bigl(\frac{c\,r}{(\log r)^{12.75}}\Bigr)\ \le\ \sigma_E(r)\ \le\ \exp(C\,r\log r).
\]
So $\sigma_E(r)=\exp(r^{1+o(1)})$. The lower bound is certified by an explicit set of $25$ relations of $\Gamma$. Moreover, for every $\Delta\in(0,1]$ there are $c_\Delta,e_\Delta>0$ such that every assignment of unitaries in
$U(D)$ to the five generators under which the $25$ relations have defect at most $e\le e_\Delta$ in the normalized Hilbert--Schmidt norm, and
$ab$ stays at distance at least $\Delta$ from $I$, has $\log D\ge c_\Delta\,e^{-2}/(\log(1/e))^{25.49}$.
\end{theoremA}

Theorem~\ref{thm:main-group} gives the precise statement, and Proposition~\ref{prop:fp} the finite presentation, which follows from
Cornulier's criterion for permutational wreath products~\cite[Theorem~1.1]{Cornulier2006} once the affine group acting on the configurations is shown to be finitely presented. The lower bound is within a polylogarithmic factor in the exponent of the most the
criterion can give, and the upper bound shows that the truth is not much larger. The upper bound comes from models that hash configurations by
random affine maps; they are far smaller than the models that Følner sets give~\cite[Proposition~3.14]{Cornulier2013}, since the Følner function
of $\Gamma$ is doubly exponential (Proposition~\ref{prop:folner-lower}). The group $\Gamma$ also embeds in Brin's higher-dimensional Thompson group $3V$
(Proposition~\ref{prop:3V}), so $3V$, a finitely presented simple group whose soficity is open, has a chunk whose profile is finite and
$\exp(r^{1+o(1)})$.

\paragraph{Configuration lamps.} In Houghton's group the copies of $S_3$ sit on points of a ray; carrying a copy a distance $k$ takes a word of
length $k$, and commuting two copies at distance $k$ costs at least of order $k^2$ relations~\cite[Proposition~3.4]{DouglasMghirbiA}. A family of
$N$ copies spreads over distance $N$, and its commutators are expensive. We place the copies instead on the finitely
supported binary configurations of the three rays of $H_3$. Houghton's group moves the points, finitely supported affine maps act on the
configurations (a translation flips one point, a transvection adds one point into another), and one copy of $S_3$ sits on each configuration. The
$2^N$ configurations supported on the first $N$ points of a ray are reached by words of length $O(N^2)$, and two of them differ by a configuration
that at most $N$ transvections clear down to a single point. Each step touches at most three points, so its cost is governed by the Dehn function
of $H_3$ at scale $N$, which is polynomial by~\cite[Theorem~3.3]{DouglasMghirbiA}. The total cost of commuting two of the $2^N$ copies is therefore
polynomial in $N$: polylogarithmic in the number of copies. This is what the criterion needs.

\paragraph{A ladder.} The construction completes a sequence of transport mechanisms, each of which the criterion turns into a profile bound.

\begin{center}
\begin{tabular}{llll}
\toprule
transport & group & $\log\sigma_E(r)\ \gtrsim$ & profile finite \\
\midrule
along a ray & $H_3$~\cite{DouglasMghirbiA} & $r^{0.129}$ & yes \\
by one element & dyadic $\mathcal G$ (\S\ref{sec:dyadic}) & $r^{1/3}$ & yes \\
exponential orbits & $\mathcal L\le V$ (\S\ref{sec:thompson}) & $r/\log^2r$ & if $F$ is amenable \\
configurations & $\Gamma$ (\S\S\ref{sec:group}--\ref{sec:profile}) & $r/(\log r)^{12.75}$ & yes \\
\bottomrule
\end{tabular}
\end{center}

\noindent Transport by the powers of one element gives exponent $1/(p+1)$ when the far commutators have area of order $k^p$; along a ray of
$H_3$ the area is at least quadratic~\cite[Proposition~3.4]{DouglasMghirbiA}, so that transport stops at $1/3$, and $\mathcal G$ attains $1/3$ with
seventeen relations.
Thompson's group $F$ transports cheaply along exponentially growing orbits, which gives nearly exponential profile in a finitely presented
subgroup $\mathcal L$ of Thompson's group $V$, provided its models exist at all: they do if $F$ is amenable, and we show that $\mathcal L$ is
hyperlinear exactly when $F$ is, so that unitary models of this chunk exist whenever $F$ is hyperlinear
(Proposition~\ref{prop:thompson-hyperlinear}). The configuration lamps keep the exponential transport and drop the condition. For a variant on the
configurations of a grid (Remark~\ref{rem:grid}) a sketched argument gives a stronger lower bound, $r/(\log r)^{7.38}$, with the same upper bound.

\paragraph{Quantum memory.} Kretschmann and Werner showed that every causal quantum process is a concatenation of memory
channels~\cite{KretschmannWerner2005}, and Rybár and Ziman asked when a single memory device, which is not reset between uses, applies the same
channel to each of its inputs~\cite{RybarZiman2008}. In~\cite{DouglasMghirbiB} we asked how much memory such a device needs to apply the same quantum channel $n$ times, when
each output must be released before the next input arrives and every retained qubit counts. The device may export qubits and import fresh ones; the \emph{exchange}
$J$ counts the qubits exchanged, and the \emph{purity}, measured as the number of qubits minus the entropy of the initial memory and of the imports,
is a resource. The error is the half trace distance to the ideal experiment, against adaptive observers that keep quantum references. For a
gadget channel in the tracial closure, $\frac12\log r_*$ qubits per use, with $r_*$ the Choi rank, is the optimal exchange
rate~\cite[Propositions~3.1 and~9.1]{DouglasMghirbiB}. A channel built from a finite set of relations inherits the approximation profile of the
group: small memory would give small approximate representations. From $H_3$ we obtained an explicit channel needing memory $(n/\log n)^{0.069}$. The
configuration lamps give nearly linear memory, and a trade-off between memory and purity that is sharp up to polylogarithmic factors.

\begin{theoremB}
The certificate of Theorem~A defines an explicit quantum channel $\Phi$ on $\mathbb C^{873}$, in the closure of channels with finite tracial
factorizations. For causal services of $n$ uses of $\Phi$ with error at most $1/16$ and large $n$:
\begin{enumerate}[label=(\alph*),itemsep=1pt]
\item with purity at most $B$, where $\log n\le B\le c_0n$, the least memory lies between $c\,n/(B(\log n)^{25.49})$, at any exchange, and
$K(n/B)\log n$, attained by exact services that refresh their memory every round, with exchange $(n-1)Q$; so the least memory at purity $B$,
times $B$, is $n^{1+o(1)}$, and purity $O(\log n)$ forces memory $n^{1-o(1)}$;
\item at the optimal exchange rate, that is, with exchange $J\le\frac n2\log r_*$ over $n$ uses, the least memory lies between $c\sqrt n/(\log n)^{12.75}$, whatever the purity, and $K\sqrt{n\log n}$; services
near the least memory spend purity $n^{1/2\pm o(1)}$, at least that much by (a) and at most $3Q+6$, so at that rate memory and purity are both
about $\sqrt n$.
\end{enumerate}
\end{theoremB}

Theorem~\ref{thm:main-channel} gives the precise statement. Its engine is a one-round theorem for every gadget channel whose relations contain
those of $S_3$ and a name for $ab$ (Theorem~\ref{thm:weighted}): a single round of any device reveals an approximate representation of the group on the device's memory, with a
relator defect of the order of the square root of the leakage and entropy production of that round. In~\cite[Theorem~5.5 and
Appendix~B]{DouglasMghirbiB} this was proved for one gadget, built from Houghton's group. The upper bounds come from sofic models, which give devices directly (Proposition~\ref{prop:sofic-upper}): at most points of a model the
pure bath state reproduces the channel exactly, so a memory that is maximally mixed on those points serves the channel with no error at all. So
the memory of $\Phi$ is caught between two readings of the profile of $\Gamma$ at scale $n$, and for $\Gamma$ both are $n^{1+o(1)}$. At the
optimal exchange rate a device can spend at most $3Q+6$ bits of purity (Proposition~\ref{prop:purity-budget}), which gives the lower bound in
part~(b). The upper bound needs a device that keeps a sofic model in memory for all $n$ rounds instead of refreshing it: a streaming theorem
(Theorem~\ref{thm:streaming}) uses the good points of a model as an exact dilation whose bath grows slightly in every round, so that the bad points
cost purity rather than accumulated error.

\paragraph{Physics.} A local lattice system whose bath starts in its maximally mixed state, run for time $T$, produces channels that are cheap
for fixed $T$: they need memory only polynomial in
$T$ and $\log n$ (Proposition~\ref{prop:light-cone}). So a local process that produces $\Phi$ to accuracy $\delta$ takes time at least about
$\delta^{-1/\nu}$ in lattice dimension $\nu$ (Corollary~\ref{cor:preparation}). Autonomous dynamics with a bath in its tracial state face a different constraint: by the mean ergodic
theorem, their long-time limits are compressions of conditional expectations and hence self-adjoint
channels~\cite[Remark~10.3]{JungeLeMerdyXu2006},~\cite[Remark~5.2]{HaagerupMusat2011}, and they satisfy a further inequality between coherence and return (Propositions~\ref{prop:autonomous}
and~\ref{prop:return}). Every channel $\Psi\Psi^*$ with $\Psi$ in the tracial closure is such a limit (Proposition~\ref{prop:nonlocal}), but for gadget channels this route gives nothing
expensive: composing a gadget channel with its adjoint erases the multiplication of the group, and $\Phi\Phi^*$ has an exact finite factorization
on at most $3r_*-2$ qubits (Proposition~\ref{prop:compositions}). Averaging the doubling of $\Phi$ with the identity does give an expensive
limit: a bounded time-independent Hamiltonian coupled to a tracial bath produces a channel on $M_{1746}$ exactly at the times $\frac\pi2k$, and in
the limit of long times, that trades memory and purity at the same rate as $\Phi$ (Proposition~\ref{prop:lazy}). The Hamiltonian is not local;
whether local autonomous dynamics can produce an expensive channel is open.

\paragraph{Dependencies on the companion papers.} Table~\ref{tab:dependencies} lists every result of~\cite{DouglasMghirbiA}
and~\cite{DouglasMghirbiB} used here, with the numbering of their version 1.1.

\begin{table}[ht]
\centering\small
\begin{tabular}{@{}p{0.47\textwidth}p{0.2\textwidth}p{0.25\textwidth}@{}}
\toprule
result & source & used in\\
\midrule
counting criterion & \cite[Thm.~8.1]{DouglasMghirbiA} & Thm.~\ref{thm:criterion}, every lower bound\\
transport by the powers of one element & \cite[Cor.~8.2]{DouglasMghirbiA} & Section~\ref{sec:criterion}\\
Johnson's presentation of $H_3$ and its conventions & \cite[Section~2]{DouglasMghirbiA} & Section~\ref{sec:houghton}\\
far-commutator area and polynomial Dehn bound for $H_3$ & \cite[Thms.~3.2, 3.3]{DouglasMghirbiA} & Prop.~\ref{prop:dehn}\\
quadratic lower bound along a ray & \cite[Prop.~3.4]{DouglasMghirbiA} & Sections~\ref{sec:intro}, \ref{sec:criterion}, \ref{sec:ladder}\\
bounded or at least $2r$ (sharpening Cornulier's Fact~3.5); embeddings in finite groups & \cite[Lemma~2.1]{DouglasMghirbiA} & Section~\ref{sec:ladder}, after Thm.~\ref{thm:main-group}\\
superpolynomial profile of $H_3$ & \cite[Thm.~1.1]{DouglasMghirbiA} & Section~\ref{sec:thompson}\\
projection products & \cite[Lemma~4.1]{DouglasMghirbiA} & Lemma~\ref{lem:projection-product}\\
\midrule
causal services, exchange, purity and error & \cite[Def.~2.3, Section~3]{DouglasMghirbiB} & Section~\ref{sec:channel}\\
gadget channels: Choi rank $r_*$, flat probe, entropy $\log r_*$ & \cite[Prop.~9.1]{DouglasMghirbiB} & Section~\ref{sec:channel-setting}\\
the rank rate is the optimal exchange rate of gadget channels in $\cK_d$ & \cite[Props.~3.1, 9.1]{DouglasMghirbiB} & Thm.~B(b)\\
sequential extraction & \cite[Thm.~5.2]{DouglasMghirbiB} & Thm.~\ref{thm:extraction}\\
purity budget $14Q+3$ at the rank rate (sharpened here) & \cite[Thm.~5.8]{DouglasMghirbiB} & Prop.~\ref{prop:purity-budget}\\
entropy variance & \cite[Lemma~5.9]{DouglasMghirbiB} & Lemma~\ref{lem:surprisal}\\
rolling-reservoir compiler & \cite[Appendix~A, Lemmas~A.1, A.4]{DouglasMghirbiB} & Thm.~\ref{thm:streaming}\\
channel profiles are group profiles; hyperlinear groups give channels in $\cK_d$ & \cite[Thm.~9.2]{DouglasMghirbiB} & Section~\ref{sec:channel}\\
one-round theorem for the Houghton gadget (generalised here) & \cite[Thm.~5.5, Appendix~B]{DouglasMghirbiB} & Thm.~\ref{thm:weighted}, Section~\ref{sec:weighted-proof}\\
weighted norms; charged transport and weighted extraction for the Houghton gadget (generalised here) & \cite[Appendix~B, Lemmas~B.2, B.4]{DouglasMghirbiB} & Lemmas~\ref{lem:weighted-extraction}, \ref{lem:charged-area}\\
weighted $S_3$ rounding & \cite[Lemma~B.3]{DouglasMghirbiB} & Lemma~\ref{lem:rounding}\\
weighted sector bound (proved here by sequential measurement) & \cite[Thm.~B.1]{DouglasMghirbiB} & Thm.~\ref{thm:sector}\\
\bottomrule
\end{tabular}
\caption{Results of the companion papers used in this paper (numbering of their version 1.1).}
\label{tab:dependencies}
\end{table}

\paragraph{Organization.} Section~\ref{sec:criterion} recalls the counting criterion. Section~\ref{sec:ladder} treats the dyadic group and the
subgroup of Thompson's group $V$. Sections~\ref{sec:group}--\ref{sec:profile} construct the configuration lamps, show that they are finitely presented and
embed in $3V$, prove the local fillings and the pair-area theorem, and prove Theorem~A. Section~\ref{sec:channel} proves Theorem~B, with the one-round theorem proved in
Section~\ref{sec:weighted-proof}. Section~\ref{sec:physics} treats local physics, and Section~\ref{sec:open} lists open problems. The certificate of
$25$ relations is in the accompanying files, together with programs that check it; logarithms are to base two unless written $\ln$.

%======================================================================
\section{Counting approximately commuting copies of \texorpdfstring{$S_3$}{S3}}\label{sec:criterion}
%======================================================================

Every lower bound in this paper comes from one principle, proved in~\cite{DouglasMghirbiA}. If finitely many relations force many displaced
copies of $S_3$ to commute cheaply, every model of a fixed chunk must be large, because approximately commuting copies of $S_3$ carry
entropy. We fix the conventions used in both halves of the paper, state the principle in the form used here, with $a,b$ in place of the
letters $u,v$ of~\cite{DouglasMghirbiA}, and note the ceiling $\exp(O(r))$ that it can never pass. The scheme, many cheaply transported
copies of a finite group, relations whose cost is controlled by area, and a dimension bound, follows Slofstra~\cite[Sections~2--3]{Slofstra2018}.
The criterion needs only pairwise approximate commutation of displaced copies of $S_3$.

\paragraph{Conventions.} Functions are composed from right to left, groups act on the left, $[x,y]=xyx^{-1}y^{-1}$, and $\log$ is to base
two. For a finite set $S$ and a finite set $R$ of words in $S^{\pm1}$, $\Area_R(w)$ is the least number of conjugates of words of $R^{\pm1}$
whose product is $w$ in the free group on $S$, and $\infty$ if there is none. A \emph{chunk} of a group is a finite subset $E\ni1$ with the
partial multiplication $xy=z$ whenever $x,y,z\in E$ and $xy=z$ in the group. Its \emph{sofic profile} $\sigma_E(r)$ is the least $n$ for which
some $f:E\to\Sym_n$ with $f(1)=\mathrm{id}$ multiplies and separates correctly up to $1/r$: $d(f(x)f(y),f(xy))\le1/r$ whenever $xy$ is defined in
$E$, and $d(f(x),f(y))\ge1-1/r$ for distinct $x,y\in E$~\cite{Cornulier2013}. Here $d$ is the normalized Hamming distance, the fraction of points
on which two permutations differ. The normalized Hilbert--Schmidt norm of a $D\times D$ matrix is $\|X\|_2=(\Tr(X^*X)/D)^{1/2}$.

The criterion is Theorem~8.1 of~\cite{DouglasMghirbiA}, and Appendix~\ref{app:imported} reproduces its proof.

\begin{theorem}[Counting criterion]\label{thm:criterion}
Let $G$ be a group, $S\subset G$ a finite set containing $a,b$, and $R$ a finite set of words in $S^{\pm1}$, each representing $1$ in $G$, with
$a^2,b^2,(ab)^3\in R$ and $ab\ne1$ in $G$. Let $l$ be the largest length of a word in $R$, and let $E$ be the chunk consisting of $1$, $S^{\pm1}$,
the prefixes of the words of $R$ and of $ab$, and the inverses of these elements. Let $w_1,\dots,w_N$ be words in $S^{\pm1}$ and $\Lambda\ge1$ such
that
\[
\Area_R\bigl([w_ixw_i^{-1},w_jyw_j^{-1}]\bigr)\le\Lambda\qquad(i\ne j,\ x,y\in\{a,b\}).
\]
Then $\log\sigma_E(r)\ge N/16$ for every $r\ge10^8N\Lambda(2l-1)$. Moreover, for $0<\Delta\le1$ and
$K_\Delta=\sqrt{80000\,(1+\log(1/\Delta))}/\Delta$, every assignment $\phi$ of unitaries in $U(D)$ to $S$ under which every word of $R$ lies
within $e$ of $I$ in the normalized Hilbert--Schmidt norm and $\|\phi(ab)-I\|_2\ge\Delta$ satisfies $\log D\ge N\Delta^2/12$ whenever
$K_\Delta\sqrt N\,\Lambda e\le1$.
\end{theorem}

The proof in~\cite{DouglasMghirbiA} has four steps, and only the last is specific to $S_3$. Conjugating the images of $a$ and $b$ by the image
of $w_i$ gives $N$ copies of $S_3$, one for each word. Each copy satisfies the relations of $S_3$ up to the defect of a relator, because those
relations are conjugates of $a^2$, $b^2$ and $(ab)^3$, so it can be rounded to an exact representation of $S_3$. Its three-cycle is a conjugate
of the image of $ab$, so it is displaced from the identity as much as that image is. Two copies commute up to $\Lambda$ times the relator defect,
because each of their commutators is a product of at most $\Lambda$ conjugates of relators. Finally, exact copies of $S_3$ that nearly commute
and have displaced three-cycles cannot share a small space: they force a fixed amount of entropy per copy, at least $\Delta^2/12$ bits for
unitaries and $1/16$ bit for permutations~\cite[Section~4]{DouglasMghirbiA}. The chunk $E$ does not depend on $N$ or on the words $w_i$; only
the areas matter.

The theorem has a ceiling. Its hypotheses force $N\le r/10^8$ and $N\le(K_\Delta e)^{-2}$, so it never gives $\log\sigma_E(r)$ beyond order $r$,
nor $\log D$ beyond order $e^{-2}$. Below the ceiling, everything depends on how cheaply the group carries copies of $S_3$ to many places.
Transport by the powers of one element, with area of order $k^p$ at distance $k$, gives exponent $1/(p+1)$; this is Corollary~8.2
of~\cite{DouglasMghirbiA}.

\begin{corollary}[Transport by the powers of one element]\label{cor:cyclic}
Let $G$, $S$, $R$, $l$ and $E$ be as in Theorem~\ref{thm:criterion}, let $t\in S$, and suppose that
$\Area_R([x,t^kyt^{-k}])\le Ck^p$ for $x,y\in\{a,b\}$ and $k\ge1$, with $C\ge1$ and $p\ge0$. Then
\[
\log\sigma_E(r)\ \ge\ \frac1{16}\Big\lfloor\Big(\frac r{10^8C(2l-1)}\Big)^{1/(p+1)}\Big\rfloor\qquad(r>1),
\]
and every unitary assignment as in Theorem~\ref{thm:criterion} has $\log D\ge\frac{\Delta^2}{12}\bigl((CK_\Delta e)^{-2/(2p+1)}-1\bigr)$.
\end{corollary}

In Houghton's group $H_3$ the copies of $S_3$ along a ray have commutators of at least quadratic area~\cite[Proposition~3.4]{DouglasMghirbiA},
so transport along a ray never passes the exponent $1/3$. A group that does better must reach many sites cheaply.

One more fact of Cornulier's is used twice, in Sections~\ref{sec:more} and~\ref{sec:ladder}. The profile of a chunk is bounded exactly when the
chunk embeds in a finite group, and otherwise it is at least $r$~\cite[Fact~3.5]{Cornulier2013}, indeed at least
$2r$~\cite[Lemma~2.1]{DouglasMghirbiA}. A chunk with unbounded profile therefore certifies that its group is not locally embeddable into finite
groups, not LEF~\cite[Fact~3.8]{Cornulier2013}.

%======================================================================
\section{Two steps up the ladder}\label{sec:ladder}
%======================================================================

Before the main construction we climb two steps that show what the criterion needs. The first keeps transport by the powers of one element
and makes it cheap: in an elementary amenable group the far commutators have at most quadratic area, which gives the exponent $1/3$.
Transport by the powers of one element gives exponent $1/(p+1)$ when the far commutators have area of order $k^p$, and along a ray of $H_3$ the
area is at least quadratic~\cite[Proposition~3.4]{DouglasMghirbiA}, so that transport stops at $1/3$. The second changes the kind of transport: Thompson's group $F$ has orbits that grow
exponentially, so words of length $O(\log N)$ reach $N$ sites, and its quadratic Dehn function fills all pairwise commutators cheaply. This comes
within a factor $\log^2r$ in the exponent of the ceiling of Theorem~\ref{thm:criterion}, but in a group that contains $F$, whose soficity is unknown.
Section~\ref{sec:group} keeps the exponential transport and removes the condition.

\subsection{An amenable group at exponent \texorpdfstring{$1/3$}{1/3}}\label{sec:dyadic}

In this subsection and the next, functions are composed from right to left, and a group acting on a set $X$ acts on
$\bigoplus_{x\in X}S_3$ by carrying the copy of $S_3$ at $x$ to the copy at $g(x)$ by the identity of $S_3$.

Let $X=\mathbb Z[\frac12]$, and let $t$ and $u$ be the permutations of $X$ given by $t(x)=x+1$ for $x\in\mathbb Z$, $t(x)=x$ for
$x\notin\mathbb Z$, and $u(x)=2x$. Let $\mathcal B=\langle t,u\rangle$ and $\mathcal G=\bigl(\bigoplus_{x\in X}S_3\bigr)\rtimes\mathcal B$, and let $a,b$ be
generating involutions of the copy of $S_3$ at $0$. Put $s=utu^{-1}$ and $v=tst^{-1}$. Then $s$ adds $2$ to the even integers and $v$ adds $2$
to the odd integers, each fixing every other point, so $t^2=sv$ and $[s,v]=1$. Moreover $u$ and $v$ fix $0$, and $s$ fixes $1=t(0)$. Since $\mathcal B$ acts transitively on $X$, the group $\mathcal G$ is generated by $a$,
$b$, $t$ and $u$. So the
following seventeen words in $a,b,t,u,s,v$, which make up the set $R_{\mathcal G}$, are relations of $\mathcal G$:
\[
a^2,\ b^2,\ (ab)^3,\quad s^{-1}utu^{-1},\ v^{-1}tst^{-1},\ t^2v^{-1}s^{-1},\ [s,v],\quad [u,x],\ [v,x],\ [s,txt^{-1}],\ [x,tyt^{-1}],
\]
where $x,y\in\{a,b\}$.

\begin{theorem}\label{thm:dyadic}
The group $\mathcal G$ is elementary amenable and not solvable, and every commutator $[x,t^kyt^{-k}]$ with $x,y\in\{a,b\}$ and
$k\ge1$ has area at most $5k^2$ with respect to $R_{\mathcal G}$. Consequently the chunk $E_{\mathcal G}$ that Theorem~\ref{thm:criterion}
associates with $R_{\mathcal G}$ has finite sofic profile with
\[
\log\sigma_{E_{\mathcal G}}(r)\ \ge\ \frac1{16}\Big\lfloor\Big(\frac r{75\cdot10^8}\Big)^{1/3}\Big\rfloor\qquad(r>1),
\]
every unitary assignment $\phi$ as in Theorem~\ref{thm:criterion} has $\log D\ge\frac{\Delta^2}{12}\bigl((5K_\Delta e)^{-2/5}-1\bigr)$, and
$\mathcal G$ is not LEF.
\end{theorem}

\begin{proof}
\emph{Amenability.} For $j\in\mathbb Z$ let $t_j=u^jtu^{-j}$; it adds $2^j$ to the points of $2^j\mathbb Z$ and fixes the others. The group $\mathcal K$
generated by all $t_j$ is normalized by $u$ and contains $t$, so it is normal in $\mathcal B$ and $\mathcal B/\mathcal K$ is cyclic. For $n\ge1$ the elements $t_j$ with
$|j|\le n$ move only points of $2^{-n}\mathbb Z$, and they commute with the translation $x\mapsto x+2^n$ of $2^{-n}\mathbb Z$, which has $4^n$
orbits. A permutation of $2^{-n}\mathbb Z$ that commutes with this translation maps each orbit onto an orbit by a translation, so these $t_j$ generate a subgroup of
$\mathbb Z^{4^n}\rtimes\Sym_{4^n}$, which is virtually abelian. Hence $\mathcal K$ is a directed union of virtually abelian groups, $\mathcal B$ is elementary
amenable, and so is $\mathcal G$, because $\bigoplus_XS_3$ is locally finite.

\emph{Not solvable.} For $n\ge1$ the elements $t_0,\dots,t_{n-1}$ preserve $\mathbb Z$ and commute with $x\mapsto x+2^n$, so they act on
$\mathbb Z/2^n$, where they permute the classes modulo $2^i$ for every $i\le n$. So they generate a subgroup $\Gamma_n$ of the iterated wreath
product $\mathsf W_n$ of $n$ copies of $\mathbb Z/2$, which has order $2^{2^n-1}$. On the even classes, identified with $\mathbb Z/2^{n-1}$ by
halving, $t_1,\dots,t_{n-1}$ act as $t_0,\dots,t_{n-2}$ act on $\mathbb Z/2^{n-1}$, and they fix the odd classes; conjugation by $t_0$ carries
this action to the odd classes, and $t_0$ exchanges the two halves. So $|\Gamma_n|\ge2|\Gamma_{n-1}|^2$, and $\Gamma_n=\mathsf W_n$ by induction from $\Gamma_1=\mathsf W_1=\mathbb Z/2$. The derived
subgroup of $\mathsf W_n$ lies in $\mathsf W_{n-1}\times\mathsf W_{n-1}$ and projects onto each factor, since the commutator of $(g,1)$ with the
exchange is $(g^{-1},g)$; by induction the $(n-1)$st derived subgroup of $\mathsf W_n$ is nontrivial. Since $\Gamma_n$ is a quotient of the subgroup
$\langle t_0,\dots,t_{n-1}\rangle$ of $\mathcal B$, the derived length of $\mathcal B$ is at least $n$ for every $n$, so neither $\mathcal B$ nor
$\mathcal G$ is solvable.

\emph{Area.} Let $A_{\mathcal G}(k)$ be the largest area of $[x,t^kyt^{-k}]$ over $x,y\in\{a,b\}$, so that $A_{\mathcal G}(1)=1$. Each of the following moves costs one
cell: replacing $t^2$ by $sv$; exchanging adjacent letters $s^{\pm1}$ and $v^{\pm1}$; replacing $s$ by $utu^{-1}$; moving a letter $x^{\pm1}$,
$x\in\{a,b\}$, across $u^{\pm1}$ or $v^{\pm1}$; and moving $txt^{-1}$ or its inverse across $s^{\pm1}$. The commutator contains four powers
$t^{\pm k}$.

For $k=2m$, write each $t^{\pm2m}$ as $(s^mv^m)^{\pm1}$, with $m$ replacements and $m(m-1)/2$ exchanges each, $2m^2+2m$ cells in all. Moving
$y^{\pm1}$ across $v^{\pm m}$ ($2m$ cells) leaves $[x,s^mys^{-m}]$. Replacing its $4m$ letters $s^{\pm1}$ gives
$[x,ut^mu^{-1}yut^{-m}u^{-1}]$, and moving $y^{\pm1}$ and $x^{\pm1}$ across $u$ (four cells) gives $u[x,t^myt^{-m}]u^{-1}$. Hence
$A_{\mathcal G}(2m)\le A_{\mathcal G}(m)+2m^2+8m+4$.

For $k=2m+1$, put $z=tyt^{-1}$, so that the commutator is freely $[x,t^{2m}zt^{-2m}]$. Write each $t^{\pm2m}$ as $(v^ms^m)^{\pm1}$, now with
$m(m+1)/2$ exchanges each, $2m^2+6m$ cells in all. Moving $z^{\pm1}$ across $s^{\pm m}$ ($2m$ cells) leaves $[x,v^mzv^{-m}]$, moving
$x^{\pm1}$ across $v^{\pm m}$ ($2m$ cells) leaves $v^m[x,z]v^{-m}$, and $[x,z]$ is a relation. Hence $A_{\mathcal G}(2m+1)\le2m^2+10m+1$.

By induction $A_{\mathcal G}(k)\le5k^2$: $A_{\mathcal G}(2)\le15$, $A_{\mathcal G}(m)+2m^2+8m+4\le7m^2+8m+4\le20m^2$ for $m\ge1$, and $2m^2+10m+1\le5(2m+1)^2$.

\emph{Profile.} The words of $R_{\mathcal G}$ have length at most $8$, so Corollary~\ref{cor:cyclic} with $C=5$, $p=2$ and $l=8$ gives the
bounds. Since $\mathcal G$ is amenable, it is sofic and every chunk has finite profile. The profile of $E_{\mathcal G}$ is unbounded, so
$E_{\mathcal G}$ does not embed in a finite group (\cite[Lemma~2.1]{DouglasMghirbiA}), and $\mathcal G$ is not LEF.
\end{proof}

The fillings in the proof have been generated and checked by free reduction for $k\le64$; the largest ratio of cell count to $k^2$ among
them is $15/4$, at $k=2$. In $H_3$ the doubling endomorphism multiplies a reweighted area by $107$ at each step, which gives the exponent $\log107$; in
$\mathcal G$ the dilation $u$ relates distance $2m$ to distance $m$ exactly, at a cost of $O(m^2)$ cells for the two commuting shifts. The orbit of $0$ under $\mathcal B$
grows exponentially, as the orbits of Thompson's group $F$ do, but we do not know whether $\mathcal G$ supports the transport of many
copies used for $F$ below (Section~\ref{sec:open}).

\subsection{Thompson's group \texorpdfstring{$V$}{V}}\label{sec:thompson}

Genevois proved that every Houghton group embeds in Thompson's group $V$~\cite{Genevois2019}. Since $H_3$ is amenable, $V$ therefore has a
chunk whose profile is finite at every $r$ and superpolynomial by~\cite[Theorem~1.1]{DouglasMghirbiA}. A different subgroup of $V$ gives a much larger bound.

Let $\mathbb D$ be the set of dyadic rationals in $(0,1)$, and let Thompson's group $F$ act on $[0,1]$ by its piecewise-linear
homeomorphisms. Let $\mathcal L=\bigl(\bigoplus_{x\in\mathbb D}S_3\bigr)\rtimes F$.

\begin{lemma}[Gap coding]\label{lem:gap-coding}
$\mathcal L$ embeds in $V$.
\end{lemma}

\begin{proof}
Let $\mathsf c$ be the substitution $0\mapsto0$, $1\mapsto11$ on finite binary words. An infinite binary word either is a concatenation of
blocks $0$ and $11$, or it has a first occurrence of $10$ at a block boundary. So Cantor space is the disjoint union of the set $\mathcal C$ of
block sequences and the cylinders $Z_w=\mathsf c(w)10\{0,1\}^{\mathbb N}$, one for each finite word $w$. Index $Z_w$ by the dyadic
$x_w=0.w1\in\mathbb D$ in binary; this is a bijection. Let $f\in F$ have a tree-pair diagram with leaves $w_1,\dots,w_n$ and $w'_1,\dots,w'_n$,
so that $f(0.w_kz)=0.w'_kz$. Define $\hat f$ by the prefix replacements $\mathsf c(w_k)\mapsto\mathsf c(w'_k)$ and
$\mathsf c(v)10\mapsto\mathsf c(v')10$, where the interior vertices $v$ of the first tree are matched in order with the interior vertices $v'$ of
the second. Both families of prefixes are complete prefix codes, because splitting a leaf $w$ replaces $\mathsf c(w)$ by $\mathsf c(w0)$,
$\mathsf c(w)10$ and $\mathsf c(w1)$; expanding a matched pair of leaves adds a matched pair of gaps and changes nothing else, so $\hat f$ depends
only on $f$. It maps $Z_w$ onto $Z_{w'}$ with $x_{w'}=f(x_w)$: below a leaf this is the formula for $f$, and interior vertices correspond to the
breakpoints between leaves. On $\mathcal C$ it acts as $f$ acts on binary sequences, so $f\mapsto\hat f$ is an injective homomorphism $F\to V$.
Inside each $Z_w$ let $S_3$ permute the three cylinders $\mathsf c(w)10y\{0,1\}^{\mathbb N}$, $y\in\{0,10,11\}$. Copies with different $w$ have
disjoint supports, and $\hat f$ carries the copy at $x$ to the copy at $f(x)$. An element of the group they generate that acts trivially acts
trivially on $\mathcal C$, so its $F$-part is trivial, and then each of its $S_3$ coordinates is trivial.
\end{proof}

Let $a,b$ generate the copy of $S_3$ at $\frac12$, let $t\in F$ be
\[
t(x)=\begin{cases}2x,&0\le x\le\frac14,\\ x+\frac14,&\frac14\le x\le\frac12,\\ \frac{x+1}2,&\frac12\le x\le1,\end{cases}
\]
so that $t(\frac12)=\frac34$, and for $m\ge1$ let $\mathbb D_m$ be the set of the $2^m-1$ points of $\mathbb D$ with denominator at most
$2^m$.

\begin{proposition}\label{prop:thompson}
The group $\mathcal L$ is finitely presented, contains $F$, and contains no Houghton group $H_m$, so the bounds below do not come from the
Houghton subgroups of $V$. There are a finite set $R$ of relations of $\mathcal L$, containing $a^2,b^2,(ab)^3$, and a constant $C$ with the
following property: for every $m\ge1$ there are words $w_p$, $p\in\mathbb D_m$, with $w_p(\frac12)=p$ and
\[
\mathrm{Area}_R\bigl([w_pxw_p^{-1},w_qyw_q^{-1}]\bigr)\le Cm^2\qquad(p\ne q\in\mathbb D_m,\ x,y\in\{a,b\}).
\]
Consequently a chunk $E$ of $\mathcal L$, and hence of $V$, has $\log\sigma_E(r)\ge c\,r/\log^2r$ for all large $r$, and for $0<\Delta\le1$
every unitary assignment as in Theorem~\ref{thm:criterion} has $\log D\ge c_\Delta\,e^{-2}/\log^4(1/e)$ for all small $e$, where $c>0$ is a
constant and $c_\Delta>0$ depends only on $\Delta$.
\end{proposition}

\begin{proof}
\emph{Finite presentation.} Thompson's group $F$ is finitely presented, and it acts transitively on increasing pairs of points of
$\mathbb D$~\cite[Lemma~4.2]{CFP1996}, so it has finitely many orbits on $\mathbb D\times\mathbb D$. The stabilizer $F_{1/2}$ of $\frac12$ is the
product of the elements supported in $[0,\frac12]$ and in $[\frac12,1]$, two copies of $F$, so it is finitely generated. By Cornulier's
criterion for permutational wreath products, $\mathcal L$ is finitely presented; this action of $F$ is~\cite[Example~3.4]{Cornulier2006}. The bounds below do not use this.

\emph{Transport.} A standard dyadic interval is one of the form $[k2^{-j},(k+1)2^{-j}]$. If two partitions of $[0,1]$ into standard dyadic
intervals have the same number $n$ of intervals, the map sending the $i$th interval of the first affinely onto the $i$th interval of the second
lies in $F$ and has a tree-pair diagram with $n-1$ carets, and every standard dyadic interval can be split into any given number of standard
dyadic intervals. For $p<q$ in $\mathbb D_m$, the binary digits of $p$ split $[0,p]$ and $[p,1]$ into at most $m$ standard dyadic
intervals each, and $[p,q]$, cut at its point $c$ with the shortest binary expansion, splits into at most $m$ standard dyadic intervals
on each side of $c$. Hence there are $w_p\in F$ with $w_p(\frac12)=p$ and fewer than $2m$ carets, and $g_{pq}\in F$ with $g_{pq}(\frac12)=p$, $g_{pq}(\frac34)=q$ and fewer than $4m$
carets. The word length in $F$ is comparable to the number of carets of the reduced tree-pair diagram~\cite{BCS2001}, so these elements
are represented by words of length $O(m)$ in the generating set of the finite presentation of $F$ used in $R$ below; the Dehn functions of
two finite presentations of $F$ are equivalent, so Guba's quadratic bound holds for it. We write $w_p$ also for such a word. An element fixing
$\frac12$ maps $[0,\frac12]$ to itself, so the left and right subtrees of its reduced diagram are diagrams for its two restrictions, which are
elements of $F$ after affine rescaling, with fewer carets in total. Hence $F_{1/2}\cong F\times F$ is undistorted in $F$: its elements of
length $n$ in $F$ are words of length $O(n)$ in a fixed finite generating set $h_1,\dots,h_q$ of $F_{1/2}$.

\emph{Relations.} Let $R$ consist of a finite presentation of $F$ on a generating set that includes $t$ and the $h_j$, the words
$a^2,b^2,(ab)^3$, the commutators $[x,h_j]$, and the commutators $[x,tyt^{-1}]$, for $x,y\in\{a,b\}$. They hold in $\mathcal L$: each $h_j$
fixes $\frac12$, and the copies at $\frac12$ and at $t(\frac12)=\frac34$ commute.

\emph{Area.} Let $p<q$ in $\mathbb D_m$ and $g=g_{pq}$. The elements $g^{-1}w_p$ and $t^{-1}g^{-1}w_q$ fix $\frac12$ and have length
$O(m)$, so they are represented by words $\omega,\omega'$ of length $O(m)$ in the $h_j$. The words $w_p^{-1}g\omega$ and $w_q^{-1}gt\omega'$ are null
in $F$
and have length $O(m)$, so Guba's quadratic Dehn function~\cite{Guba2006} fills each with $O(m^2)$ cells. Replacing the four
occurrences of $w_p^{\pm1}$ in $[w_pxw_p^{-1},w_qyw_q^{-1}]$ by $(g\omega)^{\pm1}$, and the four of $w_q^{\pm1}$ by $(gt\omega')^{\pm1}$, therefore
costs
$O(m^2)$ cells. Moving $x^{\pm1}$ across $\omega$ and $y^{\pm1}$ across $\omega'$ with the relations $[x,h_j]$ costs $O(m)$ cells and leaves
$g[x,tyt^{-1}]g^{-1}$, a conjugate of a relation. For $p>q$ the commutator is the inverse of the one with the two copies exchanged. So every
area is $O(m^2)$.

\emph{Profile.} Apply Theorem~\ref{thm:criterion} with $N=2^m-1$ and $\Lambda=Cm^2$, where $C\ge1$, and let $l$ be the largest
length of a word of $R$. Given $r$, let $m$ be the largest
integer with $2^mm^2\le r/(10^8C(2l-1))$; then $2^m$ is at least of order $r/\log^2r$, and $\log\sigma_E(r)\ge(2^m-1)/16$. Given
$e$, let $m$ be the largest integer with $2^{m/2}m^2\le1/(CK_\Delta e)$; then $2^m$ is at least of order
$(CK_\Delta)^{-2}e^{-2}/\log^4(1/e)$, and $\log D\ge(2^m-1)\Delta^2/12$.

\emph{No Houghton subgroup.} A torsion element of $\mathcal L$ maps to a torsion element of the torsion-free group $F$, so it lies in
$\bigoplus_{\mathbb D}S_3$, which has exponent six. Every $H_m$ contains a five-cycle. Finally, $\mathcal L$ embeds in $V$ by
Lemma~\ref{lem:gap-coding}.
\end{proof}

The subgroup $\mathcal L$ contains $F$, whose soficity is open, so we do not know whether this chunk has permutation models at all; if it
does, they are very large. In particular, either $V$ is not sofic, or its models of this chunk need $2^{c\,r/\log^2r}$ points, close to the
largest bound that Theorem~\ref{thm:criterion} can give. If $F$ is amenable, so is $\mathcal L$, and then this chunk has finite profile. If $F$ is
hyperlinear, so is $\mathcal L$ (Proposition~\ref{prop:thompson-hyperlinear}), and then unitary models of the chunk exist at every defect. For
unitary models of Slofstra's group, his bound~\cite{Slofstra2018} gives $\log D\gtrsim e^{-c}$ for every $c<\frac12$, and, as Slofstra
notes~\cite[Corollary~1.3 and its proof]{Slofstra2018}, it gives permutation models with $\log\sigma\gtrsim r^c$ for every $c<\frac14$; the group $\mathcal L$ gives $e^{-2}$ and $r$ up to logarithms. The
group $\mathcal L$ has the form $\bigl(\bigoplus_{\Gamma/\Gamma_0}A\bigr)\rtimes\Gamma$ of the non-sofic wreath products of Kun and Thom~\cite{KunThom2026},
which are built from groups with property (T); here the input is instead cheap transport of non-abelian copies of $S_3$.

Soficity and hyperlinearity pass to wreath products whose acting group is sofic~\cite{HayesSale2018}, and to permutational wreath products
whose acting group is sofic (respectively hyperlinear) and whose action is sofic~\cite[Theorems~3.6 and~3.8]{GKEP2025},~\cite{AlekseevBradford2026}.
Since $F$ acts on $\mathbb D$ transitively and faithfully, a sofic action of $F$ on $\mathbb D$ would already make $F$
sofic~\cite[Theorem~2.12]{GKEP2025}, so these results give nothing here; the next proposition assumes only that $F$ is hyperlinear.

\begin{proposition}\label{prop:thompson-hyperlinear}
The group $\mathcal L$ is hyperlinear if and only if $F$ is.
\end{proposition}

\begin{proof}
Subgroups of hyperlinear groups are hyperlinear, so suppose that $F$ is hyperlinear. A countable group $\Gamma$ is hyperlinear if and only if
its von Neumann algebra $L(\Gamma)$ has a trace-preserving embedding into an ultrapower $\mathcal R^\omega$ of the hyperfinite
$\mathrm{II}_1$ factor~\cite{CapraroLupini2015}, and then every separable von Neumann subalgebra of $L(\Gamma)^\omega$ embeds in
$\mathcal R^\omega$ as well, by a diagonal argument. It therefore suffices to embed $L(\mathcal L)$ into $L(F)^\omega$ preserving the trace,
for a free ultrafilter $\omega$ on $\mathbb N$. We write $\tau$ for the canonical traces and $F_J$ for the subgroup of elements supported in an
interval $J$. If $J,J'$ are disjoint, $L(F_J)$ and $L(F_{J'})$ commute and $\tau(xy)=\tau(x)\tau(y)$ for $x\in L(F_J)$ and $y\in L(F_{J'})$,
since $F_J$ and $F_{J'}$ generate their direct product. If $f,f'\in F$ agree on $J$, then $\lambda(f)x\lambda(f)^{-1}=\lambda(f')x\lambda(f')^{-1}$
for $x\in L(F_J)$, because $f'^{-1}f$ fixes $J$ pointwise and therefore commutes with $F_J$.

\emph{A localized copy of $S_3$.} Choose $\sigma\in F$ supported in $[\frac12,\frac34]$, with $\sigma(x)<x$ on $(\frac12,\frac34)$ and
$\sigma(x)=\frac12+\frac12(x-\frac12)$ near $\frac12$, and a dyadic $x_0\in(\frac12,\frac34)$. The intervals
$J_k=\sigma^k((\sigma(x_0),x_0))$, $k\in\mathbb Z$, are disjoint and accumulate at $\frac12$ as $k\to\infty$. For a nonempty open interval $J$ with
dyadic endpoints, $F_J\cong F$ has infinite conjugacy classes, since a nontrivial element has conjugates with infinitely many different
supports, so $L(F_J)$ is a $\mathrm{II}_1$ factor. In $L(F_{J_0})$ choose a unital copy of $M_6$ and, in its relative commutant, which is again a
$\mathrm{II}_1$ factor, a self-adjoint $Z_0$ with uniform distribution on $[0,1]$. Let $v_0:S_3\to U(M_6)$ be the regular representation, so
that $\tau(v_0(g))=0$ for $g\ne1$, and put $v_k=\lambda(\sigma)^kv_0\lambda(\sigma)^{-k}$ and $Z_k=\lambda(\sigma)^kZ_0\lambda(\sigma)^{-k}$,
which lie in $L(F_{J_k})$. By the factorization of traces, the $Z_k$ are independent and uniform. For a finite $K\subset\mathbb Z$ and $k\in K$,
let $p^K_k$ be the spectral projection of the event that $Z_k$ is the largest of the $Z_j$, $j\in K$. Ties have trace zero, so these projections
are orthogonal and sum to $1$, and they commute with every $v_j(g)$. Hence
\[
V_K(g)=\sum_{k\in K}p^K_kv_k(g)
\]
is a representation of $S_3$. Since $v_k(g)$ lies in a copy of $M_6$ whose relative commutant contains $Z_k$, and the other $Z_j$ lie in
algebras with factorized trace, $\tau(p^K_kv_k(g))=\tau(p^K_k)\tau(v_0(g))$; so $\tau(V_K(g))=0$ for $g\ne1$. If $K$ is an interval of $n$ integers and $|q|\le n$,
the two selected indices agree unless the largest $Z_j$ over $K\cup(K+q)$ has its index outside $K\cap(K+q)$, an event of trace
$2|q|/(n+|q|)$, and on its complement $V_K(g)=V_{K+q}(g)$. Hence $\|V_{K+q}(g)-V_K(g)\|_2^2\le8|q|/(n+|q|)$.

\emph{Invariance.} Put $V_n=V_{\{n,\dots,2n-1\}}$, and let $h\in F_{1/2}$. Near $\frac12$ on the right, $h(x)=\frac12+2^m(x-\frac12)$ for an
integer $m$, so $h$ agrees with $\sigma^{-m}$ on $J_k$ for all large $k$. Conjugation by $\lambda(h)$ therefore carries $v_k$ and $Z_k$ to
$v_{k-m}$ and $Z_{k-m}$ for large $k$, and $V_n$ to $V_{\{n,\dots,2n-1\}-m}$ for large $n$, so $\|\lambda(h)V_n(g)\lambda(h)^{-1}-V_n(g)\|_2\to0$.
Thus $V(g)=(V_n(g))_\omega$ is a representation of $S_3$ in $L(F)^\omega$ with $\tau(V(g))=0$ for $g\ne1$, and it commutes with $\lambda(F_{1/2})$.

\emph{The embedding.} For $x\in\mathbb D$ choose $f_x\in F$ with $f_x(\frac12)=x$ and put $V_x=\lambda(f_x)V\lambda(f_x)^{-1}$. By the
invariance this does not depend on the choice of $f_x$, and $\lambda(f)V_x\lambda(f)^{-1}=V_{f(x)}$ for $f\in F$. The representatives $V_n(g)$ lie
in $L(F_{I_n})$, where $I_n=\bigcup_{n\le k<2n}J_k$ shrinks to $\frac12$ from the right, so those of $V_x$ lie in $L(F_{f_x(I_n)})$, and for $x\ne y$
the intervals $f_x(I_n)$ and $f_y(I_n)$ are disjoint for large $n$. So copies at distinct points commute, and the $V_x$ together with $\lambda$
define a homomorphism $\pi:\mathcal L\to U(L(F)^\omega)$. Let $\gamma=\bigl(\prod_{x\in X}\ell_x\bigr)f$, with $X$ finite, $\ell_x$ in the copy of $S_3$
at $x$ and $f\in F$. If $f=1$ and some $\ell_x\ne1$, the factors lie in commuting algebras with factorized trace for large $n$, so
$\tau(\pi(\gamma))=\prod_x\tau(V(\ell_x))=0$. If $f\ne1$, choose a point moved by $f$ outside $X$. For large $n$ it lies outside the closure of
every $f_x(I_n)$, so $f$ is not in the group $K_n$ generated by the $F_{f_x(I_n)}$, and $\tau(a\lambda(f))=0$ for every $a\in L(K_n)$, in
particular for the representatives of $\prod_x\pi(\ell_x)$. Hence $\tau(\pi(\gamma))=0$ whenever $\gamma\ne1$, and $\pi$ extends to a
trace-preserving embedding of $L(\mathcal L)$ into $L(F)^\omega$.
\end{proof}

\subsection{From orbits to configurations}\label{sec:toward}

The two steps isolate what is needed: exponentially many sites reachable by short words, pairwise commutators of polylogarithmic area, and an
amenable ambient group. In $\mathcal G$ the orbit of $0$ grows exponentially, but we do not know how to commute two copies at points of
opposite parity cheaply. In $\mathcal L$ everything works except amenability, which is the open question of whether $F$ is amenable. The
configuration lamps of the next section obtain all three at once by a change of viewpoint: the sites are finite binary configurations of the points of Houghton's three rays
rather than points, the transport is by affine maps that touch a bounded number of coordinates at a time, and the ambient group is locally
finite-by-$\Z^2$.

%======================================================================
\section{Configuration lamps}\label{sec:group}
%======================================================================

Theorem~\ref{thm:criterion} turns $M$ copies of $S_3$ whose pairwise commutators have area at most $\Lambda$ into $\log\sigma_E(r)\ge M/16$ at
$r$ of order $M\Lambda$, so it approaches its ceiling $\exp(O(r))$ only when $\Lambda$ is tiny compared with $M$. The group of this section places
the copies of $S_3$ not on the points of Houghton's three rays but on their finitely supported binary configurations: Houghton's group moves the
points, finitely supported affine maps move the configurations, and one copy of $S_3$ sits on each configuration. Then $N$ points carry $2^N$
configurations, and hence $2^N$ copies of $S_3$, while every local move we use involves at most three points. Two of the $2^N$ copies are
compared by about $N^2$ local moves, each filled with the Dehn function of $H_3$ at scale $N$, so the area is polylogarithmic in the number of
copies.

Throughout, functions are composed from right to left, groups act on the left, $[x,y]=xyx^{-1}y^{-1}$, and $\log$ is to base two.

\paragraph{Houghton's group.} Let $Y=\{1,2,3\}\times\{1,2,\dots\}$; the point $(i,j)$ lies on ray $i$ at height $j$. Houghton's group
$H=H_3$ consists of the permutations $g$ of $Y$ that are eventually translations on each ray: there are integers $m_1(g),m_2(g),m_3(g)$ with
$g(i,j)=(i,j+m_i(g))$ for all large $j$. The map $g\mapsto(-m_2(g),-m_3(g))$ is a surjective homomorphism $H\to\Z^2$ whose kernel is the group
$\FSym(Y)$ of finitary permutations~\cite{Houghton1978}. Let $p$ be the permutation with $p(1,j)=(1,j+1)$, $p(2,1)=(1,1)$ and
$p(2,j)=(2,j-1)$ for $j\ge2$, fixing ray~$3$, and let $q$ be defined in the same way with rays $2$ and $3$ exchanged. They generate $H$
(Section~\ref{sec:houghton}), and their images in $\Z^2$ are $(1,0)$ and $(0,1)$.

\paragraph{The group.} Let $V=\F_2^{(Y)}$ be the space of finitely supported functions $Y\to\F_2$, the \emph{configurations}, with basis
$(e_y)_{y\in Y}$. The group $H$ acts on $V$ by permutation matrices, $g(e_y)=e_{g(y)}$, and we identify $H$ with its image in $\GL(V)$. Let
$\GL_{\rm fin}(Y)$ be the group of linear automorphisms of $V$ that fix all but finitely many basis vectors. It is the directed union of the finite
groups $\GL(\F_2^J)$, $J\subset Y$ finite, each extended by the identity on the other basis vectors; it is normalized by $H$; and
$H\cap\GL_{\rm fin}(Y)=\FSym(Y)$. Let $L=\GL_{\rm fin}(Y)\,H\le\GL(V)$. Write $\tau_u$ for the translation $v\mapsto v+u$ of $V$, let $A=V\rtimes L$
be the group of affine maps $v\mapsto g(v)+u$ with $g\in L$ and $u\in V$, and let
\[
\Gamma=\Bigl(\bigoplus_{v\in V}S_3\Bigr)\rtimes A,
\]
where $A$ permutes the summands through its action on $V$: writing $\lambda_v$ for the copy at $v$ of $\lambda\in S_3$, we have
$g\lambda_vg^{-1}=\lambda_{g(v)}$ for $g\in A$. Equivalently, $\Gamma$ is the group of permutations of $V\times S_3$ generated by the maps
$(v,s)\mapsto(g(v),s)$ for $g\in A$ and $(v,s)\mapsto(v,\lambda(v)s)$ for finitely supported $\lambda:V\to S_3$; this action is faithful. Groups obtained from a permutation action by adjoining the finitary linear or symmetric groups of the permuted
set, and lamps over it, appear as enrichments~\cite{Bradford2022,AlekseevBradford2026} and as halo products~\cite{GenevoisTessera2024}.

For distinct points $y,y'\in Y$ let $E_{y\leftarrow y'}$ be the transvection that sends $e_{y'}$ to $e_{y'}+e_y$ and fixes the other basis vectors.
For $g\in H$ and a linear map $M$ of $V$,
\begin{equation}\label{eq:affine-identities}
g\tau_{e_y}g^{-1}=\tau_{e_{g(y)}},\qquad gE_{y\leftarrow y'}g^{-1}=E_{g(y)\leftarrow g(y')},\qquad M\tau_uM^{-1}=\tau_{M(u)} .
\end{equation}
Linear maps fix $0$, so they commute with the copy of $S_3$ at $0$. On the first ray we abbreviate: for $x\ge1$ we write $x$ for the point $(1,x)$,
$e_x$ for $e_{(1,x)}$ and $\tau_x$ for $\tau_{e_x}$.

Let $a$ and $b$ be the transpositions $(1\,2)$ and $(2\,3)$ in the copy of $S_3$ at $0$, so that $ab$ has order three, and put
\[
t=\tau_1,\qquad d=E_{2\leftarrow1},\qquad c=d\,ab ;
\]
thus $t$ adds the point $1$ to a configuration, and $d$ adds the point $2$ to every configuration that contains the point $1$.

\begin{lemma}\label{lem:structure}
\begin{enumerate}[label=(\alph*),itemsep=1pt]
\item The element $c$ has order six, $c^3=d$ and $c^4=ab$; hence $b=ac^4$.
\item $\Gamma$ is generated by the five elements $p,q,t,c,a$.
\item Let $\pi:\Gamma\to\Z^2$ be the composite of $\Gamma\to A\to L\to L/\GL_{\rm fin}(Y)\cong H/\FSym(Y)\cong\Z^2$. Its kernel
$\bigl(\bigoplus_VS_3\bigr)\rtimes\bigl(V\rtimes\GL_{\rm fin}(Y)\bigr)$ is locally finite. Hence $\Gamma$ is locally finite-by-$\Z^2$, and in
particular elementary amenable.
\end{enumerate}
\end{lemma}

\begin{proof}
(a) The linear map $d$ fixes $0$, so it commutes with $ab$; since $d$ has order two and $ab$ order three, $c$ has order six, $c^3=d^3(ab)^3=d$ and
$c^4=d^4(ab)^4=ab$.

(b) Let $\Gamma'$ be the subgroup generated by the five elements. It contains $H$, because $H=\langle p,q\rangle$. Since $\FSym(Y)$ acts
$2$-transitively on $Y$, \eqref{eq:affine-identities} shows that $\Gamma'$ contains every $\tau_{e_y}$, hence $V$, and, since $d=c^3$, every
transvection $E_{y\leftarrow y'}$. Transvections generate every $\GL(\F_2^J)$, so $\GL_{\rm fin}(Y)\le\Gamma'$, and $A\le\Gamma'$. Finally
$b=ac^4\in\Gamma'$, and the copy of $S_3$ at $u\in V$ is $\tau_u\langle a,b\rangle\tau_u^{-1}$.

(c) Since $H\cap\GL_{\rm fin}(Y)=\FSym(Y)$, we have $L/\GL_{\rm fin}(Y)\cong H/\FSym(Y)$, and the kernel of $\pi$ is as stated. It is an
extension of $\bigoplus_VS_3$ by $V\rtimes\GL_{\rm fin}(Y)$, the directed union of the finite groups $\F_2^J\rtimes\GL(\F_2^J)$; both are locally
finite. An extension of a locally finite group $K$ by a locally finite group is locally finite: a finitely generated subgroup $\Lambda$ has finite
image modulo $K$, so $\Lambda\cap K$ has finite index in $\Lambda$, is finitely generated, and is finite. Locally finite groups are directed
unions of finite groups, and elementary amenable groups are closed under directed unions and extensions~\cite{Chou1980}.
\end{proof}

\subsection{A finite presentation}\label{sec:fp}

The certificate of Section~\ref{sec:certificate} is a finite set of relations and does not present $\Gamma$: the group it presents has
abelianization $\Z^2\times\Z/2\times\Z/4$, in which $c$ has order four (a computation with the exponent sums; for instance $a^2$,
$b^2=(ac^4)^2$ and $(ab)^3$ give $2a=8c=12c=0$), while in $\Gamma$ the element $c=d\,ab$ is a product of commutators, since
$d=[E_{2\leftarrow3},E_{3\leftarrow1}]$ and $ab$ is a three-cycle. The group itself is finitely presented.

\begin{proposition}[Finite presentation]\label{prop:fp}
The groups $L$, $A$ and $\Gamma$ are finitely presented.
\end{proposition}

\begin{proof}
\emph{Stabilizers.} For a finite set $D\subset Y$ the pointwise stabilizer $H_{(D)}$ of $D$ in $H$ is isomorphic to $H$: renumbering the points of each
ray outside $D$ in increasing order turns $Y\setminus D$ into three rays, and a permutation fixing $D$ is eventually a translation on the old rays
exactly when it is on the new ones. So $H_{(D)}$ is finitely generated. Since $H$ contains $\FSym(Y)$, it acts transitively on ordered $k$-tuples of
distinct points for every $k$.

\emph{A presentation lemma.} Let a finitely presented group $H$ act on a set $\mathcal X$ with finitely many orbits and finitely generated
stabilizers, and let $F(\mathcal X)$ be the free group on $\mathcal X$. Then $F(\mathcal X)\rtimes H$, with $H$ permuting the free generators, is
finitely presented: adjoin to a presentation of $H$ one representative $x_i$ of each orbit and the relations $[x_i,k]=1$ for $k$ in a finite
generating set of its stabilizer; then $hx_ih^{-1}$ depends only on $h\cdot x_i$, and the two groups are isomorphic (both are the iterated amalgam of $H$ with the groups
$\mathrm{Stab}(x_i)\times\langle x_i\rangle$ over the stabilizers $\mathrm{Stab}(x_i)$). Imposing, in addition, a set of
relations that is a finite union of $H$-orbits keeps the group finitely presented, since one representative of each orbit suffices.

\emph{A presentation of $A$.} For distinct $i,j\in Y$ let $x_{ij}$ be the transvection $e_j\mapsto e_j+e_i$, let $t_i=\tau_{e_i}$, and let
$\varsigma_{ij}\in\FSym(Y)\le H$ be the transposition of $i$ and $j$. Let $\widetilde A$ be the quotient of $F(\mathcal X)\rtimes H$, where
$\mathcal X=\{x_{ij}\}\sqcup\{t_i\}$ and $H$ relabels indices, by the relations
\begin{gather*}
x_{ij}^2=1,\qquad [x_{ij},x_{kl}]=1\ (i\ne l,\ j\ne k),\qquad [x_{ij},x_{jk}]=x_{ik}\ (i,j,k\text{ distinct}),\\
\varsigma_{ij}=x_{ij}x_{ji}x_{ij},\qquad t_i^2=1,\qquad [t_i,t_j]=1,\\
x_{ij}t_jx_{ij}^{-1}=t_it_j,\qquad x_{ij}t_kx_{ij}^{-1}=t_k\ (k\ne j).
\end{gather*}
They hold in $A$, so there is a surjection $\widetilde A\to A$. Each relation involves at most four points, so by transitivity the relations form
finitely many $H$-orbits; the stabilizers of $x_{ij}$ and $t_i$ are $H_{(\{i,j\})}$ and $H_{(\{i\})}$. Since $H$ is finitely presented~\eqref{eq:johnson}, $\widetilde A$ is finitely presented, and it
remains to show that $\widetilde A\to A$ is injective. The linear relations are Steinberg's~\cite{Milnor1971}; we give the short direct argument.

Fix a finite set $D\subset Y$, number its points $1,\dots,m$ (a numbering unrelated to heights), and let $U_D\le\widetilde A$ be generated by the $x_{ij}$ with $i<j$ in $D$. By the first three
relations, $x_{1m},\dots,x_{m-1,m}$ are commuting involutions, and they generate a subgroup normalized by the $x_{ij}$ with $i<j<m$; by induction
$|U_D|\le2^{m-1}2^{(m-1)(m-2)/2}=2^{m(m-1)/2}$, the order of the upper unitriangular group $\mathrm{UT}_m(\F_2)$, onto which $U_D$ maps. So
$U_D\cong\mathrm{UT}_m(\F_2)$. Let $W_D=\Sym(D)\le H$ and $\varsigma_i=\varsigma_{i,i+1}$. By the relation $\varsigma_{ij}=x_{ij}x_{ji}x_{ij}$ and relabeling, the
subgroup $G_D$ generated by all $x_{ij}$ with $i,j\in D$ contains $W_D$ and is generated by $U_D$ and the $\varsigma_i$. We claim $G_D=U_DW_DU_D$; since this
set contains $1$, it suffices that it is closed under right multiplication by $U_D$ and by each $\varsigma_i$. Every $b\in U_D$ is $x_{i,i+1}^\epsilon b_0$ with
$b_0$ in the subgroup generated by the other positive roots, the kernel of the entry $(i,i+1)$, which $\varsigma_i$ normalizes. So it suffices
to treat $wx_{i,i+1}\varsigma_i$ with $w\in W_D$; put $k=w(i)$ and $l=w(i+1)$. If $k<l$, then $wx_{i,i+1}\varsigma_i=x_{kl}w\varsigma_i\in U_DW_D$. If
$k>l$, then $w\varsigma_i=\varsigma_{kl}w$ and $x_{kl}\varsigma_{kl}=x_{lk}x_{kl}$, so $wx_{i,i+1}\varsigma_i=x_{kl}\varsigma_{kl}w=x_{lk}wx_{i,i+1}\in U_DW_DU_D$. The map $G_D\to\GL(\F_2^D)$ is onto, since transvections generate $\GL(\F_2^D)$. If $b_1wb_2$ maps to the identity of $\GL(\F_2^D)$, the permutation matrix of $w$
is upper unitriangular, so $w=1$ in $H$, and then $b_1b_2=1$ in $U_D$. Hence $G_D\cong\GL(\F_2^D)$. Every element of the subgroup $\widetilde J$
generated by all $x_{ij}$ lies in some $G_D$, so $\widetilde J\cong\GL_{\rm fin}(Y)$. By relabeling, every element of the subgroup generated by
$\widetilde J$ and $H$ has the form $jh$; if it maps to $1$, then $h\in\GL_{\rm fin}(Y)\cap H=\FSym(Y)$, so $h$ is a product of transpositions, each
equal in $\widetilde A$ to a word in the $x_{ij}$, and $jh\in\widetilde J$ maps to $1$, hence is $1$. Finally the $t_i$ generate a normal elementary
abelian subgroup, which maps isomorphically onto $V$ (it is a quotient of $\F_2^{(Y)}$, and $t_i\mapsto e_i$), and the linear part acts on it as on $V$, so every element of $\widetilde A$ is $t\ell$ with $t$ in that subgroup; if its image is
the identity affine map, evaluating at $0$ gives $t=1$, and then $\ell=1$. So $\widetilde A\cong A$. Omitting the $t_i$ gives $L$.

\emph{The lamps.} The group $A$ acts on $V$ doubly transitively: a translation carries $v$ to $0$, and a finitary linear map carries any nonzero
vector to $e_1$. The stabilizer of $0$ is $L$, which is generated by $p$, $q$ and $d$, because the $H$-conjugates of $d$ are all the transvections
$E_{y\leftarrow y'}$. So $A$ is finitely presented, its point stabilizers are finitely generated, and it has two orbits on $V\times V$. By de
Cornulier's criterion for permutational wreath products~\cite{Cornulier2006}, $\Gamma=S_3\wr_VA$ is finitely presented. Explicitly, if $\mathcal R_A$
presents $A$ on generators including $p,q,d,t$, then adding the generators $a,b$ and the relations $a^2$, $b^2$, $(ab)^3$, $[a,\ell]$ and $[b,\ell]$ for
$\ell\in\{p,q,d\}$, and $[u,tvt^{-1}]$ for $u,v\in\{a,b\}$, presents $\Gamma$: the lamp at $v\in V$ is well defined as a conjugate of
$\langle a,b\rangle$ because $L$ centralizes it, and double transitivity carries the commutation of the lamps at $0$ and $e_1$ to every pair.
\end{proof}

\subsection{An embedding in a Brin--Thompson group}\label{sec:3V}

Let $\mathfrak C=\{0,1\}^{\N}$ be Cantor space. For $n\ge1$, Brin's group $nV$~\cite{Brin2004} consists of the homeomorphisms $h$ of
$\mathfrak C^n$ given by a finite table of prefix replacements: there are two partitions of $\mathfrak C^n$ into the same number of rectangles
$u_1\mathfrak C\times\dots\times u_n\mathfrak C$, with $u_1,\dots,u_n$ finite binary words, matched in pairs, and $h$ maps each rectangle of the first
partition onto its partner by replacing the $n$ prefixes and keeping the $n$ suffixes. So $1V$ is Thompson's group $V$, and $nV$ embeds in $n'V$ for
$n\le n'$ by acting trivially on the additional coordinates. The groups $nV$ are finitely presented~\cite{HennigMatucci2012} and
simple~\cite{Brin2010}, and, as for $V$, it is not known whether they are sofic. (The letter $V$ in these names is not the space of configurations.)

\begin{proposition}\label{prop:3V}
$\Gamma$ embeds in $3V$, and hence in $nV$ for every $n\ge3$. Consequently $3V$ has a chunk $E_{3V}$ with
\[
\exp\Bigl(\frac{c\,r}{(\log r)^{12.75}}\Bigr)\ \le\ \sigma_{E_{3V}}(r)\ \le\ \exp(C\,r\log r)\qquad(r\ge2).
\]
\end{proposition}

\begin{proof}
\emph{Coding configurations.} For a finite set $S=\{j_1<\dots<j_m\}$ of positive integers, with $j_0=0$, let
\[
e(S)=\Bigl(\prod_{k=1}^m1\,0^{\,j_k-j_{k-1}-1}\,1\Bigr)\,0,
\]
so that $e(\emptyset)=0$, $e(\{1\})=110$ and $e(\{2\})=1010$. Read a word from the left: at a record boundary the letter $0$ ends the word and the
letter $1$ opens a record $10^k1$, which ends at its next letter $1$. So the words $e(S)$ form a prefix code. A finite word whose reading has not
ended can be extended to one that begins with some $e(S)$, by appending $0$ at a record boundary or $10$ inside a record. Hence the cylinders
$e(S)\mathfrak C$ are disjoint and their union is dense in $\mathfrak C$. For $v\in V$ and $i=1,2,3$ let $S_i(v)=\{j:v(i,j)=1\}$, let
\[
Z_v=e(S_1(v))\mathfrak C\times e(S_2(v))\mathfrak C\times e(S_3(v))\mathfrak C ,
\]
and let $\iota_v:\mathfrak C^3\to Z_v$ prefix the three codewords. The rectangles $Z_v$ are disjoint, and their union is dense in $\mathfrak C^3$.

\emph{Prepending a site.} Define $s_0,s_1:\mathfrak C\to\mathfrak C$ by $s_0(0\xi)=0\xi$, $s_0(1\xi)=10\xi$ and $s_1(\xi)=11\xi$. Their images
$0\mathfrak C\cup10\mathfrak C$ and $11\mathfrak C$ partition $\mathfrak C$, and, with $S+1=\{j+1:j\in S\}$,
\begin{equation}\label{eq:prepend}
s_0\bigl(e(S)\xi\bigr)=e(S+1)\,\xi,\qquad s_1\bigl(e(S)\xi\bigr)=e\bigl(\{1\}\cup(S+1)\bigr)\,\xi :
\end{equation}
an empty site in front lengthens the first gap by one, and an occupied site in front adds the record $11$.

\emph{The affine group.} For $b\in\{0,1\}$ put
\begin{gather*}
\hat P(x,s_b(y),z)=(s_b(x),y,z),\qquad \hat Q(x,y,s_b(z))=(s_b(x),y,z),\\
\hat T(s_b(x),y,z)=(s_{1-b}(x),y,z),
\end{gather*}
and let $\hat D$ fix $s_0(\mathfrak C)\times\mathfrak C^2$ pointwise and send $(11x,y,z)$ to $(11x',y,z)$, where $\hat T(x,y,z)=(x',y,z)$. These maps lie
in $3V$. With $\epsilon$ the empty word, $\hat P$ replaces the prefixes $(\epsilon,11,\epsilon)$, $(0,0,\epsilon)$, $(1,0,\epsilon)$, $(0,10,\epsilon)$ and
$(1,10,\epsilon)$ by $(11,\epsilon,\epsilon)$, $(0,0,\epsilon)$, $(10,0,\epsilon)$, $(0,1,\epsilon)$ and $(10,1,\epsilon)$; $\hat Q$ does the same with the roles
of the second and third coordinates exchanged; $\hat T$ exchanges the first-coordinate prefixes $0\leftrightarrow110$ and $10\leftrightarrow111$; and $\hat D$
exchanges $110\leftrightarrow11110$ and $1110\leftrightarrow11111$ and keeps $0$ and $10$. The permutation $p$ moves the first point of ray~$2$ to
the first point of ray~$1$ and shifts the other points of these rays, $q$ does the same with ray~$3$, $t$ changes the value of a configuration at
the point $1$, and $d$ adds the value at $1$ to the value at $2$. So~\eqref{eq:prepend} gives, for every $v\in V$,
\[
\hat P\iota_v=\iota_{p(v)},\qquad \hat Q\iota_v=\iota_{q(v)},\qquad \hat T\iota_v=\iota_{t(v)},\qquad \hat D\iota_v=\iota_{d(v)} .
\]
By the proof of Lemma~\ref{lem:structure}(b), $p,q,t,d$ generate $A$. If a word in them represents $g\in A$, the same word in $\hat P,\hat Q,\hat T,\hat D$ maps
$\iota_v(\xi)$ to $\iota_{g(v)}(\xi)$ for all $v$ and $\xi$; if $g=1$, it is the identity on the dense set $\bigcup_vZ_v$, hence the identity. So
$p,q,t,d\mapsto\hat P,\hat Q,\hat T,\hat D$ defines a homomorphism $\rho:A\to3V$ with
\begin{equation}\label{eq:3V-equivariant}
\rho(g)\,\iota_v=\iota_{g(v)}\qquad(g\in A,\ v\in V),
\end{equation}
and $\rho$ is injective: if $g\ne1$, then $g(v)\ne v$ for some $v$, and $\rho(g)$ maps $Z_v$ onto the disjoint rectangle $Z_{g(v)}$.

\emph{The lamps.} Let $\beta_1=0$, $\beta_2=10$ and $\beta_3=11$, and let $S_3$ act faithfully on $\mathfrak C^3$ by the prefix replacements
$\theta(\lambda)(\beta_i\xi,y,z)=(\beta_{\lambda(i)}\xi,y,z)$. For $\lambda\in S_3$ and $v\in V$ let $\rho(\lambda_v)$ be $\iota_v\theta(\lambda)\iota_v^{-1}$ on
$Z_v$ and the identity off $Z_v$; it lies in $3V$, since the complement of a rectangle is a finite union of rectangles. Copies at different
configurations have disjoint supports and commute, each copy is faithful, and by~\eqref{eq:3V-equivariant}
$\rho(g)\rho(\lambda_v)\rho(g)^{-1}=\rho(\lambda_{g(v)})$ for $g\in A$. These are the defining relations of $\Gamma$ as a semidirect product, so $\rho$
extends to a homomorphism $\rho:\Gamma\to3V$. Suppose that $\rho(\gamma)=1$, where $\gamma$ has lamp part $\lambda:V\to S_3$ and affine part $g$. The
lamps preserve every $Z_v$, so $\rho(g)$ does too, and $g(v)=v$ for all $v$; so $g=1$. Then $\rho(\gamma)$ restricts on $Z_v$ to
$\iota_v\theta(\lambda(v))\iota_v^{-1}$, so $\lambda(v)=1$ for all $v$. Hence $\rho$ is injective. Explicitly, $\rho(a)$ exchanges the rectangles
$00\mathfrak C\times0\mathfrak C\times0\mathfrak C$ and $010\mathfrak C\times0\mathfrak C\times0\mathfrak C$, and $\rho(c)=\hat D\rho(ab)$, where $\rho(ab)$
permutes the rectangles with prefixes $(00,0,0)$, $(010,0,0)$ and $(011,0,0)$ cyclically in this order and $\hat D$ is the identity on $Z_0$.

\emph{The chunk.} The sofic profile of a chunk depends only on its identity element and its partial multiplication
(Section~\ref{sec:criterion}), and an injective homomorphism preserves both. So the chunk $E_{3V}=\rho(E)$, for the chunk $E$ of
Theorem~\ref{thm:main-group}, has $\sigma_{E_{3V}}=\sigma_E$, and parts (a) and (b) of that theorem give the bounds.
\end{proof}

The script \texttt{verification/embedding3v\_check.py} composes the prefix tables exactly and checks the identities of the proof and
that the images of $p,q,t,c,a$ satisfy the $25$ relations of Table~\ref{tab:certificate}.

By Section~\ref{sec:thompson}, Thompson's group $V$ has chunks with finite superpolynomial profile, from its Houghton subgroups, and a chunk with
profile at least $\exp(c\,r/\log^2r)$, which is finite if $F$ is amenable. Proposition~\ref{prop:3V} gives finite, nearly exponential profile in
$3V$ with no assumption. We do not know whether $\Gamma$ embeds in $V$ or in $2V$ (Section~\ref{sec:open}).

%======================================================================
\section{Houghton inputs}\label{sec:houghton}
%======================================================================

We need three facts about $H$: short words for finite permutations, generators of point stabilizers with linear word length, and a bound on the
Dehn function. All transport below takes place among the first points of ray~$1$, so we only need words for permutations supported there. For
$n\ge1$ let $\Sigma_n\le H$ be the group of permutations of $Y$ supported in $\{(1,1),\dots,(1,n)\}$, identified with $\Sym(n)$; we write
$(x\ y)$ for the transposition of the points $x$ and $y$ of the first ray.

\paragraph{Conventions.} Put $\alpha=[p,q]$; a direct check shows $\alpha=(1\ 2)$. In~\cite{DouglasMghirbiA} permutations act on the right,
and Johnson's presentation of $H$ (see~\cite[Theorem~2.8]{Lee2012}) is written there in generators $a,b,\alpha$ that are the inverses, as maps,
of our $p,q,\alpha$. The map $g\mapsto g^{-1}$ is an isomorphism from the permutations of $Y$ with the product of~\cite{DouglasMghirbiA} to the
permutations of $Y$ under composition, and it sends $a,b,\alpha$ to $p,q,\alpha$. So a word in $a,b,\alpha$ represents $1$ in the conventions
of~\cite{DouglasMghirbiA} if and only if the same word in $p,q,\alpha$ represents $1$ here. Hence the relators of~\cite[Section~2]{DouglasMghirbiA},
read in our letters, give the presentation
\begin{equation}\label{eq:johnson}
P_J=\bigl\langle p,q,\alpha\bigm|\ \alpha^2,\ (\alpha p\alpha p^{-1})^3,\ [\alpha,p^2\alpha p^{-2}],\ [p,q]\alpha^{-1},\ p\alpha p^{-1}q\alpha^{-1}q^{-1}\bigr\rangle
\end{equation}
of $H$; in particular $H=\langle p,q\rangle$. Renaming the letters is an isomorphism of free groups that carries the relators
of~\cite{DouglasMghirbiA} to those of $P_J$, so every area computed in~\cite{DouglasMghirbiA} for Johnson's presentation holds verbatim for $P_J$.

Eliminating $\alpha$ by the relator $[p,q]\alpha^{-1}$, a Tietze transformation, gives the presentation
\begin{equation}\label{eq:PH}
P_H=\bigl\langle p,q\bigm|\ s^2,\ (sps\,p^{-1})^3,\ [s,p^2sp^{-2}],\ psp^{-1}qs^{-1}q^{-1}\bigr\rangle,\qquad s=[p,q],
\end{equation}
of $H$. Let $\delta_H$ be its Dehn function: $\delta_H(n)$ is the largest area, with respect to the four relators of $P_H$, of a word of length at most
$n$ in $p^{\pm1},q^{\pm1}$ that represents $1$ in $H$. It is nondecreasing.

\begin{lemma}[Finite permutations]\label{lem:transpositions}
Put $\sigma_k=p^k\alpha p^{-k}$ for $k\ge0$.
\begin{enumerate}[label=(\alph*),itemsep=1pt]
\item $\sigma_k=(k+1\ \,k+2)$, and for $1\le x<y$ the word
$\omega_{x,y}=p^{x-1}(\alpha p)^{y-x-1}(\alpha p^{-1})^{y-x-1}\alpha\,p^{1-x}$ represents $(x\ y)$. With $\alpha=pqp^{-1}q^{-1}$ it has length at most
$10y-8x-8<10y$ in $p,q$.
\item Every element of $\Sigma_n$ that moves $m\ge2$ points is represented by a word in $p,q$ of length less than $10(m-1)n$.
\item Let $x_1,\dots,x_l\in\{1,\dots,n\}$ be distinct. Put $\pi_0=1$ and, for $m=1,\dots,l$, $w_m=\pi_{m-1}^{-1}(x_m)$ and
$\pi_m=\pi_{m-1}(m\ w_m)$, where $(m\ m)=1$. Then $\pi[x_1,\dots,x_l]:=\pi_l$ sends $m$ to $x_m$ for $m\le l$, it is a product of at most $l$
transpositions of points in $\{1,\dots,n\}$, and the product of the corresponding words $\omega$ represents it with length less than $10ln$.
\end{enumerate}
\end{lemma}

\begin{proof}
(a) Since $p^k$ maps $(1,1),(1,2)$ to $(1,k+1),(1,k+2)$, $\sigma_k$ is the stated transposition. Hence
$(x\ y)=\sigma_{x-1}\sigma_x\cdots\sigma_{y-3}\,\sigma_{y-2}\,\sigma_{y-3}\cdots\sigma_{x-1}$, by induction on $y-x$ from
$(x\ y)=\sigma_{x-1}(x{+}1\ \,y)\sigma_{x-1}$. In the product the powers of $p$ telescope: $\sigma_k\sigma_{k+1}=p^k\alpha p\alpha p^{-k-1}$ and
$\sigma_{k+1}\sigma_k=p^{k+1}\alpha p^{-1}\alpha p^{-k}$, which gives $\omega_{x,y}$ after free reduction. Its length is
$2(x-1)+10(y-x-1)+4$.

(b) A permutation moving $m$ points is a product of at most $m-1$ transpositions of those points, each of length less than $10n$ by (a).

(c) Since $\pi_{m-1}(l)=x_l\ne x_m$ for $l<m$, we have $w_m\ge m$, so $(m\ w_m)$ fixes $1,\dots,m-1$ and $\pi_m(l)=x_l$ for $l\le m$. Each factor
has length less than $10n$ by~(a).
\end{proof}

\begin{lemma}[Point stabilizers]\label{lem:stabilizers}
For $k\ge1$ let $H^{(k)}$ be the pointwise stabilizer in $H$ of $(1,1),\dots,(1,k)$, put $\zeta_k=\sigma_{k-1}\cdots\sigma_1\sigma_0$ and
$K_k=(\zeta_kp,\ \zeta_kq)$.
\begin{enumerate}[label=(\alph*),itemsep=1pt]
\item There is an isomorphism $H\to H^{(k)}$ that sends $p,q,\alpha$ to $\zeta_kp,\zeta_kq,\sigma_k$ and the transposition $(x\ y)$ to
$(x{+}k\ \ y{+}k)$. In particular $\zeta_kp$ and $\zeta_kq$ generate $H^{(k)}$.
\item Every element of $\Sigma_n\cap H^{(k)}$ that moves $m\ge2$ points is represented by a word of length less than $10(m-1)n$ in $\zeta_kp$ and
$\zeta_kq$.
\end{enumerate}
In $p,q$ the words $\zeta_1p=[p,q]\,p$ and $\zeta_2p=p[p,q]p^{-1}[p,q]\,p$ have length at most $5$ and $11$, and likewise for $q$.
\end{lemma}

\begin{proof}
(a) Let $\theta_k:Y\to Y\setminus\{(1,1),\dots,(1,k)\}$ send $(1,j)$ to $(1,j+k)$ and fix rays $2$ and $3$. For $g\in H$ let $g^{[k]}$ act as
$\theta_kg\theta_k^{-1}$ on the image of $\theta_k$ and fix $(1,1),\dots,(1,k)$. Then $g^{[k]}$ is eventually a translation on each ray, and
$g\mapsto g^{[k]}$ is an isomorphism $H\to H^{(k)}$, since an element of $H^{(k)}$ restricts to a permutation of the image of $\theta_k$. Clearly
$(x\ y)^{[k]}=(x{+}k\ \ y{+}k)$, so $\alpha^{[k]}=\sigma_k$. The element $\zeta_k$ is the cycle $(1,1)\mapsto(1,k+1)\mapsto(1,k)\mapsto\dots\mapsto(1,2)\mapsto(1,1)$.
Hence $\zeta_kp$ fixes $(1,j)$ for $j\le k$, maps $(1,j)$ to $(1,j+1)$ for $j>k$, maps $(2,1)$ to $(1,k+1)$, and agrees with $p$ elsewhere; that is,
$\zeta_kp=p^{[k]}$. Likewise $\zeta_kq=q^{[k]}$.

(b) An element $g\in\Sigma_n\cap H^{(k)}$ moving $m$ points is $g'^{[k]}$ for some $g'\in\Sigma_{n-k}$ moving $m$ points. Write $g'$ as in
Lemma~\ref{lem:transpositions}(b) and substitute $\zeta_kp,\zeta_kq$ for $p,q$.
\end{proof}

\begin{proposition}[Dehn function]\label{prop:dehn}
Let $\delta_J$ be the Dehn function of $P_J$. Then $\delta_H\le\delta_J$, and:
\begin{enumerate}[label=(\alph*),itemsep=1pt]
\item there is an absolute constant $K$ with $\delta_H(n)\le e^{Kn}$ for $n\ge1$;
\item there is an absolute constant $C$ with $\delta_H(n)\le Cn^{\,p_H}$ for $n\ge1$, where $p_H=\log107+4\approx10.74$.
\end{enumerate}
\end{proposition}

\begin{proof}
Let $w$ be a word in $p,q$ of length at most $n$ representing $1$. In the free group on $p,q,\alpha$ it is a product of at most $\delta_J(n)$
conjugates of relators of $P_J$ and their inverses. Apply the homomorphism to the free group on $p,q$ that fixes $p,q$ and sends $\alpha$ to
$[p,q]$. It fixes $w$, sends $[p,q]\alpha^{-1}$ to the trivial word, and sends each of the other four relators to a relator of $P_H$. Hence
$\Area_{P_H}(w)\le\delta_J(n)$.

(a) Lee proved that $H_3$ satisfies an exponential isoperimetric inequality~\cite{Lee2012}. The Dehn functions $\delta,\delta'$ of two finite
presentations of the same group are equivalent: $\delta(n)\le C'\delta'(C'n)+C'n+C'$ for a constant $C'$~\cite{Bridson2002}. This bound
preserves the class of functions bounded by $e^{Kn}$.

(b) By~\cite[Theorem~3.3]{DouglasMghirbiA}, $\delta_J(n)\le400n^4(1+A^*(4n+2))$, where the far-commutator area satisfies
$A^*(M)\le135M^{\log107}$~\cite[Theorem~3.2]{DouglasMghirbiA}; so $\delta_J(n)=O(n^{\log107+4})$. By the discussion of conventions above this
holds for $P_J$ in our letters, and $\delta_H\le\delta_J$.
\end{proof}

Part~(a) uses only the published exponential bound, and part~(b) our polynomial one. Either enters the lower bounds below only through
$\delta_H$ at scale $N$, the number of points that carry the configurations.

%======================================================================
\section{A certificate of 25 relations and local fillings}\label{sec:certificate}
%======================================================================

All fillings in $\Gamma$ are built from one finite set $R$ of words in the five generators of Lemma~\ref{lem:structure}, each representing $1$
in $\Gamma$. The set $R$ does not present $\Gamma$ (Section~\ref{sec:fp}); areas are always taken with respect to $R$.

\paragraph{Abbreviations.} Put
\[
s=[p,q],\qquad B=psp^{-1},\qquad t_m=p^{m-1}\,t\,p^{1-m}\quad(m=1,2,3),\qquad d=c^3,\qquad b=ac^4 .
\]
By Lemma~\ref{lem:transpositions}, $s$ and $B$ are the transpositions $(1\ 2)$ and $(2\ 3)$, and $t_m$ represents $\tau_m$, since
$p^{m-1}(1,1)=(1,m)$. By Lemma~\ref{lem:stabilizers}, with $\zeta_1=s$ and $\zeta_2=Bs$, the pairs
\[
K_1=(sp,\ sq),\qquad K_2=(Bsp,\ Bsq)
\]
generate the pointwise stabilizers $H_1$ of the point $1$ and $H_2$ of the points $1,2$. The elements $t$ and $d$ have the \emph{prototypes}
$P_t=\{1\}$ and $P_d=\{1,2\}$, the points on which they depend: by~\eqref{eq:affine-identities}, an element of $H$ that fixes $P_f$ pointwise
commutes with $f$.

\begin{table}[ht]
\centering\small
\begin{tabular}{@{}l p{0.64\textwidth} r r@{}}
\toprule
 & relations & number & length\\
\midrule
\textsf{H} & $s^2$,\ $(sps\,p^{-1})^3$,\ $[s,p^2sp^{-2}]$,\ $psp^{-1}qs^{-1}q^{-1}$ & 4 & 64\\
\textsf{Q} & $t^2$ & 1 & 2\\
\textsf{S} & $[t,k]$ for $k\in K_1$;\ \ $[d,k]$ for $k\in K_2$ & 4 & 56\\
\textsf{Z} & $[g,x]$ for $g\in\{p,q\}$, $x\in\{a,b\}$;\ \ $[d,a]$ & 5 & 40\\
\textsf{C} & $[t_1,t_2]$ & 1 & 8\\
\textsf{A} & $d\,t_1d^{-1}(t_1t_2)^{-1}$,\ $[d,t_2]$,\ $[d,t_3]$ & 3 & 39\\
\textsf{L} & $a^2$,\ $b^2$,\ $(ab)^3$ & 3 & 30\\
\textsf{P} & $[x,tyt^{-1}]$ for $x,y\in\{a,b\}$ & 4 & 64\\
\midrule
 & total & 25 & 303\\
\bottomrule
\end{tabular}
\caption{The certificate $R$: eight families of relations in the five generators $p,q,t,c,a$, written with the abbreviations above. Each
relation is taken to be the freely reduced word obtained by expanding the abbreviations; lengths are those of these words, and the longest has
length $24$.}
\label{tab:certificate}
\end{table}

\begin{lemma}\label{lem:validity}
Every word of Table~\ref{tab:certificate} represents~$1$ in $\Gamma$.
\end{lemma}

\begin{proof}
\textsf{H}: the words of~\eqref{eq:PH} hold in $H$.
\textsf{Q}: $V$ has exponent two.
\textsf{S}: the elements of $H_1$ fix the point $1$, so they commute with $t=\tau_1$; those of $H_2$ fix the points $1$ and $2$, so they commute
with $d=E_{2\leftarrow1}$ (Lemma~\ref{lem:structure}(a) and~\eqref{eq:affine-identities}).
\textsf{Z}: $p$, $q$ and $d$ are linear, so they commute with $a$ and with $b=ac^4=a\cdot ab$, which lie in the copy of $S_3$ at~$0$.
\textsf{C}: translations commute.
\textsf{A}: $dt_md^{-1}=\tau_{d(e_m)}$, and $d(e_1)=e_1+e_2$, while $d$ fixes $e_2$ and $e_3$.
\textsf{L}: these are relations of $S_3$.
\textsf{P}: $tyt^{-1}$ lies in the copy of $S_3$ at $e_1$, which commutes with the copy at~$0$.
\end{proof}

The $25$ freely reduced words are listed in the file \texttt{verification/relators1d.txt}. The script \texttt{verification/relators1d\_check.py}
rebuilds each word from the abbreviations and checks it, independently of Lemma~\ref{lem:validity}, in the faithful action of $\Gamma$ on
$V\times S_3$. For a word $w$ of length $l$ it checks that the exponent sums of $p$ and $q$ vanish; that $w$ fixes the empty configuration and
every configuration $\{y\}$ with $y$ of height at most $l+2$ or far out on a ray; and that it fixes $(v,s)$ for every $s\in S_3$ at every
configuration $v$ at which one of its lamp letters acts. These checks are exhaustive: while $w$ is applied, a point of height greater than $l+2$
never reaches the points $1$ and $2$ and is moved only by translations along its ray, so the checks determine the affine part of $w$, and then
the lamps of $w$ sit at the configurations just listed. The script also checks random configurations, that the abbreviations act as stated, and,
as a negative control, that deleting any single letter from a sample of eleven relations produces a word that acts nontrivially.

\subsection{Local fillings}

Let $n\ge3$. We use three elementary facts about areas. The area is invariant under free equality, inversion and conjugation, and subadditive
under products. If $N$ represents~$1$, then $\mu N\nu=(\mu N\mu^{-1})\mu\nu$ for all words $\mu,\nu$; so replacing a subword $U$ by $U'$ changes
the area by at most $\Area_R(UU'^{-1})$. And $[w,z_1z_2]=[w,z_1]\cdot z_1[w,z_2]z_1^{-1}$ and $[w,z^{-1}]=z^{-1}[w,z]^{-1}z$, so
$\Area_R([w,z])$ is at most the sum of the areas $\Area_R([w,z_l])$ over the letters $z_l$ of any factorization $z=z_1^{\pm1}\cdots z_m^{\pm1}$.

\paragraph{Canonical transporters.} For points $i\ne j$ in $\{1,\dots,n\}$ put
\[
U_i=\pi[i],\qquad \tau_i=U_i\,t\,U_i^{-1},\qquad U_{ij}=\pi[j,i],\qquad X_{ij}=U_{ij}\,d\,U_{ij}^{-1},
\]
with the words of Lemma~\ref{lem:transpositions}(c). We use $\tau_i$ both for this word and for the translation by $e_i$ that it represents, since
$U_i(1)=i$; and $X_{ij}$ represents $E_{i\leftarrow j}$, since $U_{ij}$ maps $1$ to $j$ and $2$ to $i$. By Lemma~\ref{lem:transpositions},
$|U_i|<10n$ and $|U_{ij}|<20n$, and $U_i$ ($U_{ij}$) is a product of at most one (two) transpositions.

The next lemma says that it does not matter which short word transports a prototype: two such words differ by an element of a stabilizer,
which commutes with the transported element by the family~\textsf{S}. We allow the transporting word to end with a power of $p$, because the
relations of Table~\ref{tab:certificate} use $t_2=ptp^{-1}$ and $t_3=p^2tp^{-2}$.

\begin{lemma}[Change of transporter]\label{lem:exchange}
Let $n\ge3$ and $f\in\{t,d\}$. Let $U$ and $W_0$ be words in $p^{\pm1},q^{\pm1}$ of length less than $30n$ that represent elements of $\Sigma_n$,
such that $W_0^{-1}U$ moves at most $10$ points, and let $W=W_0p^j$, where $j=0$ if $f=d$ and $j\in\{0,1,2\}$ if $f=t$. If $W$ and $U$ agree
on $P_f$, then
\[
\Area_R\bigl(WfW^{-1}\cdot Uf^{-1}U^{-1}\bigr)\ \le\ \mathcal E(n):=2\,\delta_H(1500n)+130n .
\]
\end{lemma}

\begin{proof}
Let $k=|P_f|$ and let $g$ be the element represented by $W^{-1}U$; it fixes $P_f$ pointwise, so $g\in H^{(k)}$. If $j=0$, then $g\in\Sigma_n$ moves at
most $10$ points, and by Lemma~\ref{lem:stabilizers}(b) it is represented by a word $Z$ with fewer than $90n$ letters from $K_k^{\pm1}$. If
$j\ge1$, then $f=t$, $k=1$, and $g=(sp)^{-j}g'$ with $g'=(sp)^jp^{-j}\cdot W_0^{-1}U$. Here $(sp)p^{-1}=s$ and $(sp)^2p^{-2}=sB$ lie in
$\Sigma_3$, so $g'\in\Sigma_n$ moves at most $13$ points; and $g'$ fixes the point $1$, since $g$ and $sp$ do. By Lemma~\ref{lem:stabilizers}(b),
$g'$ is represented by a word $Z'$ with fewer than $120n$ letters from $K_1^{\pm1}$, and we put $Z=(sp)^{-j}Z'$. In both cases $Z$ has fewer than
$122n$ letters from $K_k^{\pm1}$, each of length at most $11$ in $p,q$.

Let $N=W^{-1}UZ^{-1}$. It represents $1$ in $H$ and has length less than $30n+2+30n+1342n\le1500n$, so by the family \textsf{H} its area is at
most $\delta_H(1500n)$. By the family \textsf{S}, $\Area_R([f,Z])<122n$. Since $U=WNZ$ freely,
\[
WfW^{-1}\cdot Uf^{-1}U^{-1}=W\,(fNf^{-1})\,[f,Z]\,N^{-1}\,W^{-1},
\]
whose area is at most $2\delta_H(1500n)+122n\le\mathcal E(n)$.
\end{proof}

\begin{lemma}[Local fillings]\label{lem:local-filling}
Let $n\ge3$ and
\[
L(n)=10\,\delta_H(1500n)+700\,n .
\]
Let $i,j,k\in\{1,\dots,n\}$ be points with $i\ne j$, and $x,y\in\{a,b\}$, where $b$ stands for the word $ac^4$. The following words represent $1$,
and each has area at most $L(n)$ with respect to $R$:
\begin{enumerate}[label=(\alph*),itemsep=1pt]
\item $\tau_i^2$, of area $1$, and $[\tau_i,\tau_j]$;
\item $X_{ij}\tau_kX_{ij}^{-1}\tau_k^{-1}$ if $k\ne j$, and $X_{ij}\tau_jX_{ij}^{-1}\tau_i^{-1}\tau_j^{-1}$;
\item $[X_{ij},x]$, of area less than $40n+1$;
\item $[x,\tau_iy\tau_i^{-1}]$, of area less than $40n+1$.
\end{enumerate}
\end{lemma}

\begin{proof}
All the words represent $1$ by~\eqref{eq:affine-identities}. Call replacing a subword $UfU^{-1}$ by $WfW^{-1}$ as in Lemma~\ref{lem:exchange} a
\emph{change of transporter}; it costs at most $\mathcal E(n)$, and $5\mathcal E(n)+1\le L(n)$. Below, $W_0$ is always a product of at most three
transpositions, of length less than $30n$, and $U\in\{U_i,U_j,U_k,U_{ij}\}$ is a product of at most two, so $W_0^{-1}U$ moves at most $10$ points.

(a) $\tau_i^2=U_it^2U_i^{-1}$ freely, a conjugate of the relation \textsf{Q}. For $[\tau_i,\tau_j]$ let $W_0=\pi[i,j]$. Replace both occurrences
of $\tau_i^{\pm1}$ by $(W_0tW_0^{-1})^{\pm1}$, and both occurrences of $\tau_j^{\pm1}$ by $(W_0t_2W_0^{-1})^{\pm1}=(W_0p\,t\,p^{-1}W_0^{-1})^{\pm1}$.
These are four changes of transporter, since $W_0$ and $U_i$ both map $1$ to $i$, and $W_0p$ and $U_j$ both map $1$ to $j$. The result
$W_0[t_1,t_2]W_0^{-1}$ is a conjugate of the relation \textsf{C}, so the area is at most $4\mathcal E(n)+1$.

(b) Let $W_0=\pi[j,i,k]$, or $\pi[j,i]$ if $k\in\{i,j\}$, so that $W_0$ maps $1$ to $j$, $2$ to $i$, and $3$ to $k$ when $k\notin\{i,j\}$; let
$m\in\{1,2,3\}$ be the point with $W_0(m)=k$, so that $m=1$ exactly when $k=j$. Replace $X_{ij}^{\pm1}$ by $(W_0dW_0^{-1})^{\pm1}$, a change of
transporter because $W_0$ and $U_{ij}$ agree on $P_d=\{1,2\}$. Replace $\tau_k^{\pm1}$ by $(W_0t_mW_0^{-1})^{\pm1}=(W_0p^{m-1}\,t\,p^{1-m}W_0^{-1})^{\pm1}$,
and $\tau_i^{\pm1}$ by $(W_0t_2W_0^{-1})^{\pm1}$, again changes of transporter. If $k\ne j$, four of them turn the word into
$W_0(dt_md^{-1}t_m^{-1})W_0^{-1}$ with $m\in\{2,3\}$, a conjugate of a relation \textsf{A}. If $k=j$, five of them turn it into
$W_0(dt_1d^{-1}t_2^{-1}t_1^{-1})W_0^{-1}$, again a conjugate of a relation \textsf{A}. So the area is at most $5\mathcal E(n)+1$.

(c) Moving $x$ across a letter of $U_{ij}$ costs one relation \textsf{Z}, and moving it across $d=c^3$ costs one relation $[d,a]$, since
$[c^3,ac^4]=[c^3,a]$ freely. Hence $\Area_R([X_{ij},x])\le2|U_{ij}|+1<40n+1$.

(d) Freely the word is $U_i[U_i^{-1}xU_i,\,t(U_i^{-1}yU_i)t^{-1}]U_i^{-1}$. Replacing the two occurrences of $U_i^{-1}x^{\pm1}U_i$ and the two of
$U_i^{-1}y^{\pm1}U_i$ by $x^{\pm1}$ and $y^{\pm1}$ costs at most $4|U_i|$ relations \textsf{Z} and leaves $U_i[x,tyt^{-1}]U_i^{-1}$, a conjugate of a
relation \textsf{P}; so the area is at most $4|U_i|+1<40n+1$.
\end{proof}

%======================================================================
\section{Pair areas and the profile}\label{sec:profile}
%======================================================================

We now bring together two of the $2^N$ copies of $S_3$ supported on the first $N$ points of ray~$1$. Their transport words are products of the
$\tau_i$; the commutator of two copies reduces to the copy at a single point after clearing the difference of the two configurations by at most
$N$ transvections. Clearing naively would let the translation word grow; we keep it normalized after every step, so that it never has more than
$N$ factors, and each step costs $O(N)$ local fillings.

\begin{theorem}[Pair areas]\label{thm:pair-area}
Let $N\ge1$ and $n=\max(3,N)$. For a configuration $u$ supported in the first $N$ points of ray~$1$ let $T_u=\tau_{i_1}\cdots\tau_{i_m}$,
where $i_1<\dots<i_m$ are the points of the support of $u$. Then $T_u$ represents $\tau_u$, so $T_uaT_u^{-1}$ and $T_ubT_u^{-1}$ generate the copy
of $S_3$ at $u$; $|T_u|\le21Nn$; and for distinct such $u,w$ and $x,y\in\{a,b\}$, with $b$ the word $ac^4$,
\[
\Area_R\bigl([T_uxT_u^{-1},\,T_wyT_w^{-1}]\bigr)\ \le\ \Lambda_N:=20N^2L(n)\ \le\ C_0N^2\bigl(\delta_H(C_0N)+N\bigr)
\]
with $C_0=4.5\cdot10^4$. Consequently there are absolute constants $K'$ and $C'$ with
\[
\Lambda_N\le e^{K'N}\qquad\text{and}\qquad \Lambda_N\le C'N^{2+p_H}\qquad(N\ge1),\qquad p_H=\log107+4 .
\]
\end{theorem}

\begin{proof}
Translations commute, so $T_u$ represents $\tau_u$, and $|T_u|\le N(2\cdot10n+1)\le21Nn$. Throughout, $L=L(n)$; by
Lemma~\ref{lem:local-filling} the local identities (a)--(d) have area at most $L$.

\emph{Step 1: one translation word.} Freely, $[T_uxT_u^{-1},T_wyT_w^{-1}]=T_u[x,ZyZ^{-1}]T_u^{-1}$ with $Z=T_u^{-1}T_w$. Replace each of the at
most $N$ factors $\tau_i^{-1}$ of $T_u^{-1}$ by $\tau_i$ (one relation \textsf{Q} each), sort the resulting product of at most $2N$ factors
$\tau_i$ into increasing order by at most $\binom{2N}2$ swaps of adjacent distinct factors (Lemma~\ref{lem:local-filling}(a)), and cancel the at
most $N$ adjacent equal pairs (one relation \textsf{Q} each). The result is $T_{u_0}$ with $u_0=u+w\ne0$, and
$\Area_R(ZT_{u_0}^{-1})\le2N+N(2N-1)L\le2N^2L$. Since $Z^{\pm1}$ occurs four times in $[x,ZyZ^{-1}]$,
\[
\Area_R\bigl([T_uxT_u^{-1},T_wyT_w^{-1}]\bigr)\le8N^2L+\Area_R\bigl([x,T_{u_0}yT_{u_0}^{-1}]\bigr).
\]

\emph{Step 2: normalized clearing.} Let $s\le N$ be the size of the support of $u_0$, choose a point $j$ in it, and list the other points of the
support as $i_1,\dots,i_{s-1}$. Put $G_l=X_{i_lj}\cdots X_{i_1j}$ and $u_l=u_0+e_{i_1}+\dots+e_{i_l}$, whose support is that of $u_0$ with
$i_1,\dots,i_l$ removed; it still contains $j$, and $u_{s-1}=e_j$. We claim
\[
\Area_R\bigl(G_lT_{u_0}G_l^{-1}\,T_{u_l}^{-1}\bigr)\le2lsL\qquad(0\le l\le s-1).
\]
For $l=0$ there is nothing to prove. Given the claim for $l$, write $X=X_{i_{l+1}j}$ and $i=i_{l+1}$. Freely
$G_{l+1}T_{u_0}G_{l+1}^{-1}=X(G_lT_{u_0}G_l^{-1})X^{-1}$, which differs from $XT_{u_l}X^{-1}$ by a conjugate of a word of area at most $2lsL$. Next,
$XT_{u_l}X^{-1}$ is freely the product of the words $X\tau_{i'}X^{-1}$ over the support of $u_l$, in order. By Lemma~\ref{lem:local-filling}(b),
at a cost of at most $sL$, each of these becomes $\tau_{i'}$, except that $X\tau_jX^{-1}$ becomes $\tau_j\tau_i$. The new factor $\tau_i$ has a
partner $\tau_i$ in $T_{u_l}$, because $i$ lies in the support of $u_l$; at most $s-1$ swaps (Lemma~\ref{lem:local-filling}(a)) make the two
adjacent, and one relation \textsf{Q} cancels them. What remains is $T_{u_{l+1}}$. This step costs at most $sL+(s-1)L+1\le2sL$, which proves the
claim. The translation word is renormalized after every transvection and never has more than $s\le N$ factors.

\emph{Step 3: a single point.} Let $G=G_{s-1}$ and let $M=GT_{u_0}G^{-1}\tau_j^{-1}$, which has area at most $2s(s-1)L\le2N^2L$ by Step~2. Freely
$T_{u_0}=G^{-1}M\tau_jG$. Substituting this into the four occurrences of $T_{u_0}^{\pm1}$ in $[x,T_{u_0}yT_{u_0}^{-1}]$ and deleting the four
occurrences of $M^{\pm1}$ costs at most $8N^2L$. By Lemma~\ref{lem:local-filling}(c), $\Area_R([G,z])\le(s-1)(40n+1)$ for $z\in\{a,b\}$, and the
same bound holds for $[G,z^{-1}]$ and $[z,G^{-1}]$. Moving $y^{\pm1}$ across $G$ and $G^{-1}$, and then $x^{\pm1}$, leaves
$G^{-1}[x,\tau_jy\tau_j^{-1}]G$, at a cost of at most $4(s-1)(40n+1)$; and $[x,\tau_jy\tau_j^{-1}]$ has area at most $40n+1$ by
Lemma~\ref{lem:local-filling}(d). Altogether, since $L\ge700n$,
\[
\Area_R\bigl([x,T_{u_0}yT_{u_0}^{-1}]\bigr)\le8N^2L+4N(40n+1)+40n+1\le9N^2L,
\]
and with Step~1 the area is at most $17N^2L\le\Lambda_N$.

\emph{Growth.} Since $n\le3N$, $L(n)\le10\delta_H(4500N)+2100N$, which gives the first bound on $\Lambda_N$. With Proposition~\ref{prop:dehn}(a)
it is at most $e^{K'N}$ for a suitable $K'$. With Proposition~\ref{prop:dehn}(b) it is at most $C_0N^2(C(C_0N)^{p_H}+N)\le C'N^{2+p_H}$, since $p_H>1$.
\end{proof}

The factor $N^2$ counts local moves: at most $N$ transvections, each followed by at most $N$ commutations. The Dehn function of $H$ enters only
through the local fillings, at the scale $N$ of the transport.

\subsection{Affine hashing and the upper bound}\label{sec:hashing}

The Følner sets of $\Gamma$ are doubly exponential (Proposition~\ref{prop:folner-lower}), but its chunks have far smaller models. A Følner set
records the lamp at every configuration of a large finite part of $V$. A model need only record the lamps that are lit, and it can store them
through a random affine hash of the configurations, which the affine group moves covariantly. Random $\F_2$-linear maps form a universal
hash family of Carter and Wegman~\cite{CarterWegman1979}, and a random translation does not change collision probabilities.

Let $W$ be a finite-dimensional vector space over $\F_2$, and let $K_W=S_3^W\rtimes\AGL(W)$, where $S_3^W$ is the group of functions
$W\to S_3$ under pointwise multiplication, and $g\in\AGL(W)$ acts by $(g\lambda)(w)=\lambda(g^{-1}w)$; thus $(\lambda,g)(\mu,h)=(\lambda\cdot g\mu,\,gh)$.
The support of $\lambda$ is the set of $w$ with $\lambda(w)\ne1$. Fix $k\ge1$ and a total order on $W$, let $\operatorname{Aff}(W,\F_2^k)$ be the
set of affine maps $W\to\F_2^k$, and let
\[
\Omega=\operatorname{Aff}(W,\F_2^k)\times S_3^{\F_2^k},\qquad |\Omega|=2^{k(\dim W+1)}\,6^{2^k}.
\]
For $\eta\in\operatorname{Aff}(W,\F_2^k)$ and $\lambda\in S_3^W$ let $\lambda^\eta:\F_2^k\to S_3$ send $z$ to the product, in the order of $W$, of the
$\lambda(w)$ with $\eta(w)=z$. For $\gamma=(\lambda,g)\in K_W$ define a permutation $\Theta_\gamma$ of $\Omega$ by
\[
\Theta_\gamma(\eta,u)=\bigl(\eta\circ g^{-1},\ \lambda^{\eta\circ g^{-1}}\cdot u\bigr);
\]
its inverse sends $(\eta',u')$ to $(\eta'\circ g,(\lambda^{\eta'})^{-1}u')$, and $\Theta_1$ is the identity.

\begin{lemma}[Affine hashing]\label{lem:hash}
Let $b\ge1$, let $\gamma_1=(\lambda,g)$ and $\gamma_2=(\mu,h)$ be elements of $K_W$ whose lamp parts $\lambda,\mu$ have supports of size at most
$b$, and let $\omega$ be uniformly distributed on $\Omega$. Then
\begin{enumerate}[label=(\roman*),itemsep=1pt]
\item $\Theta_{\gamma_1}\Theta_{\gamma_2}\omega=\Theta_{\gamma_1\gamma_2}\omega$ except with probability at most $\binom{2b}2\,2^{-k}$;
\item if $\gamma_1\ne\gamma_2$, then $\Theta_{\gamma_1}\omega\ne\Theta_{\gamma_2}\omega$ except with probability at most $\binom{2b}2\,2^{-k}$.
\end{enumerate}
\end{lemma}

\begin{proof}
Write $\omega=(\eta,u)$ and $\eta=\eta_0+z_0$ with $\eta_0$ linear; then $\eta_0$ and $z_0$ are independent and uniform. For $w\ne w'$ in $W$,
$\eta(w)-\eta(w')=\eta_0(w-w')$ is uniform on $\F_2^k$, because $\eta_0\mapsto\eta_0(w-w')$ is a surjective linear map; so
$\eta(w)=\eta(w')$ with probability exactly $2^{-k}$. Precomposition with a fixed affine bijection of $W$ permutes $\operatorname{Aff}(W,\F_2^k)$, so
$\eta\circ f$ is again uniform for every $f\in\AGL(W)$. Hence, for a set $S\subset W$ of size at most $2b$, $\eta\circ f$ is injective on $S$ except
with probability at most $\binom{2b}2\,2^{-k}$. If $\eta'$ is injective on a set $S$ containing the support of $\lambda$, then
$\lambda^{\eta'}(\eta'(w))=\lambda(w)$ for $w\in S$, and $\lambda^{\eta'}=1$ off $\eta'(S)$.

(i) Put $\eta'=\eta\circ(gh)^{-1}$. Then $\Theta_{\gamma_1}\Theta_{\gamma_2}(\eta,u)=(\eta',\lambda^{\eta'}\mu^{\eta\circ h^{-1}}u)$ and
$\Theta_{\gamma_1\gamma_2}(\eta,u)=(\eta',(\lambda\cdot g\mu)^{\eta'}u)$. Let $S$ be the union of the support of $\lambda$ and the image under $g$ of
the support of $\mu$, and suppose that $\eta'$ is injective on $S$. Since $\eta\circ h^{-1}(v)=\eta'(gv)$, the map $\eta\circ h^{-1}$ is injective on
the support of $\mu$, and $\mu^{\eta\circ h^{-1}}(\eta'(w))=\mu(g^{-1}w)$ for $w\in S$. The support of $\lambda\cdot g\mu$ lies in $S$. So all three
functions are trivial off $\eta'(S)$, and at $\eta'(w)$, $w\in S$, both sides equal $\lambda(w)\mu(g^{-1}w)$.

(ii) If $g\ne h$, choose $w_0$ with $h^{-1}g(w_0)\ne w_0$. If $\eta\circ g^{-1}=\eta\circ h^{-1}$, then $\eta=\eta\circ h^{-1}g$, so
$\eta(h^{-1}g(w_0))=\eta(w_0)$, which has probability $2^{-k}$. If $g=h$, then $\lambda\ne\mu$; put $\eta'=\eta\circ g^{-1}$ and let $S$ be the union
of the supports. If $\eta'$ is injective on $S$, then $\lambda^{\eta'}$ and $\mu^{\eta'}$ take the different values $\lambda(w)$ and $\mu(w)$ at
$\eta'(w)$ for some $w\in S$.
\end{proof}

To use the lemma in $\Gamma$ we need finite subgroups that contain the relevant cocycles. For $z\in\Z^2$ put $g_z=p^{z_1}q^{z_2}$, so that
$\pi(g_z)=z$, and for $m\ge1$ let $B_m=\{-m,\dots,m\}^2$. For $R\ge1$ let $Y_R\subset Y$ be the set of points of height at most $R$, let
$W_R=\F_2^{Y_R}\subset V$, and let $K_R\le\ker\pi$ consist of the elements whose lamps lie at configurations in $W_R$ and whose affine part lies in
$W_R\rtimes\GL(W_R)$, with $\GL(W_R)$ extended by the identity on the other basis vectors. Restricting to $W_R$ identifies $K_R$ with $K_{W_R}$.

\begin{lemma}[Cocycles]\label{lem:clock}
Let $m\ge1$, $R=2m+3$, $z\in B_m$ and $s\in\{p,q,t,c,a\}^{\pm1}$. Then $g_{z+\pi(s)}^{-1}sg_z$ and $g_zsg_{z+\pi(s)}^{-1}$ lie in $K_R$, and their
lamps sit at no more than one configuration.
\end{lemma}

\begin{proof}
Each of $p^{\pm1},q^{\pm1}$ changes heights by at most one and acts on the points of height at least $2$ as a translation along the rays, and these
translations commute. Hence for $(y_1,y_2)\in\Z^2$ the commutator $[p^{y_1},q^{y_2}]$ fixes every point of height greater than $|y_1|+|y_2|+1$:
while it is applied, such a point stays at height at least $2$. For $s=p^{\pm1}$ we have $\pi(s)=(\pm1,0)$, $g_{z+\pi(s)}^{-1}sg_z=1$ and
$g_zsg_{z+\pi(s)}^{-1}=p^{z_1}[q^{z_2},p^{\pm1}]p^{-z_1}$, which is supported in heights at most $|z_2|+2+|z_1|\le2m+2$. For $s=q^{\pm1}$ we have
$\pi(s)=(0,\pm1)$, $g_zsg_{z+\pi(s)}^{-1}=1$ and $g_{z+\pi(s)}^{-1}sg_z=q^{-z_2\mp1}[p^{-z_1},q^{\pm1}]q^{z_2\pm1}$, supported in heights at most
$|z_1|+2+|z_2|+1\le2m+3$. These are permutations of
$Y_R$, hence elements of $\GL(W_R)$. For $s\in\{t,c,a\}^{\pm1}$ we have $\pi(s)=0$, and $g_z^{\pm1}$ moves the points $1$ and $2$ to height at most
$2m+2$; so the conjugates of $t$ and $d$ by $g_z^{\pm1}$ are a translation and a transvection of $W_R$, while those of $a$ and $ab$ are $a$ and
$ab$, at the configuration~$0$.
\end{proof}

\begin{proposition}[Hash models]\label{prop:upper}
Let $E'$ be a chunk of $\Gamma$ each of whose elements is a product of at most $\ell\ge1$ of the generators $p,q,t,c,a$ and their inverses. Then
\[
\log\sigma_{E'}(r)\ \le\ 500\,\ell^2\,r\log(2+r)\qquad(r>1).
\]
\end{proposition}

\begin{proof}
Let $m=\lceil8\ell r\rceil$, $R=2m+3$, $D=|Y_R|=6m+9$ and $k=\lceil\log(4\ell^2r)\rceil$, and let $\Omega$ be the set above for $W=W_R$. For $g\in\Gamma$
and $z\in\Z^2$ put $\psi(g,z)=g_{z+\pi(g)}^{-1}gg_z\in\ker\pi$; then $\psi(gh,z)=\psi(g,z+\pi(h))\,\psi(h,z)$. If $g$ is a product of at most $\ell$
generators and inverses and $z\in B_{m-\ell}$, then $\psi(g,z)$ is the product of the $\psi(s,z')$ for the letters $s$ of the word, at points
$z'\in B_m$, so by Lemma~\ref{lem:clock} it lies in $K_R$ with lamps at no more than $\ell$ configurations.

Let $X=B_m\times\Omega$. For $g\in E'$ let $F_g$ be a permutation of $X$ that extends the injective partial map
$(z,\omega)\mapsto(z+\pi(g),\Theta_{\psi(g,z)}\omega)$, defined for $z\in B_{m-\ell}$, with $F_1=\mathrm{id}$. Let $g,h,gh\in E'$ and
$z\in B_{m-2\ell}$. Then $z$ and $z+\pi(h)$ lie in $B_{m-\ell}$, so $F_gF_h(z,\omega)=(z+\pi(gh),\Theta_{\psi(g,z+\pi(h))}\Theta_{\psi(h,z)}\omega)$ and
$F_{gh}(z,\omega)=(z+\pi(gh),\Theta_{\psi(g,z+\pi(h))\psi(h,z)}\omega)$; by Lemma~\ref{lem:hash}(i) these agree for all but a fraction
$\binom{2\ell}2\,2^{-k}$ of the $\omega$. Since $|B_{m-2\ell}|\ge(1-8\ell/(2m+1))|B_m|$,
\[
d\bigl(F_gF_h,F_{gh}\bigr)\ \le\ \varepsilon:=\frac{8\ell}{2m+1}+\binom{2\ell}2\,2^{-k}<\frac1{2r}+\frac1{2r}=\frac1r .
\]
If $g\ne h$ in $E'$ and $z\in B_{m-\ell}$, then either $\pi(g)\ne\pi(h)$ and $F_g(z,\omega)\ne F_h(z,\omega)$, or $\psi(g,z)\ne\psi(h,z)$ and
Lemma~\ref{lem:hash}(ii) applies; so $d(F_g,F_h)\ge1-\varepsilon$. Hence $\sigma_{E'}(r)\le|X|$.

Finally $\log|X|=2\log(2m+1)+k(D+1)+2^k\log6$. Here $2^k<8\ell^2r$, $k<\log(8\ell^2r)\le5\ell\log(2+r)$ and $D+1\le64\ell r$, so
$k(D+1)\le320\ell^2r\log(2+r)$; also $2^k\log6\le21\ell^2r\le14\ell^2r\log(2+r)$ and $2\log(2m+1)\le10\ell^2r\log(2+r)$. The sum is less than
$500\ell^2r\log(2+r)$.
\end{proof}

\subsection{The profile}

\begin{theorem}\label{thm:main-group}
Let $E\subset\Gamma$ be the chunk consisting of $1$, the five generators $p,q,t,c,a$ and their inverses, the elements represented by the prefixes
of the $25$ freely reduced words of Table~\ref{tab:certificate}, and the inverses of these elements. Let $p_H=\log107+4$ and
\[
\kappa=2+p_H=6+\log107<12.75 .
\]
There are absolute constants $c,C>0$ such that for every $r\ge2$:
\begin{enumerate}[label=(\alph*),itemsep=1pt]
\item $\log\sigma_E(r)\ge c\,r/(\log r)^\kappa$;
\item $\log\sigma_E(r)\le C\,r\log r$; more generally, every chunk $E'$ of $\Gamma$ satisfies $\log\sigma_{E'}(r)\le C_{E'}\,r\log r$.
\end{enumerate}
In particular $\sigma_E(r)=\exp(r^{1+o(1)})$. Moreover:
\begin{enumerate}[label=(\alph*),itemsep=1pt,resume]
\item For every $0<\Delta\le1$ there are $c_\Delta,e_\Delta>0$ such that the following holds for $0<e\le e_\Delta$. Let $\phi$ assign unitaries in
$U(D)$ to the five generators so that every word of Table~\ref{tab:certificate} is mapped within $e$ of $I$ in the normalized Hilbert--Schmidt
norm, and put $\phi(b)=\phi(a)\phi(c)^4$. If $\|\phi(a)\phi(b)-I\|_2\ge\Delta$, then
\[
\log D\ \ge\ c_\Delta\,\frac{e^{-2}}{(\log(1/e))^{2\kappa}},\qquad 2\kappa<25.49 .
\]
\item $\Gamma$ is finitely presented, elementary amenable and not LEF.
\end{enumerate}
\end{theorem}

\begin{proof}
\emph{Setting up the criterion.} Theorem~\ref{thm:criterion} asks for $a$ and $b$ among the generators, while in $R$ the letter $b$ abbreviates
$ac^4$. So let $S^+$ consist of the five generators and $b=ac^4\in\Gamma$, and let $R^+$ consist of $R$ and the three words $b^{-1}ac^4$, $b^2$
and $(ab)^3$ in the letters of $S^+$. These words represent $1$, $a^2,b^2,(ab)^3\in R^+$, $ab\ne1$, and the longest word of $R^+$ has length
$l=24$. The chunk $E^+$ that Theorem~\ref{thm:criterion} associates with $S^+$ and $R^+$ is contained in $E$: besides the elements of $E$, it
contains $b^{\pm1}$, the prefixes $b^{-1},c^{-4},c^{-3},c^{-2},c^{-1}$ of $b^{-1}ac^4$, and the prefixes $b$, $ab=c^4$, $aba=c^4a$, $abab=c^2$ and
$ababa=c^2a$ of $b^2$, $(ab)^3$ and $ab$. These lie in $E$, because $ac^4$, $a^2c^4$, $a^2c^4a$ and $a^2c^4a^2c^4a$ are prefixes of $b^2$ and
$(ab)^3$ in Table~\ref{tab:certificate}, $c^2$ and $c^3$ are prefixes of the relations $[d,k]$ of the family \textsf{S}, $a^2=1$, and $c^6=1$.
Restricting a model of $E$ to $E^+$ gives a model of $E^+$, so $\sigma_E\ge\sigma_{E^+}$.

For $N\ge1$ take the $M=2^N$ words $T_u$ of Theorem~\ref{thm:pair-area}, one for each configuration $u$ supported in the first $N$ points of
ray~$1$. A commutator $[T_uxT_u^{-1},T_wyT_w^{-1}]$ with $x,y\in\{a,b\}$ contains the letter $b^{\pm1}$ at most four times, and replacing each
occurrence by $(ac^4)^{\pm1}$ costs one relation $b^{-1}ac^4$. So its area with respect to $R^+$ is at most $\Lambda^+_N=\Lambda_N+4$, and
Theorem~\ref{thm:criterion} gives
\begin{equation}\label{eq:criterion-applied}
\log\sigma_E(r)\ge2^N/16\qquad\text{whenever}\qquad r\ge4.7\cdot10^{9}\cdot2^N\Lambda^+_N .
\end{equation}
Since $a\ne1$, a $(1-1/r)$-expansive map needs at least two points, so $\log\sigma_E(r)\ge1$; for $r$ below any fixed bound $r_0$ claim (a)
therefore holds once $c$ is small enough, and we may assume $r$ large.

(a) By Theorem~\ref{thm:pair-area}, $\Lambda^+_N\le C_1N^\kappa$ with $C_1$ absolute. Let $r'=r/(4.7\cdot10^9C_1)$ and
$N=\lfloor\log r'-\kappa\log\log r'\rfloor$, which is at least $1$ for large $r$. Then
$2^N\Lambda^+_N\le C_1r'(\log r')^{-\kappa}(\log r')^\kappa=r/(4.7\cdot10^9)$, so~\eqref{eq:criterion-applied} applies, and
$\log\sigma_E(r)\ge2^N/16\ge r'/(32(\log r')^\kappa)$. Since $\log r'\le\log r$, this is~(a).

(b) Every element of $E$ is a generator, the inverse of one, or represented by a prefix of a word of length at most $24$ or by its inverse.
Proposition~\ref{prop:upper} with $\ell=24$ gives $\log\sigma_E(r)\le500\cdot24^2r\log(2+r)$, and $\log(2+r)\le2\log r$ for $r\ge2$. The same
proposition applies to every chunk.

(c) Extend $\phi$ to $S^+$ by $\phi(b)=\phi(a)\phi(c)^4$. Then $b^{-1}ac^4$ is mapped to $I$ exactly, and $b^2$ and $(ab)^3$ are mapped to the
images of the words $b^2$ and $(ab)^3$ of Table~\ref{tab:certificate}; so every word of $R^+$ is mapped within $e$ of $I$. With
$\Lambda^+_N\le C_1N^\kappa$ as in (a), the condition $K_\Delta\sqrt M\,\Lambda^+_Ne\le1$ of Theorem~\ref{thm:criterion} holds when
$2^NN^{2\kappa}\le Y:=(C_1K_\Delta e)^{-2}$, which is the case for $N=\lfloor\log Y-2\kappa\log\log Y\rfloor$; this $N$ is at least $1$ for $e$ small
in terms of $\Delta$. Then $2^N\ge Y/(2(\log Y)^{2\kappa})$, and since $C_1K_\Delta\ge1$ we have $\log Y\le2\log(1/e)$. Hence
\[
\log D\ \ge\ \frac{\Delta^2}{12}\,2^N\ \ge\ \frac{\Delta^2}{24\cdot4^\kappa C_1^2K_\Delta^2}\cdot\frac{e^{-2}}{(\log(1/e))^{2\kappa}}.
\]

(d) Proposition~\ref{prop:fp} gives finite presentation and Lemma~\ref{lem:structure} elementary amenability. By (a) the profile of $E$ is unbounded, so $E$ does not embed in a finite group and
$\Gamma$ is not LEF (\cite[Lemma~2.1]{DouglasMghirbiA}). This also follows because $\Gamma$ contains $H_3$, which is finitely presented and not
residually finite, since the finitary alternating group of an infinite set is simple; a finitely presented LEF group is residually finite.
\end{proof}

\begin{remark}\label{rem:lee}
Part~(a) uses the polynomial Dehn bound of Proposition~\ref{prop:dehn}(b). With Lee's exponential bound alone, Theorem~\ref{thm:pair-area} gives
$\Lambda^+_N\le2^{BN}$ for a constant $B$ that depends on the constant in Lee's inequality, and $N=\lfloor\log(r/(4.7\cdot10^9))/(1+B)\rfloor$
in~\eqref{eq:criterion-applied} gives only $\log\sigma_E(r)\ge c\,r^{1/(1+B)}$: a superpolynomial profile, with an exponent we cannot make
explicit. Transport at scale $N$ is too expensive for an exponential Dehn bound. The grid host of Remark~\ref{rem:grid} transports at scale
$\sqrt N$ and would reach, by the argument sketched there, $\log\sigma\ge c\,r\,e^{-C\sqrt{\log r}}$ with Lee's bound alone.
\end{remark}

The hash models are far smaller than any Følner set. For $r\ge1$ let $\textup{Føl}(r)$ be the least size of a finite nonempty set
$F\subset\Gamma$ with $|Fs\setminus F|\le|F|/r$ for $s\in\{p,q,t,c,a\}$.

\begin{proposition}\label{prop:folner-lower}
For $r\ge12$,
\[
2^{\,2^{\lfloor r/12\rfloor-1}}\ \le\ \textup{Føl}(r)\ \le\ 2^{\,2^{3r+14}} .
\]
So the Følner function of $\Gamma$ is doubly exponential, while by Theorem~\ref{thm:main-group}(b) the sofic profile of every chunk is
$\exp(O(r\log r))$.
\end{proposition}

\begin{proof}
\emph{Lower bound.} Let $F$ be such a set and $l=\lfloor r/12\rfloor\ge1$. For $v\in\{0,1\}^l$ the word
$T_v=t^{v_1}p\,t^{v_2}p\cdots p\,t^{v_l}p^{1-l}$ is freely equal to $\prod_{j=1}^lp^{j-1}t^{v_j}p^{1-j}$, so it represents the translation by
$\sum_jv_je_j$, and it has length less than $3l$. The elements $a_v=T_vaT_v^{-1}$ are the copies of $a$ at $2^l$ distinct configurations, so they
are commuting involutions and form a basis of an elementary abelian group $Q$ of order $2^{M}$, $M=2^l$. Each $a_v$ is a product of fewer than
$6l$ generators and inverses; since $|Fxy\setminus F|\le|Fx\setminus F|+|Fy\setminus F|$ and $|Fx^{-1}\setminus F|=|Fx\setminus F|$, we get
$|Fa_v\setminus F|\le6l|F|/r\le|F|/2$, and so $\sum_v|Fa_v\setminus F|\le M|F|/2$.

On the other hand, for every finite $S\subset Q$,
\begin{equation}\label{eq:cube}
\sum_v|Sa_v\setminus S|\ \ge\ |S|\,(M-\log|S|).
\end{equation}
Indeed, the left side is $M|S|-2e(S)$, where $e(S)$ is the number of pairs $\{x,xa_v\}\subset S$, and $e(S)\le\frac12|S|\log|S|$ by the
edge-isoperimetric inequality of the cube~\cite{Harper1964}. We recall the proof, by induction on $M$: splitting $S$ by the last coordinate into
$S_0,S_1$ of sizes $s_0\le s_1$ gives $e(S)\le e(S_0)+e(S_1)+s_0\le\frac12(s_0\log s_0+s_1\log s_1)+s_0$, which is at most $\frac12s\log s$,
$s=s_0+s_1$, because the binary entropy function $h$ satisfies $h(x)\ge2x$ for $0\le x\le\frac12$. The right multiplications by the $a_v$ preserve each coset $xQ$;
applying~\eqref{eq:cube} to the sets $x^{-1}(F\cap xQ)$ and using $|F\cap xQ|\le|F|$ gives $\sum_v|Fa_v\setminus F|\ge|F|(M-\log|F|)$. Hence
$\log|F|\ge M/2=2^{l-1}$.

\emph{Upper bound.} Let $m=\lceil(r-1)/2\rceil$, $R=2m+3$ and $F=\bigcup_{z\in B_m}K_Rg_z$, a disjoint union because $\pi(g_z)=z$. By
Lemma~\ref{lem:clock}, $g_zs\in K_Rg_{z+\pi(s)}$ for $z\in B_m$ and each generator $s$, so $|Fs\setminus F|\le|K_R|(2m+1)=|F|/(2m+1)\le|F|/r$. With
$D=|Y_R|=6m+9$ we have $|K_R|=6^{2^D}\cdot2^D\cdot|\GL_D(\F_2)|$, so $\log|F|\le2^D\log6+D+D^2+2\log(2m+1)\le2^{D+2}\le2^{3r+14}$.
\end{proof}

\subsection{The grid host}\label{sec:grid}

\begin{remark}\label{rem:grid}
Placing the lamps on configurations of the cells of a grid, rather than of the points of $Y$, should give a stronger lower bound with the same upper
bound, and a nearly linear exponent $r^{1-o(1)}$ even from Lee's exponential Dehn bound alone. We only sketch the proofs, and the conclusions of this remark are conditional on them. Let $I=Y\times Y$, let $T=H\times H$ act on $I$ coordinatewise, and let
\[
\Gamma_2=\Bigl(\bigoplus_{v\in V_2}S_3\Bigr)\rtimes\bigl(V_2\rtimes(\GL_{\rm fin}(I)\rtimes T)\bigr),\qquad V_2=\F_2^{(I)} .
\]
It is locally finite-by-$\Z^4$ and generated by seven elements (commutators $[E_{x\leftarrow y},E_{y\leftarrow z}]=E_{x\leftarrow z}$ of
transvections of the third kind give those within a row or a column): the generators $p_1,q_1,p_2,q_2$ of the two copies of $H$, the translation $t$ by
the cell $o=((1,1),(1,1))$, the element $c=d\,ab$, where $d$ is the transvection that adds $e_o$ to $e_{((1,2),(1,2))}$, and $a$. A certificate of
$77$ relations, of total length $5445$ and maximum length $224$, is listed in \texttt{verification/relators.txt} and checked by
\texttt{verification/relators\_check.py}. It has the families of Table~\ref{tab:certificate}, with two copies of the Houghton relators and
their commutation, and with stabilizer and action relations for three kinds of transvection: within a row, within a column, and neither.

The gain is in the scale of transport. $N$ cells fit in a square of side $\lceil\sqrt N\,\rceil$, so the transporters have length $O(\sqrt N)$, and
the local fillings cost $C(\delta_H(C\sqrt N)+N)$; the term $N$ counts the swaps of letters from the two copies of $H$ in a change of transporter.
The clearing argument of Theorem~\ref{thm:pair-area} then bounds the pair areas by $CN^2(\delta_H(C\sqrt N)+N)$, which is at most $CN^{2+p_H/2}$ by
Proposition~\ref{prop:dehn}(b) and at most $e^{C\sqrt N}$ by Proposition~\ref{prop:dehn}(a). As in the proof of Theorem~\ref{thm:main-group}, the chunk
$E_2$ of this certificate should therefore have
\[
\log\sigma_{E_2}(r)\ \ge\ c\,\frac{r}{(\log r)^{2+p_H/2}},\qquad 2+\tfrac{p_H}2<7.38,
\]
and, using only Lee's bound, $\log\sigma_{E_2}(r)\ge c\,r\,e^{-C\sqrt{\log r}}$; unitary models need $\log D\ge c_\Delta e^{-2}/(\log(1/e))^{4+p_H}$.
In the other direction, hashing the configurations directly, as in Proposition~\ref{prop:upper}, would need a configuration space of dimension
of order $r^2$. Hashing the cells first avoids this: a cell $(y_1,y_2)$ gets the label $a_1(f_1^{-1}y_1)+a_2(f_2^{-1}y_2)\in\F_q$, where the $f_j$ are
positions in finite F{\o}lner charts of $H$ of size $e^{O(r\log r)}$ (inverses of right F{\o}lner sets, since the charts move by left
multiplication) and the $a_j$ are random tables on $O(r)$ points; this label is covariant under the
charts and separates the finitely many cells of a chunk with probability $1-O(1/q)$. The configurations are then compressed linearly to
$\F_2^{\F_q}$ and hashed affinely as before. With $q$ of order $\ell^2r$, for a chunk of words of length at most $\ell$, and the other parameters of order $r$, this would give $\log\sigma_{E'}(r)\le C_{E'}r\log(2+r)$ for
every chunk $E'$ of $\Gamma_2$, and the same for the $d$-dimensional grid host, built in the same way on the cells of $Y^d$ with $H^d$ acting,
for every fixed $d$. So $\sigma_{E_2}(r)$ would be $\exp(r^{1+o(1)})$ as well, and no fixed-dimensional
grid host would exceed $\exp(r^{1+o(1)})$. The proofs follow those above with the changes indicated, and we omit them.
\end{remark}

%======================================================================
\section{Quantum channels that need nearly linear memory}\label{sec:channel}
%======================================================================

A finite set of relations becomes a quantum channel of fixed dimension: each relation contributes a few small unitary blocks, and the channel
evaluates the blocks in the regular representation of the group and traces out the group von Neumann algebra. In~\cite{DouglasMghirbiB} we showed
that the memory needed to apply such a channel repeatedly, releasing each output before the next input arrives, is bounded below by the size of
approximate representations of the relations. This section sharpens that link, so that it loses a single square root, and applies it to the
configuration lamps. The result is an explicit channel, in the closure of finite tracial channels, whose repeated use with logarithmic purity
needs memory $n^{1+o(1)}$: at least $n/(\log n)^{O(1)}$ for $n$ uses (Theorem~\ref{thm:main-channel}), and at most $O(n\log n)$, because the sofic
models of $\Gamma$ also give devices (Proposition~\ref{prop:sofic-upper}). More precisely, memory and purity trade at a fixed rate: a device
allowed purity $B$ needs memory about $n/B$, up to polylogarithmic factors, and exact devices attain this. At the optimal exchange rate the memory
is $n^{1/2+o(1)}$, attained by devices that keep one sofic model for all rounds (Theorem~\ref{thm:streaming}).

\subsection{Setting}\label{sec:channel-setting}

We recall the framework of~\cite{DouglasMghirbiB} and keep its notation.

\paragraph{Channels.} Channels act on $M_d=M_d(\mathbb C)$. The normalized Choi state of a channel $\Psi$ is
$\mathcal J(\Psi)=(\Psi\otimes\mathrm{id})(|\Omega\rangle\langle\Omega|)$, with $|\Omega\rangle=d^{-1/2}\sum_i|ii\rangle$, and
$\ddia(\Psi,\Psi')=\frac12\|\Psi-\Psi'\|_\diamond$ is the diamond distance. We write $\tau_M$ for the normalized trace on $M_M$. A channel $\Psi$ has a
\emph{finite tracial factorization} with bath dimension $M$ if $\Psi(X)=(\mathrm{id}\otimes\tau_M)[U(X\otimes I_M)U^*]$ for a unitary $U$ on
$\mathbb C^d\otimes\mathbb C^M$, and $\cK_d$ denotes the closure of the set of channels on $M_d$ with a finite tracial factorization of some bath
dimension. It is also the closure of the convex hull of that set~\cite[Proposition~3.8(ii)]{Musat2025}, and it consists of unital channels. Channels
with a finite tracial factorization are exactly the noisy operations of~\cite{HHO2003} from $M_d$ to $M_d$, and $\cK_d$ is their closure.

\paragraph{Causal services.} A \emph{causal service} of a channel $\Phi$ on $M_d$ for $n$ rounds consists of a memory register of $Q$ qubits in a fixed
initial state $\beta_0$ and, for each round, a fixed unitary on the input and the memory, after which the $d$-dimensional output is released; after
each round but the last, a prescribed set of memory qubits is exported for good and replaced by the same number of imported qubits in a fixed
state (possibly correlated within one step), independent of everything else. An observer chooses the inputs adaptively, possibly entangled with references of its own, and the
\emph{error} of the service is the largest half trace distance between the observer's final state in this experiment and in the ideal experiment
with $n$ independent uses of $\Phi$. The resources are the memory $Q$, the \emph{exchange} $J$ (the total number of exported qubits), and the
\emph{purity} $P+C$, where $P=Q-S(\beta_0)$ and $C$ is the sum over the imports of their number of qubits minus their entropy. Every retained
qubit counts; there is no free reset, measurement record or postselection.

\paragraph{Gadget channels.} Let $G$ be a group generated by a finite set $S$, and let $R$ be a finite set of words in $S^{\pm1}$, each of length at
least two and each representing $1$ in $G$. Write $r_i=\sigma_{i,1}\cdots\sigma_{i,\ell_i}$ ($1\le i\le r_0$) for the words of $R$, and put
$s=|S|$ and $\Lambda=\sum_i\ell_i$. The \emph{symbols} are the $2s$ letters $x^{\pm1}$, $x\in S$, and the proper prefixes
$p_{i,l}=\sigma_{i,1}\cdots\sigma_{i,l}$, $2\le l\le\ell_i-1$, one for each pair $(i,l)$; we also put $p_{i,1}=\sigma_{i,1}$ and let $p_{i,\ell_i}$
be the empty word. The \emph{triangles} are the triples $(p_{i,l-1},\sigma_{i,l},p_{i,l})$, $2\le l\le\ell_i$, and $(x,x^{-1},\varnothing)$,
$x\in S$. For unitaries $\lambda(w)$ assigned to the symbols, with $\lambda(\varnothing)=I$, let
\[
U_\lambda=\mathrm{diag}\bigl(I,(\lambda(w))_{w}\bigr)\ \oplus\ \bigoplus_{(x,y,z)}\frac1{\sqrt2}\begin{pmatrix}I&\lambda(y)\\ \lambda(x)&-\lambda(x)\lambda(y)\end{pmatrix},
\]
where $w$ runs over the symbols and $(x,y,z)$ over the triangles. Each triangle block equals $\mathrm{diag}(I,\lambda(x))\,(H\otimes I)\,\mathrm{diag}(I,\lambda(y))$
with $H$ the Hadamard matrix, so $U_\lambda$ is unitary for all unitary $\lambda$. There are $2s+\Lambda-2r_0$ symbols and $\Lambda-r_0+s$
triangles, so the system dimension is
\begin{equation}\label{eq:gadget-dimension}
d=1+(2s+\Lambda-2r_0)+2(\Lambda-r_0+s)=1+4s+3\Lambda-4r_0 .
\end{equation}
The \emph{gadget channel} $\Phi_G$ on $M_d$ evaluates each symbol $w$ in the left regular representation, $\lambda(w)=\lambda([w])$ with $[w]\in G$ the
element that $w$ represents, and traces out the group von Neumann algebra $L(G)$ with its canonical trace $\tau$:
$\Phi_G(X)=(\mathrm{id}\otimes\tau)[U_\lambda(X\otimes1)U_\lambda^*]$. Since the relations hold, the lower right entry of the block of a triangle
$(x,y,z)$ is $-\lambda([z])$.

We write the gadget as $U_\lambda=\sum_{a\in\mathcal T}c_aE_a\otimes\lambda([w_a])$, where $\mathcal T$ is the set of \emph{slots} (the nonzero
entries), $E_a$ is the matrix unit of slot $a$, $c_a\in\{\pm1,\pm2^{-1/2}\}$, and the word $w_a$ is empty, a symbol, or the concatenation $xy$ in the
lower right corner of a triangle. The elements $[w_a]$ are the \emph{labels}. With $A_g=\sum_{[w_a]=g}c_aE_a$, the matrices $A_g$ have disjoint
supports, $\Phi_G(X)=\sum_gA_gXA_g^*$ and $\mathcal J(\Phi_G)=\frac1d\sum_g|A_g\rangle\!\rangle\langle\!\langle A_g|$, because
$\tau(\lambda(g)\lambda(h)^*)=\delta_{gh}$. So the Choi rank $r_*$ of $\Phi_G$ is the number of distinct labels. Every label is carried by an
\emph{anchor}: a diagonal slot with coefficient $1$ whose row and column contain no other slot, namely the first diagonal entry for the label $1$
and a symbol entry otherwise, since the lower right entry of a triangle carries the label of a prefix symbol or $1$. Feeding one half of a maximally entangled
state on the $r_*$ anchor positions gives a pure probe whose output is flat of rank $r_*$, and $\log r_*$ is the largest entropy of a complementary
output of $\Phi_G$~\cite[Proposition~9.1]{DouglasMghirbiB}.

\paragraph{One round.} Let $U$ be a unitary on $\mathbb C^d\otimes\mathbb C^M$ and $\rho$ a state on $\mathbb C^M$, and put
\[
\Psi(X)=\mathrm{Tr}_M\,U(X\otimes\rho)U^*,\qquad T(\rho)=\mathrm{Tr}_d\,U(\tau_d\otimes\rho)U^* .
\]
Relative to a channel $\Phi$ with Choi support projection $\Pi_\Phi$, the \emph{support leakage} of $\Psi$ is
$\ell=\Tr[(I-\Pi_\Phi)\mathcal J(\Psi)]$. The \emph{entropy production} of the round is $\Delta S=S(T(\rho))-S(\rho)$, which is nonnegative because
$T$ is unital. A service produces such rounds: the next result, from~\cite{DouglasMghirbiB}, finds one in which both quantities are small.

\begin{theorem}[Sequential extraction~\cite{DouglasMghirbiB}]\label{thm:extraction}
Let a causal service of a channel $\Phi$ on $M_d$ for $n$ rounds have memory $Q$, error $\epsilon<1/2$ and $P+C\le B$, and put
\[
t_n=\frac{B+\epsilon+h_2(\epsilon)+\epsilon\log n}{n(1-\epsilon)} .
\]
Then some round $t$, with its unitary $U_t$ and the memory state $\rho_t$ entering it, defines a channel $\Psi_t(X)=\mathrm{Tr}_M\,U_t(X\otimes\rho_t)U_t^*$
with bath dimension $M=2^Q$, $\ddia(\Psi_t,\Phi)\le\epsilon/(1-\epsilon)$, and support leakage $\ell_t$ and entropy production $\Delta S_t$ satisfying
$\ell_t+\Delta S_t\le t_n$.
\end{theorem}

Since the bath has $Q$ qubits, $S(\rho_t)\le Q$. Every lower bound on memory below is a lower bound on the entropy of a memory state that realizes
one good round.

\subsection{The weighted one-round theorem}\label{sec:weighted-statement}

The following theorem is the engine of this section. It reads an approximate representation of the relations directly off the memory state of
one round, weighted by that state, and turns many cheaply commuting copies of $S_3$ into entropy of the memory. Transport words $u_1,\dots,u_N$
over $S^{\pm1}$ have \emph{length at most $L$} and \emph{pairwise areas at most $A$} if $|u_i|\le L$ for all $i$ and
\[
\Area_R\bigl([u_ixu_i^{-1},\,u_kyu_k^{-1}]\bigr)\le A\qquad(i\ne k,\ x,y\in\{a,b\}),
\]
where $\Area_R$ is the area in the free group on $S$, as in Section~\ref{sec:criterion}.

\begin{theorem}[Weighted one-round theorem]\label{thm:weighted}
Let $G$, $S$ and $R$ be as above, and suppose that $S$ contains $a$, $b$ and $j$, that $a^2$, $b^2$, $(ab)^3$ and $j^{-1}ab$ lie in $R$, and that
$j\ne1$ in $G$. Put $m=\max\{6,\ell_1,\dots,\ell_{r_0}\}$. Let $u_1,\dots,u_N$ be transport words of length at most $L$ and pairwise areas at most
$A$. Let $\Psi(X)=\mathrm{Tr}_M\,U(X\otimes\rho)U^*$ be one round with $\ddia(\Psi,\Phi_G)\le1/8$, support leakage $\ell$ relative to $\Phi_G$ and
entropy production $\Delta S$, and put $\nu=\ell+\Delta S$. If
\begin{equation}\label{eq:weighted-condition}
2^{25}\,m^2\sqrt d\,\sqrt N\,(A+L+1)\sqrt\nu\ \le\ 1,
\end{equation}
then $S(\rho)\ge N/16$.
\end{theorem}

Two features matter. First, the relator defect that the round reveals is of order $\sqrt\nu$, a single square root of the quantity that
Theorem~\ref{thm:extraction} controls, so a service with purity $O(\log n)$ is read at defect about $n^{-1/2}$. Second, there is no
tracialization: the argument works with the actual memory state $\rho$, which need not be flat. Replacing the round by a nearby finite tracial
mixture, as the general correspondence of~\cite{DouglasMghirbiB} does, costs a second square root, so that correspondence reads the profile at
defect about $n^{-1/4}$ instead of $n^{-1/2}$. The weighted route was
carried out in~\cite{DouglasMghirbiB} for a single gadget built from Houghton's group; Theorem~\ref{thm:weighted} holds for every relation
set that contains the relations of $S_3$ and a name for $ab$. It is proved in Section~\ref{sec:weighted-proof}.

\begin{remark}[The constant term]\label{rem:plus-one}
The term $+1$ in~\eqref{eq:weighted-condition} cannot be dropped. Take $N=1$, the empty transport word, so that $L=0$, and $A=0$, which is
allowed since there are no pairs. The gadget channel has a Stinespring dilation with a pure environment of dimension $r_*$, so there is a round
with $M=r_*$, $\rho$ pure and $\Psi=\Phi_G$ exactly. Then $\ddia(\Psi,\Phi_G)=0$ and the condition with $A+L$ in place of $A+L+1$ holds, since its left
side vanishes, but $S(\rho)=0<1/16$. With the $+1$, the condition forces $\nu\le2^{-50}m^{-4}d^{-1}$, while this round has $\ell=0$ and
$\Delta S=S(T(\rho))$, the entropy of the complementary output at the maximally mixed input. That output is diagonal with entries
$\|A_g\|_F^2/d\ge1/d$, and there are at least two labels, since $ab\ne1$; so $\Delta S\ge h_2(1/d)$, far above the threshold.
\end{remark}

Applied to a family of transport words that grows with a parameter, Theorem~\ref{thm:weighted} gives a lower bound on the entropy of one round
in terms of $\nu$ alone.

\begin{corollary}[Amplification]\label{cor:amplification}
Keep the assumptions of Theorem~\ref{thm:weighted} on $G$, $S$, $R$ and the round, and let $K\ge1$ and $\kappa\ge0$.
\begin{enumerate}[label=(\alph*),itemsep=1pt]
\item Suppose that for every integer $k\ge1$ there are $2^k$ transport words of length at most $L_k$ and pairwise areas at most $A_k$, with
$A_k+L_k+1\le Kk^\kappa$. Put $C_R=2^{50}m^4dK^2$. There is $y_0$ depending only on $\kappa$ such that
\[
S(\rho)\ \ge\ \frac{y}{32\,(\log y)^{2\kappa}}\qquad\text{for every }y\ge y_0\text{ with }C_R\,\nu\,y\le1 .
\]
\item Suppose instead that $A_k+L_k+1\le2^{K\sqrt k}$ for every $k\ge1$, and put $C'_R=2^{50}m^4d$. There is $y_1$ depending only on $K$ such that
\[
S(\rho)\ \ge\ \frac{y}{32}\,2^{-2K\sqrt{\log y}}\qquad\text{for every }y\ge y_1\text{ with }C'_R\,\nu\,y\le1 .
\]
\end{enumerate}
\end{corollary}

\begin{proof}
(a) Let $y_0\ge4$ be such that $k=\lfloor\log y-2\kappa\log\log y\rfloor\ge1$ for $y\ge y_0$. Then $k\le\log y$ and $2^k\le y(\log y)^{-2\kappa}$, so
$2^kk^{2\kappa}\le y$. With $N=2^k$, the left side of~\eqref{eq:weighted-condition} is at most
$2^{25}m^2\sqrt d\,K\,(2^kk^{2\kappa})^{1/2}\sqrt\nu\le(C_R\nu y)^{1/2}\le1$. So $S(\rho)\ge2^k/16$, and $2^k>y(\log y)^{-2\kappa}/2$.

(b) Let $y_1$ be such that $k=\lfloor\log y-2K\sqrt{\log y}\rfloor\ge1$ for $y\ge y_1$. Then $K\sqrt k\le K\sqrt{\log y}$, so
$2^k(A_k+L_k+1)^2\le2^{k+2K\sqrt{\log y}}\le y$, and~\eqref{eq:weighted-condition} holds with $N=2^k$ as in (a). Hence
$S(\rho)\ge2^k/16>y\,2^{-2K\sqrt{\log y}}/32$.
\end{proof}

\subsection{The channel of the configuration lamps}\label{sec:lamp-channel}

Theorem~\ref{thm:weighted} needs $a$, $b$ and a name for $ab$ among the generators. Let $R$ be the certificate of Table~\ref{tab:certificate}, in
the five generators $p,q,t,c,a$, with $b$ standing for the word $ac^4$. Let $S_\Phi=\{p,q,t,c,a,b,j\}$, where the new letters represent
$b=ac^4$ and $j=ab$ in $\Gamma$, and let $R_\Phi$ consist of the $25$ words of $R$ and the four words
\[
b^{-1}ac^4,\qquad b^2,\qquad (ab)^3,\qquad j^{-1}ab
\]
in the letters of $S_\Phi$. They represent $1$ in $\Gamma$, $a^2\in R$, and $j\ne1$ because $ab$ has order three (Lemma~\ref{lem:structure}). The
channel of the configuration lamps is the gadget channel $\Phi=\Phi_\Gamma$ of $R_\Phi$.

\paragraph{Dimension and rank.} Here $s=7$, $r_0=29$, and the total length is $\Lambda=303+6+2+6+3=320$, so~\eqref{eq:gadget-dimension} gives
\[
d=1+4\cdot7+3\cdot320-4\cdot29=873 ,
\]
with $276$ symbols and $298$ triangles; the longest word has length $24$, so $m=24$. The Choi rank $r_*$ is the number of distinct labels, at most
$1+276=277$. It is smaller: each of the $26$ words of $R_\Phi$ of length at least three has a longest proper prefix that represents the inverse of
its last letter, so these $26$ prefix symbols carry labels already carried by letters, and $a^2=b^2=t^2=1$ identifies the labels of
$a^{-1},b^{-1},t^{-1}$ with those of $a,b,t$. Hence $2\le r_*\le248$, where the lower bound holds because $1$ and $j$ are distinct labels.

\paragraph{Transport.} For $N\ge1$, the proof of Theorem~\ref{thm:main-group} provides $2^N$ words $T_u$, one for each configuration $u$
supported in the first $N$ points of ray~$1$, with $|T_u|\le21N\max(3,N)\le63N^2$ (Theorem~\ref{thm:pair-area}) and
$\Area_{R_\Phi}([T_uxT_u^{-1},T_wyT_w^{-1}])\le\Lambda_N+4\le C_1N^{\kappa}$ for $u\ne w$ and $x,y\in\{a,b\}$, where
\[
p_H=\log107+4,\qquad \kappa=2+p_H<12.75,
\]
and $C_1$ is absolute; the letter $b$ is converted to $ac^4$ by the word $b^{-1}ac^4$, at most four times, and areas with respect to
$R_\Phi\supseteq R$ are no larger than with respect to $R$. So the hypothesis of Corollary~\ref{cor:amplification}(a) holds with this $\kappa$ and
$K=C_1+64$.

\begin{theorem}[Theorem B]\label{thm:main-channel}
The channel $\Phi$ on $M_{873}$ lies in $\cK_{873}$, and its Choi rank satisfies $2\le r_*\le248$. Let $\kappa=2+p_H=6+\log107$. There are
constants $c,c_0,K>0$ and $n_0$ such that the following hold; all services have error at most $1/16$.
\begin{enumerate}[label=(\alph*),itemsep=1pt]
\item \emph{Memory against purity.} For $n\ge n_0$ and $\log n\le B\le c_0n$, every causal service of $\Phi$ for $n$ rounds with purity
$P+C\le B$ has, at any exchange,
\[
Q\ \ge\ c\,\frac{n}{B\,(\log n)^{2\kappa}} .
\]
For every $n\ge1$ and $0<B\le n$ there is an exact causal service, with error $0$, purity $P+C<B$, exchange $J=(n-1)Q$ and memory
$Q\le K(n/B)\log(2n/B)$. In particular, for every $A_0\ge0$ there is $c'>0$ such that services with $P+C\le A_0\log n$ have
$Q\ge c'n/(\log n)^{2\kappa+1}$ for all large $n$, where $2\kappa+1=2p_H+5<26.49$, and there are services with $P=C=0$ and $Q\le Kn\log n$.
\item \emph{The rank rate.} For $n\ge n_0$, every causal service of $\Phi$ for $n$ rounds at the rank rate, $J\le\frac n2\log r_*$, has, whatever
its purity,
\[
Q\ \ge\ c\,\frac{\sqrt n}{(\log n)^{\kappa}} .
\]
For every $n\ge2$ there is a causal service at the rank rate with $C=0$ and $P,Q\le K\sqrt{n\log n}$.
\end{enumerate}
So memory and purity trade at a fixed rate: for $n\ge n_0$ and $\log n\le B\le c_0n$ the least memory of a service with purity at most $B$ lies
between $c\,n/(B(\log n)^{2\kappa})$ and $K(n/B)\log(2n/B)$, and the product of memory and purity is $n^{1+o(1)}$. At the rank rate the least
memory is $n^{1/2+o(1)}$, and it is attained with purity of the same order: the rank rate sits where memory and purity are equal.
\end{theorem}

The exponent $2\kappa$ in (a) is the unitary exponent of Theorem~\ref{thm:main-group}(c): a service with purity $B$ is read at relator defect about
$(B/n)^{1/2}$. The lower bound in (b) needs no assumption on purity, because at the rank rate purity is paid for with memory
(Proposition~\ref{prop:purity-budget}): there $P+C<3Q+6$, and the argument of (a) with $B$ of order $Q$ gives $Q^2\gtrsim n$ up to logarithms. The
exact services in (a) refresh the whole memory in every round, so their exchange is far above the rank rate; the services in (b) keep a sofic
model in memory for all $n$ rounds and pay for its bad points with purity (Theorem~\ref{thm:streaming}). We prove Theorem~\ref{thm:main-channel} in Section~\ref{sec:main-channel-proof}, after the purity budget
and the upper bound.

\begin{remark}[The grid host]\label{rem:grid-channel}
The grid host $\Gamma_2$ of Remark~\ref{rem:grid} should give a stronger lower bound with the same upper bound; both rest on sketches. Its certificate has $7$ generators and $77$
relations of total length $5445$ and maximum length $224$. Adding $b$ and $j$ and the same four words gives $s=9$, $r_0=81$ and $\Lambda=5462$, so
its gadget channel $\Phi_2$ acts on $M_d$ with
\[
d=1+4\cdot9+3\cdot5462-4\cdot81=16099,
\]
has $m=224$, and has Choi rank at most $1+18+5300-78-3=5238$ by the count above ($78$ of its $81$ words have length at least three). With the
transport of Remark~\ref{rem:grid}, words of length $O(N^{3/2})$ and pair areas $O(N^2(\delta_H(C\sqrt N)+N))$, the hypothesis of
Corollary~\ref{cor:amplification}(a) holds with $\kappa=2+p_H/2$, and part (b) of that corollary holds with Lee's exponential Dehn
bound~\cite{Lee2012} alone. The proof below then gives, for $\Phi_2$: memory at least $c\,n/(\log n)^{p_H+5}$ at purity $O(\log n)$, where
$p_H+5<15.75$; memory at least $n\,2^{-C\sqrt{\log n}}$ at purity $O(\log n)$ using only Lee's bound; and memory at least
$c\sqrt n/(\log n)^{2+p_H/2}$ at the rank rate, whatever the purity. These rest on the pair-area bound for $\Gamma_2$, whose proof
Remark~\ref{rem:grid} only sketches. By the two-stage hashing sketched in Remark~\ref{rem:grid}, $\log\sigma(r)=O(r\log r)$ for $\Gamma_2$ as well, and so an upper bound of order
$n\log n$. For $\Gamma$ itself, Lee's bound alone gives pair areas $2^{O(N)}$, and the argument gives only memory $n^{c}$ for some $c>0$.
\end{remark}

\subsection{A purity budget at the rank rate}\label{sec:purity-budget}

At the rank rate a device cannot buy memory with purity: the purity it can consume is linear in its memory. We sharpen the budget $14Q+3$
of~\cite{DouglasMghirbiB} to $3Q+6$ by combining its rank count with a one-sided second-moment bound.

\begin{proposition}[Purity budget]\label{prop:purity-budget}
Let $\Phi$ be a channel on $M_d$, and suppose that some pure state $\psi_*$ of the input and a reference has output
$\omega=(\Phi\otimes\mathrm{id})(\psi_*)$ flat of rank $r$. Every causal service of $\Phi$ for $n$ rounds with error $\epsilon\le1/16$ and exchange
$J\le\frac n2\log r$ satisfies
\[
P+C\ <\ 3Q+6 .
\]
\end{proposition}

Every gadget channel has such a probe with $r=r_*$ (Section~\ref{sec:channel-setting}). The proof uses the variance of the surprisal of a
state with small deficit, which is~\cite[Lemma~5.9]{DouglasMghirbiB}; we include the short argument.

\begin{lemma}[{Surprisal variance~\cite[Lemma~5.9]{DouglasMghirbiB}}]\label{lem:surprisal}
Let $\beta$ be a state on $q\ge1$ qubits with spectrum $(\beta_i)$ and deficit $D=q-S(\beta)$. The surprisal $-\log\beta_i$, drawn with probability
$\beta_i$, has variance at most $(q+2/\ln2)\,D$.
\end{lemma}

\begin{proof}
Put $\varphi(x)=x\ln x-x+1\ge0$. We claim $x(\ln x)^2\le(2+\ln_+x)\varphi(x)$ for $x\ge0$. For $x=e^t$ with $t\ge0$ the difference of the two sides
is $f(t)=e^t(t-2)+t+2$, with $f(0)=f'(0)=0$ and $f''(t)=te^t\ge0$. For $t\le0$ it is $g(t)=2(te^t-e^t+1)-t^2e^t$, with $g(0)=0$ and
$g'(t)=-t^2e^t\le0$; and $x=0$ is clear. Now let $Q_0=2^q$ and $x_i=Q_0\beta_i\le Q_0$, the index running over all $Q_0$ eigenvalues. Then
$\frac1{Q_0}\sum_ix_i=1$ and $D\ln2=\sum_i\beta_i\ln(Q_0\beta_i)=\frac1{Q_0}\sum_i\varphi(x_i)$. The variance is at most
\[
\sum_i\beta_i\bigl(\log(Q_0\beta_i)\bigr)^2=\frac1{Q_0(\ln2)^2}\sum_ix_i(\ln x_i)^2\le\frac{2+q\ln2}{(\ln2)^2}\,D\ln2=\Bigl(q+\frac2{\ln2}\Bigr)D .
\]
\end{proof}

\begin{proof}[Proof of Proposition~\ref{prop:purity-budget}]
Put $b=P+C$ and suppose, for a contradiction, that $b\ge3Q+6$.

\emph{Preloading.} The imports are fixed states, independent of everything else, and the exported qubits are never used again. So the observer's
final state does not change if every import is prepared at the start in a register of its own, and each replacement swaps it with the exported
qubits, which then stay in the device. The device becomes a register of $L=Q+J$ qubits in the product state
$\sigma_0=\beta_0\otimes\bigotimes_f\omega_f$ of the initial memory state and the imported states, and it applies a fixed unitary in each round.

\emph{The surprisal.} Let $(\sigma_i)$ be the spectrum of $\sigma_0$ and let $Z=L+\log\sigma_i$, with $i$ drawn with probability $\sigma_i$. Then
$\mathbb EZ=L-S(\sigma_0)=b$. Also $Z$ is the sum of the independent variables $Z_f=q_f+\log\beta^{(f)}_i$ attached to the factors of $\sigma_0$,
where $q_f$ is the number of qubits of the factor and $\beta^{(f)}$ its spectrum. Every factor has at most $Q$ qubits, since an import replaces
exported memory qubits, and the deficits $b_f=\mathbb EZ_f$ of the factors add up to $b$. By Lemma~\ref{lem:surprisal},
\[
\mathrm{Var}\,Z=\sum_f\mathrm{Var}\,Z_f\le\Bigl(Q+\frac2{\ln2}\Bigr)b .
\]

\emph{A tail bound from the rank.} Let the observer feed the input halves of $n$ copies of $\psi_*$ and keep the references. Its final state is
$\tau_O=\sum_i\sigma_i\tau_i$, where $\tau_i$ is its final state when the register starts in the eigenvector of $\sigma_i$. Each $\tau_i$ is a
marginal of a pure state of the observer's system and the $L$-qubit register, so $\rank\tau_i\le2^L$. Fix a real number $k$ and let $I_k$ be the
set of indices with $\sigma_i\ge2^{k-L}$, so that $|I_k|\le2^{L-k}$. Let $E$ be the support projection of $\sum_{i\in I_k}\tau_i$. Then
$\rank E\le2^{2L-k}$ and $\Tr(E\tau_O)\ge\sum_{i\in I_k}\sigma_i=\Pr(Z\ge k)$. The ideal final state $\omega^{\otimes n}$ is flat of rank $r^n$, so
$\Tr(E\omega^{\otimes n})\le2^{2L-k}r^{-n}\le2^{2Q-k}$, using $2J\le n\log r$. Since the error is at most $\epsilon$,
\begin{equation}\label{eq:rank-tail}
\Pr(Z\ge k)\ \le\ \epsilon+2^{2Q-k}\qquad(k\in\mathbb R).
\end{equation}

\emph{The contradiction.} Put $u=b-2Q-3$, so that $u\ge b/3+1>0$. By the one-sided Chebyshev inequality,
$\Pr(Z\ge b-u)\ge u^2/(\mathrm{Var}\,Z+u^2)$, while~\eqref{eq:rank-tail} with $k=b-u=2Q+3$ gives $\Pr(Z\ge b-u)\le\frac1{16}+\frac18=\frac3{16}$. Hence
$13u^2\le3\,\mathrm{Var}\,Z$, and with $u>b/3$,
\[
\frac{13}9\,b^2<3\Bigl(Q+\frac2{\ln2}\Bigr)b,\qquad\text{that is}\qquad b<\frac{27}{13}\Bigl(Q+\frac2{\ln2}\Bigr)<2.08\,Q+5.993 ,
\]
which contradicts $b\ge3Q+6$.
\end{proof}

\subsection{An upper bound from sofic models}\label{sec:sofic-upper}

Sofic models of the labels give devices. The point is that the triangle blocks can be made exactly unitary for arbitrary permutations, and that at
most points of a model the pure bath vector reproduces the target channel exactly, not approximately. Recall from Section~\ref{sec:criterion}
that $\sigma_E(r)$ is the least $M$ for which some $f:E\to\Sym(\Omega)$, $|\Omega|=M$, with $f(1)=\mathrm{id}$, is a
\emph{$(1-r^{-1})$-expansive $r^{-1}$-morphism}: $f(x)f(y)$ and $f(xy)$ differ on at most $M/r$ points whenever $x,y,xy\in E$, and $f(x)$ and
$f(y)$ agree on at most $M/r$ points whenever $x\ne y$ in $E$.

\begin{proposition}[Sofic models give devices]\label{prop:sofic-upper}
Let $G$, $S$ and $R$ be as in Section~\ref{sec:channel-setting}, with $t_\triangle$ triangles and $r_*$ labels, let $E$ be a chunk of $G$ containing
all labels, and put $K_G=t_\triangle+\binom{r_*}2$. Every $(1-r^{-1})$-expansive $r^{-1}$-morphism $f:E\to\Sym(\Omega)$, $|\Omega|=M$, gives a unitary $U_f$ on
$\mathbb C^d\otimes\mathbb C^\Omega$ whose tracial channel $\Phi_f(X)=(\mathrm{id}\otimes\tau_M)[U_f(X\otimes I)U_f^*]$ satisfies
\[
\Phi_f=(1-\vartheta)\Phi_G+\vartheta\,\Xi
\]
for some channel $\Xi$ and some $0\le\vartheta\le K_G/r$. Consequently, for all $n\ge1$ and $0<\epsilon\le1$ with $\sigma_E(2K_Gn/\epsilon)<\infty$,
$\Phi_G$ has
\begin{enumerate}[label=(\alph*),itemsep=1pt]
\item a causal service with error at most $\epsilon$, $P=C=0$, $J=(n-1)Q$ and $Q=\bigl\lceil\log\sigma_E(2K_Gn/\epsilon)+\log(2n/\epsilon)\bigr\rceil$;
\item a causal service with error at most $\epsilon$, $J=P=C=0$ and $Q=\bigl\lceil n\log\sigma_E(2K_Gn/\epsilon)+\log(2/\epsilon)\bigr\rceil$;
\end{enumerate}
and, for all $n\ge1$ and $0<B\le n$ with $\sigma_E(4K_Gn/B)<\infty$,
\begin{enumerate}[label=(\alph*),itemsep=1pt,resume]
\item an exact causal service, with error $0$, purity $P+C<B$, $J=(n-1)Q$ and $Q=\bigl\lceil\log\sigma_E(4K_Gn/B)+\log(4n/B)\bigr\rceil$.
\end{enumerate}
\end{proposition}

\begin{proof}
\emph{The unitary.} For a symbol $w$ let $P_w$ be the permutation matrix of $f([w])$ on $\mathbb C^\Omega$, so that symbols with the same label get the
same matrix, and let $U_f$ be the gadget $U_\lambda$ with $\lambda(w)=P_w$. Each triangle block is
$\frac1{\sqrt2}\bigl(\begin{smallmatrix}I&P_y\\P_x&-P_xP_y\end{smallmatrix}\bigr)=\mathrm{diag}(I,P_x)(H\otimes I)\mathrm{diag}(I,P_y)$, so $U_f$ is unitary
although no relation holds among the $P_w$. Writing $\pi_a$ for the permutation of $\Omega$ in slot $a$, namely $f([w_a])$ for every slot except the
lower right corner of a triangle $(x,y,z)$, where it is $f([x])f([y])$, we have $U_f=\sum_ac_aE_a\otimes P_{\pi_a}$.

\emph{Pure bath points.} Since $\tau_M$ is the average of the vector states of the points of $\Omega$, $\Phi_f=\frac1M\sum_{\omega\in\Omega}\Phi_\omega$ with
$\Phi_\omega(X)=\mathrm{Tr}_\Omega\,U_f(X\otimes|\omega\rangle\langle\omega|)U_f^*$. The isometry $|\psi\rangle\mapsto U_f(|\psi\rangle\otimes|\omega\rangle)$ is
$\sum_ac_aE_a\otimes|\pi_a(\omega)\rangle$, so
\begin{gather*}
\Phi_\omega(X)=\sum_{a,b\in\mathcal T}c_a\overline{c_b}\;\mathbf 1[\pi_a(\omega)=\pi_b(\omega)]\;E_aXE_b^*,\\
\Phi_G(X)=\sum_{a,b\in\mathcal T}c_a\overline{c_b}\;\mathbf 1\bigl[[w_a]=[w_b]\bigr]\;E_aXE_b^*,
\end{gather*}
the second because $\tau(\lambda(g)\lambda(h)^*)=\delta_{gh}$. Call $\omega$ \emph{good} if $f([x])f([y])(\omega)=f([z])(\omega)$ for every triangle
$(x,y,z)$ and $f(g)(\omega)\ne f(h)(\omega)$ for all distinct labels $g,h$. At a good point $\pi_a(\omega)=f([w_a])(\omega)$ for every slot, since
$[w_a]=[z]$ at the lower right corner of a triangle $(x,y,z)$; so $\pi_a(\omega)=\pi_b(\omega)$ exactly when $[w_a]=[w_b]$, and $\Phi_\omega=\Phi_G$.

\emph{Bad points.} The labels of a triangle satisfy $[x][y]=[z]$ with all three in $E$, so $f([x])f([y])$ and $f([z])$ differ on at most $M/r$ points;
and distinct labels $g,h\in E$ have $f(g)(\omega)=f(h)(\omega)$ for at most $M/r$ points. So the fraction $\vartheta$ of bad points is at most $K_G/r$,
and $\Phi_f=(1-\vartheta)\Phi_G+\vartheta\,\Xi$, where $\Xi$ is the average of $\Phi_\omega$ over the bad points (any channel if there are none).

\emph{Exchange.} Let $r=2K_Gn/\epsilon$, let $f$ have degree $M=\sigma_E(r)$, and let $Q$ be as in (a), so that $M\le2^Q\epsilon/(2n)$. Split the
$Q$-qubit memory as $(\mathbb C^\Omega\otimes\mathbb C^{k_0})\oplus\mathbb C^{g_0}$ with $k_0=\lfloor2^Q/M\rfloor$ and $g_0=2^Q-k_0M<M$. In every round
apply $U_f\otimes I_{k_0}$ on the first summand and the identity on the second, and after every round but the last export all $Q$ memory qubits and
import $Q$ fresh qubits in the maximally mixed state; the initial state is maximally mixed as well, so $P=C=0$. Every round then meets a fresh
flat memory, and the observer faces $n$ independent uses of the channel $(1-g_0/2^Q)\Phi_f+(g_0/2^Q)\mathrm{id}$, which is within
$\vartheta+g_0/2^Q\le\epsilon/(2n)+\epsilon/(2n)$ of $\Phi_G$ in diamond distance, by convexity. A hybrid argument over the rounds bounds the error by $\epsilon$.

\emph{Exact service.} Let $0<B\le n$, put $\eta=B/(4n)\le\frac14$ and $r=K_G/\eta$, let $f$ have degree $M=\sigma_E(r)$, and let $Q$ be as in (c), so
that $M\le\eta2^Q$. Let $\Omega_0\subseteq\Omega$ be the set of good points, so that $|\Omega_0|\ge(1-\eta)M$. Split the memory as
$(\mathbb C^\Omega\otimes\mathbb C^{k_0})\oplus\mathbb C^{g_0}$ as in the exchange case, with $k_0=\lfloor2^Q/M\rfloor$ and $g_0<M\le\eta2^Q$, and let
$\beta$ be the maximally mixed state on $\mathrm{span}\{|\omega\rangle:\omega\in\Omega_0\}\otimes\mathbb C^{k_0}$. Since $k_0M=2^Q-g_0$, its deficit is
\[
Q-\log(|\Omega_0|k_0)=\log\frac{2^Q}{2^Q-g_0}+\log\frac M{|\Omega_0|}\le2\log\frac1{1-\eta}\le\frac{8\eta}{3\ln2}<\frac Bn .
\]
Start the memory in $\beta$, apply $U_f\otimes I_{k_0}$ on the first summand in every round, and after every round but the last export all $Q$
memory qubits and import $Q$ fresh qubits in the state $\beta$. The $n$ registers used have total deficit $P+C<B$. In every round the memory is a
fresh copy of $\beta$, a uniform mixture of vector states $|\omega\rangle\otimes|j\rangle$ with $\omega$ good, and $\Phi_\omega=\Phi_G$ at good points;
so every round applies exactly $\Phi_G$, independently of the other rounds, and the error is $0$.

\emph{No exchange.} Let $f$ be as in the exchange case and $Q$ as in (b), so that $M^n\le2^Q\epsilon/2$. Split the memory as
$(\mathbb C^{\Omega^n}\otimes\mathbb C^{k_0})\oplus\mathbb C^{g_0}$ with $k_0=\lfloor2^Q/M^n\rfloor$ and $g_0<M^n$, start it maximally mixed, and in round $s$
apply $U_f$ to the input and the $s$-th factor $\mathbb C^\Omega$ of the first summand, and the identity on the second. With probability
$1-g_0/2^Q\ge1-\epsilon/2$ the memory starts in the first summand, uniformly distributed, and then the $n$ factors are independent flat baths, each
used once: the observer faces $n$ independent uses of $\Phi_f$, which is within $n\vartheta\le\epsilon/2$ of the ideal experiment. The remaining
branch has weight at most $\epsilon/2$, so the error is at most $\epsilon$.
\end{proof}

In particular, if $\sigma_E(r)$ is finite for every $r$, as it is when $G$ is sofic, then $\Phi_G\in\cK_d$, since
$\ddia(\Phi_f,\Phi_G)\le K_G/r$ for models at every accuracy. This is a quantitative form, for sofic groups, of the fact that gadget channels of
hyperlinear groups lie in $\cK_d$~\cite{DouglasMghirbiB}.

\subsection{Streaming at the rank rate}\label{sec:streaming}

The services of Proposition~\ref{prop:sofic-upper} refresh the whole memory in every round. At the rank rate the memory has to be kept, and a sofic
model cannot simply be reused, since its bad points would then accumulate error. The following theorem keeps the model for all rounds and pays
for its bad points in purity instead. It uses the good points of a model as an exact \emph{rectangular} dilation, whose input bath is smaller than
its output bath, and lets the part of the memory in use grow slightly in every round. The mixing and export steps are those of the rolling
reservoir of~\cite[Appendix~A]{DouglasMghirbiB}.

\begin{theorem}[Streaming at the rank rate]\label{thm:streaming}
Let $\Phi$ be a channel on $M_d$ with linearly independent Kraus operators $A_1,\dots,A_r$, $r\ge2$, and let $1\le k\le k'$. Suppose that
$W=\sum_iA_i\otimes B_i$, with $B_i:\mathbb C^k\to\mathbb C^{k'}$, is an isometry from $\mathbb C^d\otimes\mathbb C^k$ to $\mathbb C^d\otimes\mathbb C^{k'}$ and that
\begin{equation}\label{eq:rectangular}
\Tr(B_i^*B_j)=k\,\delta_{ij},\qquad \sum_iB_i^*B_i=r\,I_k .
\end{equation}
Then for all $n\ge2$ and $0<\epsilon<1$ there is a causal service of $\Phi$ for $n$ rounds with error at most $\epsilon$, $C=0$,
$J=\lfloor\frac{n-1}2\log r\rfloor$ and
\[
P\le n\log\frac{k'}k+2,\qquad Q\le\log k+n\log\frac{k'}k+O_r\Bigl(\log\frac n\epsilon+\log(2+\log k)\Bigr).
\]
\end{theorem}

The first identity in~\eqref{eq:rectangular} says that $W$ with a flat input bath realizes $\Phi$ exactly:
$\Phi(X)=\mathrm{Tr}_{k'}[W(X\otimes I_k/k)W^*]=\sum_iA_iXA_i^*$. The second says that every Kraus history costs the same. For a channel with a
flat probe of rank $r$ it follows from $W^*W=I$ alone: apply $\Tr(\rho_*\,\cdot\,)$ to the first factor of $\sum_{i,j}A_i^*A_j\otimes B_i^*B_j=I$,
where $\rho_*$ is the input marginal of the probe, for which $\Tr(\rho_*A_i^*A_j)=\delta_{ij}/r$.

\begin{proof}
We use the frames of~\cite[Appendix~A]{DouglasMghirbiB}. A \emph{frame} is a linear map $F:\mathbb C^M\to\mathbb C^D\otimes\mathbb C^S$; its columns
are the vectors, in the memory $\mathbb C^D$ and an inaccessible exterior $\mathbb C^S$, attached to the Kraus histories. The exterior holds the
purifier of the initial memory, the exported qubits and the purifiers of the imported ones. Put $g(F)=\|F^*F\|$ and
$\alpha(F)=\|\mathrm{Tr}_{\mathbb C^D}FF^*\|$. Two facts are used.
\begin{itemize}[leftmargin=1.2em,itemsep=1pt]
\item \emph{Concentration}~\cite[Lemma~A.1]{DouglasMghirbiB}. For Hermitian $X$ on $\mathbb C^D$ and Haar-random $V\in U(D)$, the matrix
$Z=F^*((V^*XV)\otimes I)F-\tau_D(X)F^*F$ satisfies $\Pr\{\|Z\|\ge s\}\le2M\exp[-Ds^2/(4\alpha(F)g(F)\|X\|^2)]$.
\item \emph{Export}~\cite[Lemma~A.4]{DouglasMghirbiB}. If $\mathbb C^D=\mathbb C^{D/e}\otimes\mathbb C^e$, there is a unitary of $\mathbb C^D$ after
which exporting the factor $\mathbb C^e$ and importing a maximally mixed $e$-dimensional register give a frame $F'$ with $F'^*F'=F^*F$ and, for every
$\eta>0$, $\alpha(F')\le(1+\eta)\alpha(F)/e^2+4e^2g(F)\ln(8Se^2)/(\eta D)$.
\end{itemize}

\emph{Schedule and dimensions.} Put $h=\log r$, $E_p=2^{\lceil h/2\rceil}$, $J_t=\lfloor th/2\rfloor$, $j_t=J_t-J_{t-1}\le\lceil h/2\rceil$ and
$e_t=2^{j_t}$, a divisor of $E_p$. Choose $D_0$, a multiple of $kE_p$, with $D_0\ge nkE_p$ and
\begin{equation}\label{eq:D0}
D_0^2\ \ge\ \frac{32\,A_*\,r^4n^2\ln(8r^{n+2})}{\epsilon^2},\qquad\text{where}\quad A_*=12\bigl(1+8E_p^2n^2L_*\bigr),\quad L_*=\ln(8D_0r^nE_p^2).
\end{equation}
Put $q_0=k'/k$. For $t\ge1$ let $\widetilde D_t=q_0D_{t-1}$, a multiple of $k'E_p$, and let $D_t$ be the least multiple of $kE_p$ with
$D_t\ge\widetilde D_t$. Then $D_t\le q_0D_{t-1}+kE_p$, so $D_n\le q_0^n(D_0+nkE_p)\le2q_0^nD_0$. The memory has $Q=\lceil\log D_n\rceil$ qubits, of which the last $\lceil h/2\rceil$
form a \emph{port}. A space whose dimension $D'$ is a multiple of $E_p$ is realized as a subspace of the first $Q-\lceil h/2\rceil$ qubits of dimension
$D'/E_p$, tensored with the port; the \emph{active space} of round $t$ is such a space of dimension $D_{t-1}$, identified with
$\mathbb C^k\otimes\mathbb C^{D_{t-1}/k}$.

\emph{The device.} The memory starts maximally mixed on the active space of round $1$, of dimension $D_0$. In round $t$ the device applies a unitary
$V_t$ of the active space and then $W\otimes I_{D_{t-1}/k}$, which maps the input and the active space isometrically onto the output and a space of
dimension $\widetilde D_t$; it releases the output, then applies a unitary $V'_t$ of that space (it acts on the memory only, so it could equally act before the release), and, if $t<n$, exports $j_t$ port qubits and imports $j_t$
qubits in the maximally mixed state in their place. The resulting space of dimension $\widetilde D_t$ lies in the active space of round $t+1$. These
maps are completed arbitrarily to unitaries of the input and the whole memory; the completion never acts, since the memory never leaves the
active spaces. So $C=0$, $J=J_{n-1}$, and $P=Q-\log D_0\le n\log q_0+2$.

\emph{Frames.} Purify the initial memory: its frame $F_0$ is one unit vector, with $\alpha(F_0)=1/D_0$. Let $F$ be the frame before round $t$, with one
column $F_y$ for each history $y\in[r]^{t-1}$. After the call the frame has the columns $(B_i\otimes I)(V_t\otimes I)F_y$, indexed by the histories
$(y,i)$. Its Gram matrix has the blocks $F^*\bigl((V_t^*(B_i^*B_j\otimes I)V_t)\otimes I\bigr)F$, whose Haar means are
$\tau(B_i^*B_j\otimes I)F^*F=\delta_{ij}F^*F$ by the first identity in~\eqref{eq:rectangular}. By the second identity and cyclicity of the partial trace
over the memory, the exterior marginal after the call is exactly $r$ times the one before, whatever $V_t$ is; so the call multiplies $\alpha$ by $r$.

\emph{Choice of the unitaries.} We choose $V_1,V'_1,V_2,\dots$ in turn so that the frame $F_t$ after round $t$ satisfies
\begin{equation}\label{eq:streaming-induction}
\|F_t^*F_t-I\|\le t\epsilon/n,\qquad \alpha(F_t)\le A_*/D_0 .
\end{equation}
This holds for $t=0$. Suppose it holds for $t-1$; then $g(F_{t-1})\le2$. Since $B_i^*B_i\le rI$, the Hermitian and anti-Hermitian parts of the $B_i^*B_j$
have norm at most $r$. Apply the concentration lemma to these $2r^2$ operators, tensored with $I_{D_{t-1}/k}$, with $D=D_{t-1}\ge D_0$, at most $r^n$
columns and level $s=\epsilon/(2rn)$: by~\eqref{eq:D0} the failure probabilities add up to at most
$4r^{n+2}\exp[-D_0^2\epsilon^2/(32A_*r^4n^2)]\le\frac12$. So some $V_t$ makes every block of the new Gram matrix deviate from its mean by at most
$\epsilon/(rn)$, hence the $r\times r$ block matrix by at most $\epsilon/n$, and the Gram matrix after the call is within $t\epsilon/n$ of $I$; exports and
imports preserve it. Choose $V'_t$ by the export lemma with $e=e_t$, $\eta=1/n$ and $D=\widetilde D_t\ge D_0$. The exterior has dimension
$S\le D_0\,2^{2J_{t-1}}\le D_0r^n$, so $\ln(8Se_t^2)\le L_*$, and
\[
\alpha(F_t)\ \le\ \Bigl(1+\frac1n\Bigr)\frac{r\,\alpha(F_{t-1})}{e_t^2}+\frac{8E_p^2nL_*}{D_0}.
\]
Put $c_t=2^{2J_t}r^{-t}$, so that $c_0=1$, $\frac14<c_t\le1$ and $c_tr/e_t^2=c_{t-1}$. Multiplying by $c_t$ and unrolling gives
$c_t\alpha(F_t)\le(1+\frac1n)^t(1+8E_p^2n^2L_*)/D_0$, hence $\alpha(F_t)\le4e(1+8E_p^2n^2L_*)/D_0\le A_*/D_0$. In round $n$ there is no export, and only the
Gram bound is needed.

\emph{Error.} Purify the observer. In the ideal experiment, with the Stinespring isometry $\xi\mapsto\sum_iA_i\xi\otimes|i\rangle$ in each round, its
final vector is $\sum_y\psi_y\otimes|y\rangle$; in the service it is $\sum_y\psi_y\otimes F_n|y\rangle$ with the same vectors $\psi_y$, because the device
acts on the input and the memory through $W=\sum_iA_i\otimes B_i$ and unitaries of the memory, and the observer never touches the memory or the
exterior. The polar part $T=F_n(F_n^*F_n)^{-1/2}$ is an isometry with $\|F_n-T\|=\|(F_n^*F_n)^{1/2}-I\|\le\epsilon$. So the actual final vector is within
$\epsilon$ of the ideal one mapped by $I\otimes T$, and after tracing out the memory and the exterior the observer's final states are within
$\epsilon$ in trace distance.

\emph{Memory.} Finally $Q\le\log D_n+1\le\log D_0+n\log q_0+2$. The conditions on $D_0$ hold for $D_0=kE_p\lceil K_r(n/\epsilon)^4(1+\ln k)\rceil$ with $K_r$
large enough in terms of $r$, and then $\log D_0=\log k+O_r(\log(n/\epsilon)+\log(2+\log k))$.
\end{proof}

\begin{corollary}[The lamp channel at the rank rate]\label{cor:streaming-lamps}
There is $K$ such that for every $n\ge2$ the channel $\Phi$ of Theorem~\ref{thm:main-channel} has a causal service for $n$ rounds with error at most
$1/16$, $C=0$, $J=\lfloor\frac{n-1}2\log r_*\rfloor$ and $P,Q\le K\sqrt{n\log n}$.
\end{corollary}

\begin{proof}
Let $E$ be the chunk of labels of Proposition~\ref{prop:sofic-upper}, let $s\ge2K_G$, and let
$f:E\to\Sym(\Omega)$ be a $(1-s^{-1})$-expansive $s^{-1}$-morphism of degree $M=\sigma_E(s)$. Let $\Omega_0$ be its set of good points, so that
$k=|\Omega_0|\ge(1-K_G/s)M$, and for each label $g$ let $B_g:\mathbb C^{\Omega_0}\to\mathbb C^\Omega$ send $|\omega\rangle$ to $|f(g)\omega\rangle$. At a good point
$\omega$ every slot of $U_f$ acts through the permutation of its label, so $U_f(\xi\otimes|\omega\rangle)=\sum_gA_g\xi\otimes|f(g)\omega\rangle$, and
$W=\sum_gA_g\otimes B_g$ is the restriction of the unitary $U_f$ to $\mathbb C^d\otimes\mathbb C^{\Omega_0}$, an isometry. Each $B_g^*B_g$ is $I_k$, and
$\Tr(B_g^*B_h)=0$ for $g\ne h$ because distinct labels do not collide at good points; the $A_g$ are linearly independent, having disjoint
supports. So Theorem~\ref{thm:streaming} applies with $r=r_*$ and $k'=M$. Here $\log(M/k)\le-\log(1-K_G/s)\le2K_G/(s\ln2)$, and
$\log k\le\log\sigma_E(s)\le C_Es\log s$ by Theorem~\ref{thm:main-group}(b). With $\epsilon=1/16$ the theorem gives
$Q\le C_Es\log s+2K_Gn/(s\ln2)+O(\log n)$ and $P\le2K_Gn/(s\ln2)+2$; take $s=\max(2K_G,\lceil\sqrt{n/\log n}\,\rceil)$.
\end{proof}

\subsection{Proof of Theorem~\ref{thm:main-channel}}\label{sec:main-channel-proof}

\emph{Closure and services.} By Theorem~\ref{thm:main-group}(b) the chunk $E$ of $\Gamma$ that consists of $1$ and the labels of the gadget of
$R_\Phi$ has $\log\sigma_E(r)\le C_Er\log r$. Proposition~\ref{prop:sofic-upper} therefore gives finite tracial channels within $K_G/r$ of $\Phi$ for
every $r$, so $\Phi\in\cK_{873}$; this also follows from~\cite{DouglasMghirbiB}, since $\Gamma$ is amenable (Lemma~\ref{lem:structure}) and hence
hyperlinear. Parts (c) and (a) of the proposition give the exact services and the zero-purity services in (a), since
$\log\sigma_E(R)\le C_ER\log R$ for $R\ge2$ and $\log(4K_Gn/B)\le(1+\log2K_G)\log(2n/B)$, because $\log(2n/B)\ge1$; for the zero-purity services use
$\epsilon=1/16$. Corollary~\ref{cor:streaming-lamps} gives the services in (b). The bounds on $r_*$ were proved in Section~\ref{sec:lamp-channel}.

\emph{Part (a).} Let a service have error $\epsilon\le1/16$ and $P+C\le B$, where $\log n\le B$. Theorem~\ref{thm:extraction} gives a round with
bath dimension $2^Q$, $\ddia(\Psi_t,\Phi)\le\epsilon/(1-\epsilon)\le1/15$ and $\nu_t=\ell_t+\Delta S_t\le t_n$. Since $\epsilon+h_2(\epsilon)<1$,
$\epsilon\log n\le\log n$ and $1\le\log n\le B$,
\[
t_n\le\frac{16}{15}\cdot\frac{B+1+\log n}{n}\le\frac{16B}{5n}\qquad(n\ge2).
\]
Let $C_R\ge1$ be the constant of Corollary~\ref{cor:amplification}(a) for $R_\Phi$ and the transport of Section~\ref{sec:lamp-channel}, and put
$y_n=5n/(16C_RB)$, so that $C_R\nu_ty_n\le1$ and $y_n\le n$. Put $c_0=5/(16C_Ry_0)$. If $B\le c_0n$, then $y_n\ge y_0$, and the corollary gives
\[
Q\ \ge\ S(\rho_t)\ \ge\ \frac{y_n}{32(\log y_n)^{2\kappa}}\ \ge\ \frac{5}{512\,C_R}\cdot\frac{n}{B(\log n)^{2\kappa}} .
\]
For the second statement take $B=\max(1,A_0)\log n$, which lies between $\log n$ and $c_0n$ for large $n$.

\emph{Part (b).} Proposition~\ref{prop:purity-budget}, with the flat probe of rank $r_*$, gives $P+C\le B:=3Q+6$. Theorem~\ref{thm:extraction} now
gives a round with $\nu_t\le t_n\le\frac{16}{15}(3Q+7+\log n)/n$. Put $y_n=15n/(16C_R(3Q+7+\log n))\le n$. If $y_n<y_0$, then
$3Q+7+\log n>15n/(16C_Ry_0)$, so $Q\ge c'n$ for large $n$. Otherwise Corollary~\ref{cor:amplification}(a) gives $Q\ge y_n/(32(\log n)^{2\kappa})$,
that is,
\[
Q\,(3Q+7+\log n)\ \ge\ \frac{15}{512\,C_R}\cdot\frac{n}{(\log n)^{2\kappa}} .
\]
If $Q\le\log n+7$ the left side is at most $4(\log n+7)^2$, which is smaller than the right side for large $n$. So $Q>\log n+7$, the left side is
at most $4Q^2$, and $Q\ge c\sqrt n/(\log n)^\kappa$. \qed

\begin{remark}[Memory between two readings of the profile]\label{rem:sandwich}
For every gadget channel of a sofic group that satisfies the hypotheses of Theorem~\ref{thm:weighted}, the least memory of services with error
$1/16$ and purity $O(\log n)$ lies between two readings of the group at scale $n$: below, the entropy forced by as many cheaply commuting copies of
$S_3$ as Corollary~\ref{cor:amplification} fits at relator defect about $\sqrt{\log n/n}$; above, $\log\sigma_E(Cn)+O(\log n)$ by
Proposition~\ref{prop:sofic-upper}, already at zero purity and with exchange $(n-1)Q$. For the configuration lamps both readings are
$n^{1+o(1)}$. The upper reading is the sofic profile, not the F{\o}lner function: the latter is doubly exponential for $\Gamma$
(Proposition~\ref{prop:folner-lower}), so models built from F{\o}lner sets~\cite[Proposition~3.14]{Cornulier2013} would only give memory
exponential in $n$, while the hash models
behind Theorem~\ref{thm:main-group}(b) give $O(n\log n)$. For the grid host of Remark~\ref{rem:grid-channel}, whose proofs are only sketched, the two readings would be $n/(\log n)^{15.75}$ and
$O(n\log n)$.
\end{remark}

\begin{remark}[The rank rate is the diagonal of the trade-off]\label{rem:sqrt-limit}
At the rank rate a device with error at most $1/16$ can spend at most $3Q+6$ bits of purity, so Theorem~\ref{thm:main-channel}(a) with $B$ of order $Q$ forces $Q\gtrsim\sqrt n$,
and the streaming services attain this with purity of the same order. In general, for every channel with a flat probe, padding a rank-rate
service with error at most $1/16$ with $2Q+6$ maximally mixed qubits turns it into a service whose purity is at most its memory; so the least memory at purity at most
memory is at most three times the least memory at the rank rate, plus six. For $\Phi$ both are $n^{1/2+o(1)}$. The one-round method cannot give
more than $\sqrt n$ at the rank rate, since condition~\eqref{eq:weighted-condition} certifies at most about $1/\nu$ copies; Theorem~\ref{thm:streaming}
shows that for $\Phi$ nothing more is true.
\end{remark}

%======================================================================
\section{Proof of the weighted one-round theorem}\label{sec:weighted-proof}
%======================================================================

The proof has four steps. First, a round with small leakage and entropy production yields unitaries on the bath under which every relation
holds up to $O(\sqrt\nu)$, with errors weighted by the memory state (Lemma~\ref{lem:weighted-extraction}). Second, weighted errors add up over van
Kampen diagrams, at a price for each face and each vertex that is paid once (Lemma~\ref{lem:charged-area}). Third, each transported copy of $S_3$ is rounded to an
exact representation on the same space, so the state and its entropy do not change (Lemma~\ref{lem:rounding}). Fourth, measuring the rounded
three-cycles one after another on one half of a purification of the memory state produces about one bit per copy that the other half also knows
(Theorem~\ref{thm:sector}).

The section generalises the argument of~\cite[Appendix~B]{DouglasMghirbiB}, which proves the one-round theorem for the Houghton
gadget~\cite[Theorem~5.5]{DouglasMghirbiB}, to every gadget channel that satisfies the hypotheses of Theorem~\ref{thm:weighted}. The weighted norms
are those of~\cite[Appendix~B]{DouglasMghirbiB}; Lemmas~\ref{lem:weighted-extraction} and~\ref{lem:charged-area} extend Lemmas~B.4 and~B.2
there from the Houghton gadget to every such gadget; Lemma~\ref{lem:rounding} is Lemma~B.3 there; and Theorem~\ref{thm:sector} is a form of
Theorem~B.1 there, proved by sequential measurement.

\subsection{Weighted norms}

Throughout, $\|\cdot\|_F$ is the unnormalized Hilbert--Schmidt norm on $M_M$, $\rho$ is a state on $\mathbb C^M$ and $\xi=\sqrt\rho$, so
$\|\xi\|_F=1$. For operators $W$ on $\mathbb C^M$ put
\[
E_\rho(W)=\|(W-I)\xi\|_F,\qquad\chi_\rho(W)=\|W\xi-\xi W\|_F,
\]
the \emph{displacement} and the \emph{centrality defect} of $W$ at $\rho$. We drop the subscript $\rho$. For unitaries $u,v$ and any operator $F$:
\begin{enumerate}[label=(W\arabic*),itemsep=1pt]
\item $E(uv)\le E(u)+E(v)$ and $E(u^*)=E(u)$;
\item $\chi(uv)\le\chi(u)+\chi(v)$ and $\chi(u^*)=\chi(u)$;
\item $\|Fv\xi\|_F\le\|F\xi\|_F+\|F\|\,\chi(v)$;
\item $|E(uWu^*)-E(W)|\le2\chi(u)$ for every $W$ with $\|W\|\le1$;
\item $\|\xi F\|_F=\|F^*\xi\|_F$, and $\|\xi F\|_F=\|F\xi\|_F$ if $F$ is normal.
\end{enumerate}
Indeed, $(uv-I)\xi=u(v-I)\xi+(u-I)\xi$ and $(u^*-I)\xi=-u^*(u-I)\xi$; $uv\xi-\xi uv=u(v\xi-\xi v)+(u\xi-\xi u)v$ and
$u^*\xi-\xi u^*=-u^*(u\xi-\xi u)u^*$; $Fv\xi=F\xi v+F(v\xi-\xi v)$; $(uWu^*-I)\xi=u(W-I)u^*\xi$, to which (W3) applies with $\|W-I\|\le2$, in both
directions; and $\|F^*\xi\|_F^2=\Tr(\xi FF^*\xi)$, which equals $\Tr(\xi F^*F\xi)$ when $F$ is normal. A word $w$ in $S^{\pm1}$ is evaluated by an
assignment $\phi$ of unitaries to $S$ through $\phi(x^{-1})=\phi(x)^*$; by (W2), $\chi(\phi(w))\le|w|\max_{x\in S}\chi(\phi(x))$.

\subsection{Extraction}

\begin{lemma}[Weighted extraction]\label{lem:weighted-extraction}
Let $G$, $S$ and $R$ be as in Section~\ref{sec:channel-setting}, with $a,b,j\in S$, $j^{-1}ab\in R$ and $j\ne1$ in $G$, and let $m$ be as in
Theorem~\ref{thm:weighted}. Let $\Psi(X)=\mathrm{Tr}_M\,U(X\otimes\rho)U^*$
be a round with $u=\ddia(\Psi,\Phi_G)\le1/8$, leakage $\ell$ relative to $\Phi_G$ and entropy production $\Delta S$, and put $\nu=\ell+\Delta S$,
$\eta=\sqrt{d\nu}$ and $e=256\,m\,\eta$. There are unitaries $\phi(x)$, $x\in S$, on $\mathbb C^M$ such that
\begin{enumerate}[label=(\roman*),itemsep=1pt]
\item $E(\phi(r))\le e$ for every $r\in R$;
\item $\chi(\phi(x))\le e$ for every $x\in S$;
\item $E(\phi(a)\phi(b))\ge1-e$.
\end{enumerate}
\end{lemma}

\begin{proof}
\emph{Normalization.} Let $(\zeta_i)$ and $(\zeta'_i)$ be orthonormal eigenbases of $\rho$ and $T(\rho)$ with decreasing eigenvalues $\lambda_i$ and $\mu_i$.
Replacing $U$ by $(I\otimes W)U$, where the unitary $W$ maps $\zeta'_i$ to $\zeta_i$, changes neither $\Psi$ nor its leakage and conjugates $T(\rho)$
by $W$. So we may assume $T(\rho)=\sum_i\mu_i|\zeta_i\rangle\langle\zeta_i|$. Input and output bath are now the same space, and
$\xi=\sum_i\sqrt{\lambda_i}|\zeta_i\rangle\langle\zeta_i|$ acts on both; this does not assert $T(\rho)=\rho$.

\emph{Entropy production controls commutators.} Write $U_{ji}=(I\otimes\langle\zeta_j|)U(I\otimes|\zeta_i\rangle)\in M_d$ and $t_{ji}=\frac1d\|U_{ji}\|_F^2$.
Since the columns and the rows of $U$ are orthonormal, $t$ is doubly stochastic, and the diagonal of $T(\rho)$ gives $\mu_j=\sum_it_{ji}\lambda_i$. The
distributions $p_{ji}=t_{ji}\lambda_i$ and $q_{ji}=t_{ji}\mu_j$ satisfy
\[
D(p\|q)=\sum_{j,i}t_{ji}\lambda_i\log\frac{\lambda_i}{\mu_j}=-S(\rho)+S(T(\rho))=\Delta S .
\]
By Jensen's inequality $-2\ln\sum\sqrt{p_{ji}q_{ji}}\le D(p\|q)\ln2$, and $-2\ln x\ge2(1-x)$, so
$\sum_{j,i}t_{ji}(\sqrt{\lambda_i}-\sqrt{\mu_j})^2\le\Delta S\ln2$. The matrix $t$ is a convex combination of permutation matrices (Birkhoff), and both sequences
decrease, so the rearrangement inequality gives $\sum_{j,i}t_{ji}\sqrt{\lambda_i\mu_j}\le\sum_j\sqrt{\lambda_j\mu_j}$; expanding the squares and using
double stochasticity, $\sum_j(\sqrt{\lambda_j}-\sqrt{\mu_j})^2\le\sum_{j,i}t_{ji}(\sqrt{\lambda_i}-\sqrt{\mu_j})^2$. The bath block $(j,i)$ of $U(I\otimes\xi)-(I\otimes\xi)U$
is $(\sqrt{\lambda_i}-\sqrt{\lambda_j})U_{ji}$, and $(\sqrt{\lambda_i}-\sqrt{\lambda_j})^2\le2(\sqrt{\lambda_i}-\sqrt{\mu_j})^2+2(\sqrt{\mu_j}-\sqrt{\lambda_j})^2$.
Therefore, with $U=\sum_{o,o'}|o\rangle\langle o'|\otimes A_{oo'}$ in system blocks $A_{oo'}\in M_M$,
\begin{equation}\label{eq:block-commutators}
\sum_{o,o'}\|[A_{oo'},\xi]\|_F^2=\|[U,I\otimes\xi]\|_F^2\le4d\,\Delta S\ln2=:\epsilon_S^2 .
\end{equation}

\emph{Leakage controls residuals.} The Choi state $\mathcal J(\Psi)$ is the marginal on $\mathbb C^d\otimes\mathbb C^d$ of the unit vector
$d^{-1/2}\sum_{o,o'}|oo'\rangle\otimes A_{oo'}\xi$ in $\mathbb C^d\otimes\mathbb C^d\otimes\mathcal H$, where $\mathcal H$ is $M_M$ with the Hilbert--Schmidt
inner product: indeed $\langle A_{o_2o_2'}\xi,A_{o_1o_1'}\xi\rangle=\Tr(A_{o_1o_1'}\rho A_{o_2o_2'}^*)=\langle o_1|\Psi(|o_1'\rangle\langle o_2'|)|o_2\rangle$. The support
of $\mathcal J(\Phi_G)$ is spanned by the orthogonal vectors $|A_g\rangle\!\rangle$, of squared norms $\mu_g=\|A_g\|_F^2$. Projecting onto it, put
\[
X_g=\mu_g^{-1}\sum_{[w_a]=g}\overline{c_a}A_a,\qquad R_{oo'}=A_{oo'}-\sum_g(A_g)_{oo'}X_g,
\]
where $A_a$ is the block at slot $a$. Then $\sum_{o,o'}\|R_{oo'}\xi\|_F^2=d\ell=:\epsilon_\ell^2$. At a slot $A_a=c_aX_{[w_a]}+R_a$, and at every other position
$A_{oo'}=R_{oo'}$. We have $\epsilon_\ell\le\eta$, $\epsilon_S\le2\sqrt{\ln2}\,\eta$, and, by Cauchy--Schwarz, $\epsilon_\ell+\epsilon_S\le\sqrt{1+4\ln2}\,\eta<2\eta$ and
$2\epsilon_\ell^2+2\epsilon_S^2\le8\ln2\,\eta^2$.

\emph{Anchors.} Every symbol $w$, and the empty word, occupies a diagonal position $z=z(w)$ whose row and column contain no other slot; there the
coefficient is $1$, so $B_w:=A_{zz}=X_{[w]}+R_{zz}$. Column unitarity gives $\Tr\rho(I-B_w^*B_w)=\sum_{o\ne z}\|R_{oz}\xi\|_F^2\le\epsilon_\ell^2$. Row
unitarity gives $\Tr\rho(I-B_wB_w^*)=\sum_{o'\ne z}\|\xi A_{zo'}\|_F^2\le\sum_{o'\ne z}(\|R_{zo'}\xi\|_F+\|[A_{zo'},\xi]\|_F)^2\le2\epsilon_\ell^2+2\epsilon_S^2$.
Let $B_w=V_w|B_w|=|B_w^*|V_w$ with $V_w$ unitary. Since $(1-x)^2\le1-x^2$ on $[0,1]$,
\[
\|(V_w-B_w)\xi\|_F\le\epsilon_\ell,\qquad\|\xi(V_w-B_w)\|_F=\|\xi(I-|B_w^*|)\|_F\le\sqrt{8\ln2}\,\eta,
\]
and so $\chi(V_w)\le\|(V_w-B_w)\xi\|_F+\|[B_w,\xi]\|_F+\|\xi(B_w-V_w)\|_F\le\epsilon_\ell+\epsilon_S+\sqrt{8\ln2}\,\eta<4.3\eta$, since
$\epsilon_\ell+\epsilon_S\le\sqrt{1+4\ln2}\,\eta$.

\emph{Triangle slots.} In the block of a triangle $(x,y,z)$, compare each slot $a$ with the anchor of the symbol $w'(a)$ that carries the same
label: the empty word at the upper left, $y$ at the upper right, $x$ at the lower left and $z$ at the lower right. Put $F_a=A_a-c_aV_{w'(a)}$. Then
$F_a=R_a-c_aR_{z'z'}+c_a(B_{w'(a)}-V_{w'(a)})$ with $z'=z(w'(a))$, and since $a$ and $(z',z')$ are different positions, Cauchy--Schwarz gives
\begin{gather*}
\|F_a\xi\|_F\le\sqrt{3/2}\,\epsilon_\ell+2^{-1/2}\epsilon_\ell<2\eta,\qquad\|F_a\|<2,\\
\|\xi F_a\|_F\le\|F_a\xi\|_F+\|[A_a,\xi]\|_F+2^{-1/2}\chi(V_{w'(a)})<7\eta .
\end{gather*}
For two slots $a_1,a_2$ in the same row of the block, with $A_{a_i}=c_iV_i+F_i$,
\[
A_{a_1}^*A_{a_2}-\overline{c_1}c_2V_1^*V_2=A_{a_1}^*F_2+c_2F_1^*V_2,
\]
whose right-weighted norm is at most $\|F_2\xi\|_F+2^{-1/2}(\|\xi F_1\|_F+\|F_1\|\chi(V_2))<2\eta+2^{-1/2}\cdot17\eta<15\eta$ by (W3) and (W5). The two
columns of the block are orthogonal columns of $U$, so the products for the two rows of the block add up to $-C_1^*C_2$, where $C_1,C_2$ are the
parts of the two columns in the other rows. These are residuals, so $\|C_1^*C_2\xi\|_F\le\|C_2\xi\|_F\le\epsilon_\ell$. The coefficients are
$\overline{c_1}c_2=\frac12$ for the pair (upper left, upper right) and $-\frac12$ for the pair (lower left, lower right), so
\[
\tfrac12\|(V_\varnothing^*V_y-V_x^*V_z)\xi\|_F<15\eta+15\eta+\eta .
\]
Put $W_w=V_\varnothing^*V_w$ for every symbol $w$, and $W_\varnothing=I$. Since $W_xW_y-W_z=V_\varnothing^*V_x(V_\varnothing^*V_y-V_x^*V_z)$,
\begin{equation}\label{eq:triangle}
\|(W_xW_y-W_z)\xi\|_F\le64\eta\quad\text{for every triangle }(x,y,z),\qquad\chi(W_w)<10\eta .
\end{equation}

\emph{Relations.} Let $\phi(x)=W_x$ for $x\in S$. For a letter $x\in S$, the triangle $(x,x^{-1},\varnothing)$ gives
$\|(W_{x^{-1}}-\phi(x^{-1}))\xi\|_F=\|W_x^*(W_xW_{x^{-1}}-I)\xi\|_F\le64\eta$. Fix a word $r=\sigma_1\cdots\sigma_\ell$ of $R$ and let
$D_l=\|(\phi(\sigma_1\cdots\sigma_l)-W_{p_l})\xi\|_F$, with $p_l$ its prefix symbols, $p_1=\sigma_1$ and $W_{p_\ell}=I$. Then $D_1\le64\eta$, and
\[
\phi(\sigma_1\cdots\sigma_l)-W_{p_l}=\bigl(\phi(\sigma_1\cdots\sigma_{l-1})-W_{p_{l-1}}\bigr)\phi(\sigma_l)+W_{p_{l-1}}\bigl(\phi(\sigma_l)-W_{\sigma_l}\bigr)+\bigl(W_{p_{l-1}}W_{\sigma_l}-W_{p_l}\bigr).
\]
By (W3), \eqref{eq:triangle} and the triangle $(p_{l-1},\sigma_l,p_l)$, $D_l\le D_{l-1}+2\cdot10\eta+64\eta+64\eta$. So
$E(\phi(r))=D_\ell\le148\,\ell\,\eta\le e$, which is (i); and (ii) holds because $\chi(\phi(x))<10\eta\le e$.

\emph{Displacement.} Let $z=z(j)$ and let $0$ be the position of the empty word. The labels there are $[j]\ne1$ and $1$, so
$\langle z|\Phi_G(|z\rangle\langle0|)|0\rangle=\tau(\lambda([j]))=0$, while $\langle z|\Psi(|z\rangle\langle0|)|0\rangle=\Tr(B_j\rho B_\varnothing^*)$. Feed
$(|zz\rangle+|00\rangle)/\sqrt2$ into $\Psi\otimes\mathrm{id}$ and $\Phi_G\otimes\mathrm{id}$. The difference $\Delta$ of the outputs is Hermitian and traceless with
$\|\Delta\|_1\le2u$, so each of its entries has modulus at most $\|\Delta\|\le\|\Delta\|_1/2\le u$; the entry at $(zz,00)$ is half the difference
above, so $|\Tr(\rho B_\varnothing^*B_j)|\le2u$. By the anchor bounds, $|\Tr(\rho V_\varnothing^*V_j)-\Tr(\rho B_\varnothing^*B_j)|\le2\epsilon_\ell$, and so
\[
E(\phi(j))^2=\|(V_j-V_\varnothing)\xi\|_F^2=2-2\,\mathrm{Re}\,\Tr(\rho V_\varnothing^*V_j)\ge2-4u-4\epsilon_\ell\ge\tfrac32-4\eta .
\]
By (i) for the word $j^{-1}ab$, $\|(\phi(a)\phi(b)-\phi(j))\xi\|_F=E(\phi(j)^*\phi(a)\phi(b))\le e$, so $E(\phi(a)\phi(b))\ge\sqrt{3/2-4\eta}-e\ge1-e$ if
$\eta\le1/8$. If $\eta>1/8$ then $e>1$ and (iii) is trivial.
\end{proof}

\subsection{Charged area}

Relator errors add up over a van Kampen diagram, but in the weighted norms every conjugation costs twice the centrality defect of the
conjugator. Writing a null word naively as a product of conjugates of relations would charge each conjugating word once for every relation it
encloses. Peeling the diagram one boundary face at a time charges instead each face a bounded amount and each vertex once.

\begin{lemma}[Charged area]\label{lem:charged-area}
Let $S$ and $R$ be finite, with every word of $R$ of length at most $m$, and let $\phi$ assign unitaries to $S$ with $E(\phi(r))\le e$ for $r\in R$
and $\chi(\phi(x))\le e$ for $x\in S$. Every word $w$ of finite area satisfies
\[
E(\phi(w))\ \le\ \bigl(3m\Area_R(w)+|w|\bigr)\,e\ \le\ \bigl((4m+1)\Area_R(w)+2|w|\bigr)\,e .
\]
\end{lemma}

\begin{proof}
Since $\phi$ is a homomorphism on the free group, $E(\phi(w))$ depends only on the free reduction of $w$, which is not longer and has the same
area; so let $w$ be freely reduced. By van Kampen's lemma~\cite[Chapter~V]{LyndonSchupp1977} there is a van Kampen diagram for $w$ with at most
$\Area_R(w)$ faces: a finite, connected and simply connected planar $2$-complex with edges labelled by letters of $S$, whose boundary cycle, read
from a base vertex, is $w$, and in which the boundary cycle of every face, read from a suitable vertex in a suitable direction, is a word of $R$.
We show, by induction on the number of edges, that every such diagram $D$ with $V$ vertices and $A$ faces satisfies
\begin{equation}\label{eq:charge}
E(\phi(w_D))\le\bigl((2m-1)A+2(V-1)\bigr)e ,
\end{equation}
where $w_D$ is its boundary word. If $D$ is a single vertex, $w_D$ is empty. Otherwise let $\varepsilon$ be the first edge of the boundary cycle,
from the base vertex $v_0$ to $v_1$, with label $x^{\pm1}$.

If $\varepsilon$ lies on no face, the boundary cycle traverses it twice, and removing it leaves diagrams $D_0\ni v_0$ and $D_1\ni v_1$ with
$w_D=x^{\pm1}w_{D_1}x^{\mp1}w_{D_0}$, the words read from $v_1$ and $v_0$. By (W1) and (W4), $E(\phi(w_D))\le2e+E(\phi(w_{D_1}))+E(\phi(w_{D_0}))$, and
the vertices and faces of $D_0$ and $D_1$ add up to those of $D$.

Otherwise $\varepsilon$ lies on the boundary of a face $f$. Let $c$ be the boundary cycle of $f$ read from $v_0$ and beginning with $\varepsilon$,
so that $c=\varepsilon\pi$ with $\pi$ a path in $\partial f$ from $v_1$ to $v_0$. Removing $f$ and $\varepsilon$ merges $f$ with the outer
region and leaves a van Kampen diagram $D'$ with the same vertices and one face fewer, whose boundary word read from $v_0$ is
$w_{D'}=\pi^{-1}\varrho$, where $w_D=\varepsilon\varrho$. Hence
$w_D=c\,w_{D'}$ freely. The word $c$ is a cyclic permutation of some $r^{\pm1}$, $r\in R$, that is $c=y^{-1}r^{\pm1}y$ freely with $|y|\le m-1$, so
$E(\phi(c))\le e+2(m-1)e$ by (W4) and (W2), and $E(\phi(w_D))\le(2m-1)e+E(\phi(w_{D'}))$ by (W1).

In both cases~\eqref{eq:charge} follows from the induction hypothesis. Finally $V-1$ is at most the number of edges of $D$, since $D$ is
connected, and counting the two sides of every edge gives twice the number of edges as at most $mA+|w|$, the sum of the lengths of the face
boundaries and of the boundary cycle. So~\eqref{eq:charge} gives $E(\phi(w))\le((2m-1)A+mA+|w|)e$.
\end{proof}

\subsection{Exact \texorpdfstring{$S_3$}{S3} rounding in place}

An approximate copy of $S_3$ is rounded to an exact one by functional calculus on the same space. The state is untouched, so no entropy is
gained or lost.

\begin{lemma}[{Exact rounding~\cite[Lemma~B.3]{DouglasMghirbiB}}]\label{lem:rounding}
Let $X$ and $Y$ be unitaries on $\mathbb C^M$, and suppose that the displacements of $X^2$, $Y^2$ and $(XY)^3$ and the centrality defects of
$X$ and $Y$ are at most $\delta$. There are unitaries $x,c$ on $\mathbb C^M$ with $x^2=c^3=I$ and $xcx=c^{-1}$, such that
\[
\|(x-X)\xi\|_F\le\delta,\qquad\chi(x)\le3\delta,\qquad\|(c-XY)\xi\|_F\le29\delta,\qquad\chi(c)\le56\delta .
\]
\end{lemma}

\begin{proof}
Let $\mathrm{sgn}\,t=1$ for $t\ge0$ and $-1$ for $t<0$, and let $x=\mathrm{sgn}(\mathrm{Re}\,X)$ and $y=\mathrm{sgn}(\mathrm{Re}\,Y)$ by functional calculus;
they are self-adjoint unitaries. On the unit circle $|z-\mathrm{sgn}\,\mathrm{Re}\,z|\le|z^2-1|$: for $z=e^{i\vartheta}$ with $|\vartheta|\le\pi/2$ this reads
$2|\sin(\vartheta/2)|\le2|\sin\vartheta|$, which holds as $\cos(\vartheta/2)\ge\frac12$, and the case $\mathrm{Re}\,z<0$ follows by symmetry. Hence
$|x-X|^2\le|X^2-I|^2$ as functions of the normal operator $X$, and $\|(x-X)\xi\|_F\le E(X^2)\le\delta$; by (W5) also $\|\xi(x-X)\|_F\le\delta$, so
$\chi(x)\le\chi(X)+2\delta\le3\delta$. The same holds for $y$.

Put $Z=XY$ and $w=xy$, so that $xwx=yx=w^*$. Then $\chi(Z)\le2\delta$, $\chi(w)\le6\delta$, and, by (W3),
$\|(w-Z)\xi\|_F\le\|x(y-Y)\xi\|_F+\|(x-X)Y\xi\|_F\le\delta+3\delta=4\delta$. From $w^3-Z^3=w^2(w-Z)+w(w-Z)Z+(w-Z)Z^2$ and (W3),
\[
E(w^3)\le E(Z^3)+4\delta+(4\delta+2\cdot2\delta)+(4\delta+2\cdot4\delta)\le25\delta .
\]
Let $f$ send each point of the unit circle to a nearest cube root of unity, with the ties $e^{\pm i\pi/3}$ and $-1$ sent to $1$. Then
$f(\bar z)=\overline{f(z)}$, and $|z-f(z)|\le|z^3-1|$: for $z=-1$ both sides equal $2$, and otherwise $z=f(z)e^{i\vartheta}$ with
$|\vartheta|\le\pi/3$, and the inequality reads $2|\sin(\vartheta/2)|\le2|\sin(3\vartheta/2)|$. Let $c=f(w)$. Then $c^3=I$, and $xcx=f(xwx)=f(w^*)=c^*=c^{-1}$. As before
$\|(c-w)\xi\|_F\le E(w^3)\le25\delta$ and $\|\xi(c-w)\|_F\le25\delta$. Hence $\|(c-Z)\xi\|_F\le29\delta$ and
$\chi(c)\le\chi(w)+\|(c-w)\xi\|_F+\|\xi(c-w)\|_F\le56\delta$.
\end{proof}

\subsection{A weighted sector bound}

Rounded copies of $S_3$ that nearly commute in the weighted sense force entropy. We work in the Hilbert space $\mathcal H$ of operators on
$\mathbb C^M$ with the Hilbert--Schmidt inner product $\langle Y,Z\rangle=\Tr(Y^*Z)$, identified with $\mathbb C^M\otimes\mathbb C^M$. Left multiplications
act on the first factor, which we call the system, and right multiplications on the second, the reference; they commute. The unit vector $\xi$
purifies $\rho$ on the system, and its reference marginal has the same nonzero spectrum, so it has entropy $S(\rho)$.

\begin{lemma}[Projection products~{\cite[Lemma~4.1]{DouglasMghirbiA}}]\label{lem:projection-product}
Let $P_1,\dots,P_k$ be orthogonal projections on a Hilbert space. Every vector $y$ satisfies
$\|y\|^2-\mathrm{Re}\langle y,P_k\cdots P_1y\rangle\le2\sum_l\|(I-P_l)y\|^2$.
\end{lemma}

\begin{proof}
By homogeneity let $\|y\|=1$. Put $y_0=y$, $y_l=P_ly_{l-1}$ and $s_l=y_{l-1}-y_l$, which lies in the range of $I-P_l$. Then
$\sum_l\|s_l\|^2=1-\|y_k\|^2=:\beta$, and $\alpha:=1-\mathrm{Re}\langle y,y_k\rangle=\mathrm{Re}\sum_l\langle(I-P_l)y,s_l\rangle\le\sqrt{E\beta}$ with
$E=\sum_l\|(I-P_l)y\|^2$. Since $0\le\|y-y_k\|^2=2\alpha-\beta$, we get $\alpha\le\sqrt{2E\alpha}$, that is $\alpha\le2E$.
\end{proof}

Let $c_1,\dots,c_N$ be unitaries on $\mathbb C^M$ with $c_i^3=I$, and write $c_i=\sum_{k\in\Z/3}\omega^kP_i^k$ with $\omega=e^{2\pi i/3}$ and spectral
projections $P_i^k$. For every $Y\in\mathcal H$,
\begin{equation}\label{eq:spectral-identities}
\|(c_i-I)Y\|_F^2=3\sum_{k\ne0}\|P_i^kY\|_F^2,\qquad\|[c_i,Y]\|_F^2=3\sum_{k\ne l}\|P_i^kYP_i^l\|_F^2,
\end{equation}
because $|\omega^k-\omega^l|^2=3$ for $k\ne l$ and the pieces $P_i^kYP_i^l$ are orthogonal. For $\mathbf a=(a_1,\dots,a_{i-1})\in(\Z/3)^{i-1}$ put
$Q_{\mathbf a}=P_{i-1}^{a_{i-1}}\cdots P_1^{a_1}$. Summing over the coordinates one at a time gives
$\sum_{\mathbf a}Q_{\mathbf a}^*Q_{\mathbf a}=\sum_{\mathbf a}Q_{\mathbf a}Q_{\mathbf a}^*=I$.

\begin{lemma}[Transfer to the reference]\label{lem:transfer}
For every $Y\in\mathcal H$ and $1\le i\le N$,
\[
\sum_{\mathbf a}\|Q_{\mathbf a}Y-YQ_{\mathbf a}\|_F^2\ \le\ \frac43\sum_{k<i}\|[c_k,Y]\|_F^2 .
\]
\end{lemma}

\begin{proof}
The pinchings $\mathcal E_k(Y)=\sum_lP_k^lYP_k^l$ are orthogonal projections on $\mathcal H$, and by~\eqref{eq:spectral-identities}
$\|(I-\mathcal E_k)Y\|_F^2=\sum_{l\ne l'}\|P_k^lYP_k^{l'}\|_F^2=\frac13\|[c_k,Y]\|_F^2$. Both $\sum_{\mathbf a}\|Q_{\mathbf a}Y\|_F^2$ and
$\sum_{\mathbf a}\|YQ_{\mathbf a}\|_F^2$ equal $\|Y\|_F^2$, and summing over the coordinates one at a time,
$\sum_{\mathbf a}\langle YQ_{\mathbf a},Q_{\mathbf a}Y\rangle=\sum_{\mathbf a}\Tr(Q_{\mathbf a}^*Y^*Q_{\mathbf a}Y)=\langle\mathcal E_1\cdots\mathcal E_{i-1}Y,Y\rangle$.
So the left side equals $2\|Y\|_F^2-2\,\mathrm{Re}\langle Y,\mathcal E_1\cdots\mathcal E_{i-1}Y\rangle$, and Lemma~\ref{lem:projection-product} bounds it by
$4\sum_{k<i}\|(I-\mathcal E_k)Y\|_F^2$.
\end{proof}

A three-valued distribution that is nearly symmetric under $k\mapsto-k$ carries about one bit for each unit of weight off $0$. Let $h_2$ be the
binary entropy. With $t=(1+y)/2$,
\[
\ln2\,\bigl(1-h_2(t)\bigr)=\sum_{k\ge1}\frac{y^{2k}}{2k(2k-1)}\ \le\ \sum_{k\ge1}\binom{2k}k\frac{y^{2k}}{(2k-1)4^k}=1-\sqrt{1-y^2}=1-2\sqrt{t(1-t)},
\]
comparing coefficients: $4^k\le2k\binom{2k}k$ holds for $k=1$, and the right side grows by the factor $2(2k+1)/k\ge4$ from $k$ to $k+1$. For a
distribution $(r_0,r_1,r_2)$ with $s=r_1+r_2>0$ this gives, with $t=r_1/s$,
\begin{equation}\label{eq:binary-hellinger}
H(r_0,r_1,r_2)\ \ge\ s\,h_2(r_1/s)\ \ge\ s-\frac{(\sqrt{r_1}-\sqrt{r_2})^2}{\ln2},
\end{equation}
which also holds when $s=0$.

\begin{theorem}[Weighted sector bound]\label{thm:sector}
Let $c_i,x_i$, $1\le i\le N$, be unitaries on $\mathbb C^M$ with $c_i^3=x_i^2=I$ and $x_ic_ix_i=c_i^{-1}$, with no condition relating different
indices. For a state $\rho$ with $\xi=\sqrt\rho$ put $\Delta_i=E(c_i)$, $\theta_i=\chi(c_i)$, $\theta'_i=\chi(x_i)$, and
\[
T_i=\Bigl(\sum_{k<i}\theta_k^2\Bigr)^{1/2},\qquad K_i=\Bigl(\sum_{k<i}\|[c_i,c_k]\xi\|_F^2\Bigr)^{1/2},\qquad K'_i=\Bigl(\sum_{k<i}\|[x_i,c_k]\xi\|_F^2\Bigr)^{1/2},
\]
\[
w_i=\frac13\Bigl(\Delta_i-\frac4{\sqrt3}T_i\Bigr)_+^2,\qquad\beta_i=\theta'_i+\frac2{\sqrt3}(2T_i+K'_i),\qquad
\pi_i=\frac13\Bigl(\theta_i+\frac2{\sqrt3}(2T_i+K_i)\Bigr)^2 .
\]
If $\pi_i\le2/3$ for all $i$, then
\[
S(\rho)\ \ge\ \sum_{i=1}^N\Bigl[w_i-\frac{\beta_i^2}{2\ln2}-h_2(\pi_i)-\pi_i\Bigr].
\]
\end{theorem}

\begin{proof}
\emph{Sequential measurement.} Measure $c_1,\dots,c_N$ in turn on the system, with the projections $P_i^k$, and let $A_i\in\Z/3$ be the outcomes.
Then $\Pr(A_{<i}=\mathbf a,A_i=k)=\|P_i^kQ_{\mathbf a}\xi\|_F^2=:p(\mathbf a,k)$. The measurements do not touch the reference, whose entropy is
$S(\rho)$, so
\[
S(\rho)\ \ge\ I(A_1\cdots A_N:\mathrm{Ref})=\sum_{i=1}^NI(A_i:\mathrm{Ref}\mid A_{<i}),
\]
since the conditional entropy of the reference given the classical outcomes is nonnegative. Fix $i$.

\emph{The copy stays displaced.} By~\eqref{eq:spectral-identities}, $\Pr(A_i\ne0)=\frac13\sum_{\mathbf a}\|(c_i-I)Q_{\mathbf a}\xi\|_F^2$. By
Lemma~\ref{lem:transfer} with $Y=\xi$, the vectors $Q_{\mathbf a}\xi$ and $\xi Q_{\mathbf a}$ differ by at most $\frac2{\sqrt3}T_i$ in $\ell^2$ over $\mathbf a$,
and $\sum_{\mathbf a}\|(c_i-I)\xi Q_{\mathbf a}\|_F^2=\Delta_i^2$ because $\sum_{\mathbf a}Q_{\mathbf a}Q_{\mathbf a}^*=I$. Since $\|c_i-I\|\le2$, the triangle
inequality in $\ell^2$ gives $\Pr(A_i\ne0)\ge w_i$.

\emph{The sign is a fair coin.} Since $x_iP_i^kx_i=P_i^{-k}$ and right multiplication by $x_i$ is unitary,
$\sqrt{p(\mathbf a,-k)}=\|P_i^kx_iQ_{\mathbf a}\xi\|_F$ and $\sqrt{p(\mathbf a,k)}=\|P_i^kQ_{\mathbf a}\xi x_i\|_F$, so
$\sum_{\mathbf a,k}(\sqrt{p(\mathbf a,k)}-\sqrt{p(\mathbf a,-k)})^2\le\sum_{\mathbf a}\|x_iQ_{\mathbf a}\xi-Q_{\mathbf a}\xi x_i\|_F^2$. Write
\[
x_iQ_{\mathbf a}\xi-Q_{\mathbf a}\xi x_i=x_i(Q_{\mathbf a}\xi-\xi Q_{\mathbf a})+\bigl((x_i\xi)Q_{\mathbf a}-Q_{\mathbf a}(x_i\xi)\bigr)+Q_{\mathbf a}(x_i\xi-\xi x_i).
\]
In $\ell^2$ over $\mathbf a$ the first term is at most $\frac2{\sqrt3}T_i$ and the third is $\theta'_i$. For the second apply Lemma~\ref{lem:transfer} with
$Y=x_i\xi$: since $[c_k,x_i\xi]=[c_k,x_i]\xi+x_i[c_k,\xi]$, it is at most $\frac2{\sqrt3}(K'_i+T_i)$. So the sum is at most $\beta_i^2$, and since the
terms $k=1$ and $k=2$ contribute equally, \eqref{eq:binary-hellinger} gives
\[
H(A_i\mid A_{<i})\ \ge\ \Pr(A_i\ne0)-\frac1{\ln2}\sum_{\mathbf a}\bigl(\sqrt{p(\mathbf a,1)}-\sqrt{p(\mathbf a,2)}\bigr)^2\ \ge\ w_i-\frac{\beta_i^2}{2\ln2}.
\]

\emph{The reference knows the outcome.} Let the reference measure $c_i$ from the right, with the projections $Y\mapsto YP_i^l$, and call the outcome
$B_i$. Left and right multiplications commute, so $\Pr(A_{<i}=\mathbf a,A_i=k,B_i=l)=\|P_i^kQ_{\mathbf a}\xi P_i^l\|_F^2$, and
by~\eqref{eq:spectral-identities} $\Pr(A_i\ne B_i)=\frac13\sum_{\mathbf a}\|[c_i,Q_{\mathbf a}\xi]\|_F^2$. As above,
\[
[c_i,Q_{\mathbf a}\xi]=c_i(Q_{\mathbf a}\xi-\xi Q_{\mathbf a})+\bigl((c_i\xi)Q_{\mathbf a}-Q_{\mathbf a}(c_i\xi)\bigr)+Q_{\mathbf a}(c_i\xi-\xi c_i),
\]
and Lemma~\ref{lem:transfer} with $Y=\xi$ and $Y=c_i\xi$ gives $\Pr(A_i\ne B_i)\le\pi_i$. By data processing and Fano's inequality for three values,
\[
I(A_i:\mathrm{Ref}\mid A_{<i})\ \ge\ I(A_i:B_i\mid A_{<i})\ \ge\ H(A_i\mid A_{<i})-h_2(\pi_i)-\pi_i ,
\]
where we used that $h_2(\pi)+\pi$ increases on $[0,2/3]$. Summing over $i$ proves the theorem.
\end{proof}

\begin{corollary}\label{cor:sector}
In Theorem~\ref{thm:sector}, suppose $\Delta_i\ge1/2$ for all $i$, and that $\theta_i$, $\theta'_i$, $\|[c_i,c_k]\xi\|_F$ and $\|[x_i,c_k]\xi\|_F$ are at
most $b$ for all $i\ne k$, where $\sqrt N\,b\le10^{-3}$. Then $S(\rho)\ge N/16$.
\end{corollary}

\begin{proof}
Now $T_i,K_i,K'_i\le\sqrt N\,b\le10^{-3}$ and $\theta_i,\theta'_i\le10^{-3}$. So $w_i\ge\frac13(\frac12-\frac4{\sqrt3}10^{-3})^2>0.0825$, and $\beta_i$ and
$\sqrt{3\pi_i}$ are at most $10^{-3}+2\sqrt3\cdot10^{-3}<4.5\cdot10^{-3}$. Hence $\beta_i^2/(2\ln2)<1.5\cdot10^{-5}$ and $\pi_i<6.8\cdot10^{-6}$, and with
$h_2(\pi)\le\pi\log(e/\pi)$ we get $h_2(\pi_i)<1.3\cdot10^{-4}$. Each summand exceeds $0.0823>1/16$.
\end{proof}

\subsection{Completion}

\begin{proof}[Proof of Theorem~\ref{thm:weighted}]
Let $\phi$ be given by Lemma~\ref{lem:weighted-extraction}, with $e=256\,m\sqrt{d\nu}$, and put
\[
b_*=100\,m\,(A+L+1)\,e=25600\,m^2(A+L+1)\sqrt{d\nu},\qquad\delta_L=(2L+1)e .
\]
Condition~\eqref{eq:weighted-condition} gives $\sqrt N\,b_*\le25600\cdot2^{-25}<10^{-3}$.

\emph{Copies.} Let $\psi_i=\phi(u_i)$, so $\chi(\psi_i)\le Le$, and put $X_i=\psi_i\phi(a)\psi_i^*$, $Y_i=\psi_i\phi(b)\psi_i^*$ and $Z_i=X_iY_i$. Since
$a^2$, $b^2$ and $(ab)^3$ lie in $R$, (W4) and Lemma~\ref{lem:weighted-extraction} give $E(X_i^2),E(Y_i^2),E(Z_i^3)\le e+2Le=\delta_L$, and (W2) gives
$\chi(X_i),\chi(Y_i)\le\delta_L$. Lemma~\ref{lem:rounding} gives unitaries $x_i,c_i$ with $x_i^2=c_i^3=I$, $x_ic_ix_i=c_i^{-1}$ and
\[
\|(x_i-X_i)\xi\|_F\le\delta_L,\qquad\chi(x_i)\le3\delta_L,\qquad\|(c_i-Z_i)\xi\|_F\le29\delta_L,\qquad\chi(c_i)\le56\delta_L .
\]
By (W4) and Lemma~\ref{lem:weighted-extraction}(iii), $E(Z_i)\ge E(\phi(a)\phi(b))-2Le\ge1-\delta_L$, so $\Delta_i=E(c_i)\ge1-30\delta_L$.

\emph{Raw commutators.} Let $i\ne k$ and $x,y\in\{a,b\}$, and let $P=\psi_i\phi(x)\psi_i^*$ and $P'=\psi_k\phi(y)\psi_k^*$. The group commutator
$g=[u_ixu_i^{-1},u_kyu_k^{-1}]$ is a word of length at most $8L+4$ and area at most $A$, so Lemma~\ref{lem:charged-area} gives
$E(\phi(g))\le(3mA+8L+4)e$. Since $PP'-P'P=(\phi(g)-I)P'P$ and $\chi(P'P)\le2\delta_L$, (W3) gives
\[
\|(PP'-P'P)\xi\|_F\le\gamma:=(3mA+8L+4)e+4\delta_L=(3mA+16L+8)e .
\]
Expanding $[X_i,X_kY_k]=[X_i,X_k]Y_k+X_k[X_i,Y_k]$ and $[X_iY_i,Z_k]=X_i[Y_i,Z_k]+[X_i,Z_k]Y_i$ and using (W3),
$\|[X_i,Z_k]\xi\|_F,\|[Y_i,Z_k]\xi\|_F\le2\gamma+2\delta_L$ and $\|[Z_i,Z_k]\xi\|_F\le4\gamma+6\delta_L$.

\emph{Rounded commutators.} For unitaries $U,V,U_0,V_0$,
\[
[U,V]-[U_0,V_0]=U(V-V_0)+(U-U_0)V_0-V(U-U_0)-(V-V_0)U_0,
\]
so by (W3) the weighted norms of the two commutators differ by at most
$2\|(U-U_0)\xi\|_F+2\|(V-V_0)\xi\|_F+2\chi(U_0)+2\chi(V_0)$. With $(U,U_0,V,V_0)=(x_i,X_i,c_k,Z_k)$ and $(c_i,Z_i,c_k,Z_k)$ this gives
\begin{gather*}
\|[x_i,c_k]\xi\|_F\le2\gamma+68\delta_L=(6mA+168L+84)e,\\
\|[c_i,c_k]\xi\|_F\le4\gamma+130\delta_L=(12mA+324L+162)e .
\end{gather*}
Since $m\ge6$, these and $\chi(c_i)\le(112L+56)e$ and $\chi(x_i)\le(6L+3)e$ are at most $b_*$, and $\Delta_i\ge1-(60L+30)e\ge1-b_*>1/2$.
Corollary~\ref{cor:sector} gives $S(\rho)\ge N/16$.
\end{proof}

The constant $2^{25}$ is not optimized; the proof uses only $25600\cdot2^{-25}<10^{-3}$, and Corollary~\ref{cor:sector} has room to spare.

%======================================================================
\section{Local physics}\label{sec:physics}
%======================================================================

Theorem~B concerns an abstract channel. Here we record what it implies for physical systems that could produce the channel, and what
autonomous dynamics can and cannot produce in the long run. We call a channel \emph{cheap} if for all $n$ and $\epsilon$ it has causal
services for $n$ rounds with error at most $\epsilon$, $P=C=0$ and memory polynomial in $\log(n/\epsilon)$, and \emph{expensive} if for some
$A,c>0$ every service for $n$ rounds with error at most $1/16$ and $P+C\le A\log(n+1)$ needs memory at least $n^c$ for all large $n$. Neither notion
constrains the exchange.

\paragraph{Local reservoirs.} A \emph{local reservoir} consists of a finite box of the lattice $\Z^\nu$, with a system of dimension at most $q$ at
each site, the probe $\mathbb C^d$ attached to the origin, and a Hamiltonian $H(t)=\sum_Xh_X(t)$, possibly time dependent, whose terms are supported
in sets $X$ of diameter at most $R_0$ and satisfy $\|h_X(t)\|\le J_0$. The reservoir starts in its tracial state, the probe is coupled to it for a
time $T$, and $\Phi_T(X)=(\mathrm{id}\otimes\tau)[W_T(X\otimes I)W_T^*]$ is the resulting channel on the probe, where $W_T$ is the time-ordered
evolution and $\tau$ the normalized trace of the reservoir. We call $\nu,q,R_0,J_0$ and $d$ the \emph{parameters} of the reservoir; the size of the
box is arbitrary.

\begin{proposition}[Local reservoirs are cheap at every finite time]\label{prop:light-cone}
There is $K$, depending only on the parameters, such that for every local reservoir, every $T\ge0$, every $n\ge1$ and every $0<\epsilon\le1$, the
channel $\Phi_T$ has a causal service for $n$ rounds with error at most $\epsilon$, $P=C=0$, $J=(n-1)Q$ and
\[
Q\ \le\ K\,\bigl[1+T+\log(n/\epsilon)\bigr]^\nu .
\]
\end{proposition}

\begin{proof}
For $r\ge0$ let $\Phi_T^{(r)}$ be the channel produced by the terms $h_X$ with $X$ inside the ball $B_r$ of radius $r$ about the origin. The
Lieb--Robinson bound~\cite{LiebRobinson1972,BHV2006,NSY2019} gives constants $C_1,v,\mu>0$, depending only on the parameters, such that the Heisenberg
evolutions of an operator $A$ at the origin under the full and the truncated dynamics differ in norm by at most $C_1\|A\|e^{vT-\mu r}$. The bound
does not depend on the dimension of the site at the origin, so it holds for operators on the probe tensored with a reference system that the
dynamics does not touch. Hence the Heisenberg-picture maps of $\Phi_T$ and $\Phi_T^{(r)}$ differ in completely bounded norm by at most
$C_1e^{vT-\mu r}$, and by duality $\ddia(\Phi_T,\Phi_T^{(r)})\le C_1e^{vT-\mu r}$. The truncated evolution acts only on the probe and the sites
of $B_r$, and the other sites start in their tracial state, so $\Phi_T^{(r)}$ has a finite tracial factorization with bath dimension
$m\le q^{(2r+1)^\nu}$.

Choose $r=\lceil(vT+\ln(2C_1n/\epsilon))/\mu\rceil$, so that $\ddia(\Phi_T,\Phi_T^{(r)})\le\epsilon/(2n)$, and let $Q=\lceil\log m+\log(2n/\epsilon)\rceil$.
The exchange construction in the proof of Proposition~\ref{prop:sofic-upper}, with $U_f$ replaced by the unitary of this factorization and
$\vartheta=0$, gives a service of $\Phi_T^{(r)}$ in which every round is within $\epsilon/(2n)$ of $\Phi_T^{(r)}$, hence within $\epsilon/n$ of
$\Phi_T$; a hybrid argument bounds the error by $\epsilon$. Finally $Q\le(2r+1)^\nu\log q+\log(2n/\epsilon)+1\le K[1+T+\log(n/\epsilon)]^\nu$.
\end{proof}

With Theorem~B this says that the channel of the configuration lamps cannot be produced quickly by local physics.

\begin{corollary}[Preparation time]\label{cor:preparation}
Let $\Phi$ be the channel of Theorem~\ref{thm:main-channel}, and fix the parameters of a local reservoir with $d=873$. There are $c,\delta_0>0$ such
that, whenever such a reservoir produces at time $T$ a channel within $\delta\le\delta_0$ of $\Phi$ in diamond distance,
\[
T\ \ge\ c\,\delta^{-1/\nu}\bigl(\log(1/\delta)\bigr)^{-(2\kappa+1)/\nu},\qquad 2\kappa+1<26.49 .
\]
\end{corollary}

\begin{proof}
Let $n=\lfloor1/(32\delta)\rfloor$. Proposition~\ref{prop:light-cone} gives a service of $\Phi_T$ with error at most $1/32$, zero purity and
$Q\le K[1+T+\log(32n)]^\nu$. Replacing the $n$ uses of $\Phi_T$ by uses of $\Phi$ one at a time changes the observer's final state by at most
$n\delta\le1/32$, so the same device serves $\Phi$ with error at most $1/16$ and zero purity. For $\delta_0$ small, $n$ is large, and
Theorem~\ref{thm:main-channel}(a) with $A_0=0$ gives $Q\ge c'n/(\log n)^{2\kappa+1}$. Hence
$1+T+\log(32n)\ge\bigl(c'n/(K(\log n)^{2\kappa+1})\bigr)^{1/\nu}$, and since $n\ge1/(64\delta)$ this gives the claim for $\delta_0$ small.
\end{proof}

The maximally mixed start of the reservoir is essential: $\Phi$ has a Stinespring dilation with a pure environment of dimension $r_*$, so
from a pure state a bounded number of sites produce $\Phi$ exactly in bounded time. The same argument applies to a circuit of depth $T$ of gates of
bounded range, whose light cone is exact, and so bounds the depth of any local
circuit that prepares $\Phi$. The corollary does not exclude that local dynamics produce $\Phi$ in the limit of long times; it says how slowly.

\paragraph{Autonomous dynamics.} Let $\mathcal B$ be a von Neumann algebra with a faithful normal tracial state $\tau_{\mathcal B}$, put
$\mathcal M=M_d\otimes\mathcal B$ with the trace $\tau=\mathrm{tr}_d\otimes\tau_{\mathcal B}$, where $\mathrm{tr}_d$ is the normalized trace, and let
$\iota(a)=a\otimes1$ and $E=\mathrm{id}\otimes\tau_{\mathcal B}$. An \emph{autonomous dynamics} is a trace-preserving automorphism $\alpha$ of $\mathcal M$,
and the channel after $t\in\N$ steps is $\Phi_t=E\circ\alpha^t\circ\iota$; for $\alpha=\mathrm{Ad}\,U$ with $\mathcal B=M_m$ this is the tracial
factorization $\Phi_t(X)=(\mathrm{id}\otimes\tau_m)[U^t(X\otimes I)U^{-t}]$. Let $\mathcal N$ be the fixed-point algebra of $\alpha$ and
$E_{\mathcal N}$ the $\tau$-preserving conditional expectation onto it.

The next proposition combines the mean ergodic theorem with the observation of Junge, Le Merdy and Xu~\cite[Remark~10.3]{JungeLeMerdyXu2006}
and of Haagerup and Musat~\cite[Remark~5.2]{HaagerupMusat2011} that compressions $E\circ E_{\mathcal N}\circ\iota$ of trace-preserving conditional
expectations are self-adjoint and positive; both of its clauses are standard.

\begin{proposition}[{Autonomous limits are self-adjoint; the mean ergodic theorem and~\cite[Remark~10.3]{JungeLeMerdyXu2006},
\cite[Remark~5.2]{HaagerupMusat2011}}]\label{prop:autonomous}
The averages $\frac1T\sum_{t=0}^{T-1}\Phi_t$ converge in diamond norm to $\overline\Phi=E\circ E_{\mathcal N}\circ\iota$, and if $\Phi_t$ converges, its
limit is $\overline\Phi$. The channel $\overline\Phi$ is self-adjoint and positive semidefinite for the Hilbert--Schmidt inner product on $M_d$. In
particular, a channel $\Psi$ with $\Psi^*\ne\Psi$ is not the long-time limit of any autonomous dynamics with a tracial bath, local or not.
\end{proposition}

\begin{proof}
The automorphism extends to a unitary $V$ on $L^2(\mathcal M,\tau)$. By the mean ergodic theorem~\cite{Krengel1985}, $\frac1T\sum_{t<T}V^t$ converges
strongly to the orthogonal projection $P$ onto the invariant vectors. For $x\in\mathcal M$ the averages of $x$ lie in the ball of radius $\|x\|$ of
$\mathcal M$, which is closed in $L^2$, so $Px\in\mathcal M$, and $Px$ is invariant; hence $P(\mathcal M)\subseteq\mathcal N$. Since
$\mathcal N$ consists of invariant vectors and $\mathcal M$ is dense in $L^2(\mathcal M)$, $P$ is the orthogonal projection onto
$L^2(\mathcal N)$, whose restriction to $\mathcal M$ is $E_{\mathcal N}$. The map $\iota$ is an isometry from $(M_d,\mathrm{tr}_d)$ into
$L^2(\mathcal M,\tau)$ with adjoint $E$, so $\Phi_t=\iota^*V^t\iota$, and the averages converge to $\iota^*P\iota=\overline\Phi$ on every $a\in M_d$,
hence in diamond norm, since $M_d$ is finite dimensional. An ordinary limit equals the limit of the averages. Finally
$\mathrm{tr}_d(b^*\overline\Phi(a))=\langle P\iota(b),P\iota(a)\rangle$ is Hermitian in $(a,b)$ and nonnegative for $a=b$.
\end{proof}

Self-adjointness and positivity are not the only constraints. The limit $\overline\Phi$ is a compression of a conditional expectation, and this ties
the decay of a coherence to the return probability of a state.

\begin{proposition}[A return inequality]\label{prop:return}
Let $T=E\circ E_{\mathcal N}\circ\iota$ be as in Proposition~\ref{prop:autonomous}, on $M_D$. If $T(e_{00})=e_{00}$ and $T(e_{0j})=a\,e_{0j}$ for some
$j\ne0$ and $0\le a\le1$, then
\[
\langle j|T(e_{jj})|j\rangle\ \ge\ a^2+\frac{(1-a)^2}{D-1}.
\]
\end{proposition}

\begin{proof}
Write $P=E_{\mathcal N}$, the orthogonal projection onto $L^2(\mathcal N)$, and identify $x\in M_D$ with $\iota(x)$; then
$\mathrm{tr}_D(x^*T(y))=\langle x,Py\rangle$ with $\langle x,y\rangle=\tau(x^*y)$. Put $p_0=e_{00}$, $p_j=e_{jj}$ and $x=e_{0j}$. Since
$\|Pp_0\|_2^2=\langle p_0,T(p_0)\rangle=\|p_0\|_2^2$, we have $Pp_0=p_0$, so $p_0\in\mathcal N$. Let $y=P(x)$, $X=P(p_j)$ and $B=y^*y$. As $P$ is
$\mathcal N$-bimodular, $p_0y=P(p_0x)=y$, $yp_0=P(xp_0)=0$ and $p_0X=P(p_0p_j)=0$; so $yy^*$ is supported in $p_0$, and $B$ and $X$ in $1-p_0$. The
Schwarz inequality for conditional expectations gives $B\le P(x^*x)=X$. Moreover $\tau(X)=\tau(p_j)=1/D$ and $\tau(B)=\|Px\|_2^2=\langle x,T(x)\rangle=a/D$,
so $R=X-B\ge0$ has $\tau(R)=(1-a)/D$. By the Cauchy--Schwarz inequality, $\tau(B^2)=\tau((yy^*)^2)\ge\tau(yy^*)^2/\tau(p_0)=a^2/D$ and
$\tau(R^2)\ge\tau(R)^2/\tau(1-p_0)=(1-a)^2/(D(D-1))$, and $\tau(BR)\ge0$. Hence
\[
\frac1D\,\langle j|T(e_{jj})|j\rangle=\langle p_j,Pp_j\rangle=\tau(X^2)=\tau(B^2)+2\tau(BR)+\tau(R^2)\ \ge\ \frac{a^2}D+\frac{(1-a)^2}{D(D-1)} .\qedhere
\]
\end{proof}

\begin{remark}[Self-adjoint and positive is not enough]\label{rem:qutrit}
Haagerup and Musat give self-adjoint positive maps that are not such compressions because they are not
factorizable~\cite[Theorem~5.4]{HaagerupMusat2011}. The following channel is an exact finite factor and is still excluded. Let $a=\sqrt3/2$ and $b=\frac12$, and let $\Lambda$ on $M_3$ have the Kraus operators $\mathrm{diag}(1,a,a)$, $b\,e_{12}$ and $-b\,e_{21}$. With
$\omega=e^{2\pi i/3}$ and $W_k=1\oplus\bigl(\begin{smallmatrix}a&b\omega^k\\-b\omega^{-k}&a\end{smallmatrix}\bigr)$ one has
$\Lambda=\frac13\sum_{k=0}^2\mathrm{Ad}\,W_k$, so $\Lambda$ has a finite tracial factorization with bath dimension three. It is self-adjoint, and positive
for the Hilbert--Schmidt inner product, with eigenvalues between $\frac12$ and $1$. But $\Lambda(e_{00})=e_{00}$, $\Lambda(e_{01})=a\,e_{01}$ and
$\langle1|\Lambda(e_{11})|1\rangle=a^2$, which violates Proposition~\ref{prop:return}. So $\Lambda$ is not the long-time limit of any autonomous dynamics
with a tracial bath.
\end{remark}

Conversely, every channel of the form $\Psi\Psi^*$ with $\Psi$ in the tracial closure is a compression of a conditional
expectation (apply~\cite[Theorem~5.9]{HaagerupMusat2011} to $\Psi^*$, which is again factorizable), and in fact an autonomous long-time limit, of a dynamics that the construction does not make
local; the construction adapts Kümmerer's tensor dilations with a Bernoulli shift~\cite{Kummerer1985}.

\begin{proposition}[Autonomy without locality]\label{prop:nonlocal}
Let $\mathcal N_0$ be a von Neumann algebra with a faithful normal tracial state $\tau_0$, let $U$ be a unitary in $M_d\otimes\mathcal N_0$, and put
$\Psi(X)=(\mathrm{id}\otimes\tau_0)[U(X\otimes1)U^*]$. Let $\mathcal B=\overline\bigotimes_{k\in\Z}(\mathcal N_0,\tau_0)$, with $U$ acting on the factor
$k=0$, and let $\beta$ be the Bernoulli shift of $\mathcal B$. Then the autonomous dynamics $\alpha=\mathrm{Ad}\,U\circ(\mathrm{id}\otimes\beta)\circ\mathrm{Ad}\,U^*$
has $\Phi_t\to\Psi\Psi^*$, where $\Psi^*(a)=(\mathrm{id}\otimes\tau_0)[U^*(a\otimes1)U]$ is the Hilbert--Schmidt adjoint of $\Psi$.
\end{proposition}

\begin{proof}
We have $\alpha^t=\mathrm{Ad}\,U\circ(\mathrm{id}\otimes\beta^t)\circ\mathrm{Ad}\,U^*$. The shift is mixing: $\tau_{\mathcal B}(w\,\beta^t(z))\to\tau_{\mathcal B}(w)\tau_{\mathcal B}(z)$ for
$w,z\in\mathcal B$, since equality holds as soon as $w$ and $\beta^t(z)$ are elementary tensors with disjoint supports, and such tensors of norm
at most $\|w\|$ and $\|z\|$ approximate $w$ and $z$ in $2$-norm (Kaplansky's density theorem). Let $a\in M_d$ and $y=U^*(a\otimes1)U$, and write
$y=\sum_{i,j}e_{ij}\otimes y_{ij}$. For every $x=\sum_{k,l}e_{kl}\otimes x_{kl}$ in $M_d\otimes\mathcal B$,
\[
\tau\bigl(x\,(\mathrm{id}\otimes\beta^t)(y)\bigr)=\frac1d\sum_{k,l}\tau_{\mathcal B}\bigl(x_{kl}\beta^t(y_{lk})\bigr)\ \longrightarrow\ \tau\bigl(x\,(E(y)\otimes1)\bigr).
\]
Taking $x=U^*(e_{ji}\otimes1)U$ shows that every matrix entry of $\Phi_t(a)=E[U(\mathrm{id}\otimes\beta^t)(y)U^*]$ converges to the corresponding entry
of $E[U(E(y)\otimes1)U^*]=\Psi(\Psi^*(a))$. Convergence on a basis of $M_d$ gives convergence in diamond norm.
\end{proof}

By a theorem of Haagerup and Musat~\cite[Theorem~3.6]{HaagerupMusat2015}, every channel in $\cK_d$ has the form $\Psi$ of
Proposition~\ref{prop:nonlocal} for some $\mathcal N_0$, and the channel of Theorem~B has it with $\mathcal N_0=L(\Gamma)$.

Nothing in Proposition~\ref{prop:nonlocal} makes the dynamics local: $U$ is an arbitrary unitary of $M_d\otimes\mathcal N_0$, and we do not know
whether a dynamics of this kind can be given by a local rule.

\begin{remark}[Self-adjoint expensive channels]\label{rem:doubling}
Self-adjointness is no obstacle to expense. Let $z$ be the unitary of order three in $L(\Z_3)$, so that $\tau(z)=\tau(z^2)=0$, let $U$ be the gadget
unitary of the channel $\Phi$ of Theorem~B, and let
$V=\bigl(\begin{smallmatrix}0&U\otimes z\\U^*\otimes z^*&0\end{smallmatrix}\bigr)$, a self-adjoint unitary over $\mathbb C^2\otimes\mathbb C^{873}$ and
$L(\Gamma\times\Z_3)$. The channel it defines on $M_{1746}$ is
$\Theta\bigl(\begin{smallmatrix}x&y\\u&w\end{smallmatrix}\bigr)=\mathrm{diag}\bigl(\Phi(w),\Phi^*(x)\bigr)$, which is self-adjoint. It keeps the lower bounds of
Theorem~B. For part~(a), a service of $\Theta$ becomes a service of $\Phi$ with one more memory qubit and one more bit of purity, at the same
exchange, by feeding each input together with a retained tag qubit that starts in $|1\rangle$ and is flipped after every round, since $\Theta$ maps
$|1\rangle\langle1|\otimes\rho$ to $|0\rangle\langle0|\otimes\Phi(\rho)$. For part~(b), $\Theta$ has a flat probe of rank $2r_*$, with anchors in both blocks,
so Proposition~\ref{prop:purity-budget} gives $P+C<3Q+6$ at its own rank rate $J\le\frac n2\log(2r_*)$; the tag transfer and part~(a) then give
$Q\ge c\sqrt n/(\log n)^\kappa$. But $\Theta$ maps $\mathrm{diag}(I,-I)$ to $-\mathrm{diag}(I,-I)$, so it is not
positive and not an autonomous limit, while $\Theta^2=\mathrm{diag}(\Phi\Phi^*,\Phi^*\Phi)$ is one, by Proposition~\ref{prop:nonlocal}. The next
proposition shows that $\Theta^2$ is cheap.
\end{remark}

Composing a gadget channel with its adjoint erases the multiplication of the group. In a triangle block the unitarity of the gadget forces the
coefficient of the product label to be the product of the coefficients, and this is what the one-round theorem reads. In $\Phi\Phi^*$ a triangle only
requires a permutation $\pi_x$ of the labels that carries $1$ to $x$ and $y$ to $xy$, and no relation between the permutations $\pi_x$ for
different $x$ is needed.

\begin{proposition}[Adjoint compositions of gadgets are exact finite factors]\label{prop:compositions}
Let $\Phi_G$ be a gadget channel with $r_*$ labels. Then $\Phi_G\Phi_G^*$ and $\Phi_G^*\Phi_G$ have finite tracial factorizations with baths of at
most $3r_*-2$ qubits. Consequently, for every $n$, they have exact causal services with $Q\le3r_*-2$, $P=C=0$ and $J=(n-1)Q$, and so does
$\Theta^2$ in Remark~\ref{rem:doubling}, with one more qubit.
\end{proposition}

\begin{proof}
Let $\mathcal F$ be the set of labels, and write $A_g=\sum_{[w_a]=g}c_aE_a$, so that $\Phi_G\Phi_G^*(X)=\sum_{g,h}A_gA_h^*XA_hA_g^*$. A product $A_gA_h^*$
is a sum of terms $c_a\overline{c_b}E_aE_b^*$ over slots $a,b$ in a common column, so it vanishes unless $g$ and $h$ label two slots of one column. In a
triangle block with labels $1,y$ in the first row and $x,z$ in the second, where $z=xy$, the columns carry the labels $\{1,x\}$ and $\{y,z\}$; so
$A_gA_h^*\ne0$ forces $hg^{-1}\in K=\{1\}\cup\{x^{\pm1}\}$, the union over the triangles. For $k\in K$ let $\pi_k$ be a permutation of $\mathcal F$ that
extends the partial map $g\mapsto kg$ (defined when $g,kg\in\mathcal F$), with $\pi_1=\mathrm{id}$ and $\pi_{k^{-1}}=\pi_k^{-1}$: for $k\ne k^{-1}$ extend
one of the two maps and invert it, and for an involution $k$ the partial map is an involution on its domain, which we extend by the identity.

The bath is $r_*+|K|-1\le3r_*-2$ qubits. On the first $r_*$, indexed by $\mathcal F$, let $V_g$ be the Pauli matrix $Z$ on the qubit $g$, so that
$\tau(V_gV_h)=\delta_{gh}$, and let $P_k$ permute these qubits by $\pi_k$, so that $P_kV_gP_k^*=V_{kg}$ whenever $g,kg\in\mathcal F$. On the other
qubits choose diagonal unitaries $\chi_k$ with $\chi_1=I$, $\chi_{k^{-1}}=\chi_k^*$ and $\tau(\chi_k^*\chi_l)=\delta_{kl}$: for an involution, $Z$ on one
qubit, and for a pair $k\ne k^{-1}$, $\mathrm{diag}(1,i,-1,-i)$ and its adjoint on two qubits, different pairs on different qubits. Put
$R_k=P_k\otimes\chi_k$, so that $R_{k^{-1}}=R_k^*$ and $R_kV_gR_k^*=V_{kg}$ for $g,kg\in\mathcal F$, and put $C_{g,h}=V_gR_{hg^{-1}}^*$. These are
orthonormal for the normalized trace: if $hg^{-1}\ne h'g'^{-1}$ the factor $\chi$ has trace zero, and otherwise
$\tau(C_{g,h}^*C_{g',h'})=\tau(R_kV_gV_{g'}R_k^*)=\delta_{gg'}$, which then forces $h=h'$. In the same way
$\tau(C_{g,h}C_{g',h'}^*)=\tau(V_gV_{g'})=\delta_{gg'}$ when $hg^{-1}=h'g'^{-1}$, and it vanishes otherwise.

Let $W=\sum A_gA_h^*\otimes C_{g,h}$, summed over the pairs with $hg^{-1}\in K$; the other products $A_gA_h^*$ vanish. On the block of a symbol
with label $g$, including the first block, where $g=1$, it is $V_g$. On a triangle block, with $H$ the Hadamard matrix and using
$R_xV_1R_x^*=V_x$ and $R_xV_yR_x^*=V_z$,
\[
W=\frac12\begin{pmatrix}V_1+V_y&(V_1-V_y)R_x^*\\ R_x(V_1-V_y)&V_x+V_z\end{pmatrix}
=\begin{pmatrix}I&0\\0&R_x\end{pmatrix}(H\otimes I)\begin{pmatrix}V_1&0\\0&V_y\end{pmatrix}(H\otimes I)\begin{pmatrix}I&0\\0&R_x^*\end{pmatrix},
\]
a self-adjoint unitary; triangles with coinciding labels are covered by the same computation. So $W$ is unitary, and by orthonormality of the
$C_{g,h}$, $(\mathrm{id}\otimes\tau)[W(X\otimes I)W^*]=\sum_{g,h}A_gA_h^*XA_hA_g^*=\Phi_G\Phi_G^*(X)$. For $\Phi_G^*\Phi_G$ use right translations
$g\mapsto gk$ with $k\in\{1\}\cup\{y^{\pm1}\}$, the operators $C_{g,h}=V_gR_{g^{-1}h}^*$, and rows instead of columns; on a triangle block the rows carry
$\{1,y\}$ and $\{x,z\}$, and the same computation applies with $x$ and $y$ exchanged. For $\Theta^2$, one further maximally mixed qubit dephases the
two blocks, and controlled on the block the two factorizations act on a common bath. Refreshing the whole bath in every round gives the services.
\end{proof}

These services refresh the whole bath in every round, so their exchange is far above the rank rate; whether $\Phi\Phi^*$ has services with $P=C=0$ and
memory polynomial in $\log(n/\epsilon)$ at its own rank rate we do not know.

So a self-running system with a tracial bath cannot reach expense through $\Phi\Phi^*$. It can through a lazy version of the doubling.

\paragraph{A lazy doubling.} The doubled channel $\Theta$ of Remark~\ref{rem:doubling} is expensive but not positive, and its square is cheap.
Averaging $\Theta$ with the identity restores positivity and keeps the expense. Put
\[
\Theta_+=\tfrac12(\mathrm{id}+\Theta),\qquad
\Theta_+\begin{pmatrix}x&y\\u&w\end{pmatrix}=\frac12\begin{pmatrix}x+\Phi(w)&y\\u&w+\Phi^*(x)\end{pmatrix}\qquad\text{on }M_{1746}.
\]
In Remark~\ref{rem:doubling} only $zz^*=1$ and $\tau(z^2)=0$ are used, so $z$ may also be the unitary $\mathrm{diag}(1,i)$ of a qubit with its
normalized trace. Since $\Theta$ is a unital channel that is self-adjoint for the Hilbert--Schmidt inner product, its spectrum lies in $[-1,1]$ (unital channels are contractions for the Hilbert--Schmidt norm), and
$\Theta_+$ is self-adjoint and positive semidefinite, as Proposition~\ref{prop:autonomous} requires of an autonomous limit.

\begin{proposition}[An expensive autonomous limit]\label{prop:lazy}
Let $\Phi$, $V$ and $\Theta$ be as in Remark~\ref{rem:doubling}, and $\kappa=6+\log107$.
\begin{enumerate}[label=(\alph*),itemsep=1pt]
\item \emph{Autonomy.} Let $Z$ be the coordinate of $L^\infty[-1,1]$ with the normalized Lebesgue measure, put
$\mathcal B=L(\Gamma\times\Z_3)\,\overline\otimes\,L^\infty[-1,1]$ with the product trace, and
$W=\exp\bigl(i\tfrac\pi2V\otimes Z\bigr)=\cos(\tfrac\pi2Z)+i\sin(\tfrac\pi2Z)\,V\in M_{1746}\otimes\mathcal B$. The autonomous dynamics
$\alpha=\mathrm{Ad}\,W$ has $\Phi_t=\Theta_+$ for every $t\ge1$. More generally, $\exp(isV\otimes Z)$ gives the channel
$\Xi_s=\frac{1+\mathrm{sinc}\,2s}2\,\mathrm{id}+\frac{1-\mathrm{sinc}\,2s}2\,\Theta$, with $\mathrm{sinc}\,x=\sin x/x$, and
$\ddia(\Xi_s,\Theta_+)=|\sin2s|/(4s)$.
\item \emph{Memory against purity, lower bound.} There are $c,c_0>0$ and $n_0$ such that for $n\ge n_0$ and $\log n\le B\le c_0n$ every causal
service of $\Theta_+$ for $n$ rounds with error at most $1/16$ and purity $P+C\le B$ has, at any exchange,
\[
Q\ \ge\ c\,\frac n{B\,(\log n)^{2\kappa}} .
\]
In particular, for every $A_0\ge0$ there is $c'>0$ such that services with $P+C\le A_0\log n$ have $Q\ge c'n/(\log n)^{2\kappa+1}$ for all large $n$.
\item \emph{Upper bound.} There is $K$ such that for every $n\ge1$ and $0<B\le n$, $\Theta_+$ has an exact causal service with purity $P+C<B$,
$J=(n-1)Q$ and $Q\le K(n/B)\log(2n/B)$. Moreover $\Theta_+\in\cK_{1746}$.
\end{enumerate}
\end{proposition}

So $\Theta_+$ is a long-time limit, indeed an exact finite-time value, of an autonomous dynamics with a tracial bath; it is expensive in the
sense defined at the beginning of this section, and it trades memory and purity at the same fixed rate as $\Phi$ (Theorem~\ref{thm:main-channel}(a)). Nothing in the construction is local. Part (b) loses nothing against
Theorem~\ref{thm:main-channel}(a) because the round extracted from a service of $\Theta_+$ is converted into a round of $\Phi$ with a
\emph{linear} loss in leakage and in the commutator $\|[U,1\otimes\sqrt\rho\,]\|_F^2$, which replaces entropy production in the one-round theorem
(Lemmas~\ref{lem:commutator-form} and~\ref{lem:block-extraction}), on the same bath.

\begin{proof}[Proof of (a)]
Since $V^2=1$ and $Z$ commutes with $V$, $\exp(isV\otimes Z)=\cos(sZ)+i\sin(sZ)V$, so that $W^t=\exp(it\frac\pi2V\otimes Z)$ and
\[
W^t(X\otimes1)W^{-t}=\cos^2(\tfrac{t\pi}2Z)\,X+\sin^2(\tfrac{t\pi}2Z)\,V(X\otimes1)V+\tfrac i2\sin(t\pi Z)\bigl(V(X\otimes1)-(X\otimes1)V\bigr).
\]
The trace of $\mathcal B$ is the product of the trace of $L(\Gamma\times\Z_3)$ and the integral over $Z$. The last term integrates to $0$ because
$\sin(t\pi Z)$ is odd, $\int\cos^2(\frac{t\pi}2Z)=\frac12(1+\mathrm{sinc}\,t\pi)=\frac12$ for $t\ge1$, and $(\mathrm{id}\otimes\tau)[V(X\otimes1)V]=\Theta(X)$.
So $\Phi_t=\frac12(\mathrm{id}+\Theta)$. The same computation with $s$ in place of $\frac{t\pi}2$ gives $\Xi_s$. Finally
$\Xi_s-\Theta_+=\frac12\mathrm{sinc}(2s)(\mathrm{id}-\Theta)$, and $\|\mathrm{id}-\Theta\|_\diamond=2$: the input $|0\rangle\langle0|\otimes\sigma$ is mapped
by $\Theta$ into the block $|1\rangle\langle1|$. The automorphism $\mathrm{Ad}\,W$ is inner, hence trace preserving.
\end{proof}

The same channel is reached, for every $t\ge1$, by a Bernoulli-shift construction like that of Proposition~\ref{prop:nonlocal}: with
$\mathcal B=L(\Gamma\times\Z_3)\,\overline\otimes\,\overline\bigotimes_{k\in\Z}(M_2,\tau_2)$, $\beta$ the shift of the chain and $p_0$ the projection $|1\rangle\langle1|$ of the qubit at site $0$, the automorphism
$\mathrm{Ad}\bigl((1-p_0)+p_0V\bigr)\circ(\mathrm{id}\otimes\beta)$ raises $V$ to the parity of $t$ fair bits.

We now work with a general unital channel $\Phi$ on $M_d$, its doubling $\Theta(\begin{smallmatrix}x&y\\u&w\end{smallmatrix})=\mathrm{diag}(\Phi(w),\Phi^*(x))$
and $\Theta_+=\frac12(\mathrm{id}+\Theta)$ on $M_{2d}$. Let $\mathcal S_\Phi\subseteq M_d$ be the span of the Kraus operators of $\Phi$, so that the Choi
support of $\Phi$ is $\{|A\rangle\!\rangle:A\in\mathcal S_\Phi\}$, and let $\Pi_\Phi$ also denote the Hilbert--Schmidt projection of $M_d$ onto
$\mathcal S_\Phi$. Since $\Theta(X)=\sum_g(e_{01}\otimes A_g)X(e_{01}\otimes A_g)^*+(e_{10}\otimes A_g^*)X(e_{10}\otimes A_g^*)^*$ for Kraus operators
$A_g$ of $\Phi$, the Choi support of $\Theta_+$ corresponds to the span of the three pairwise orthogonal subspaces
\begin{equation}\label{eq:lazy-support}
\mathbb C\,1_{2d},\qquad e_{01}\otimes\mathcal S_\Phi,\qquad e_{10}\otimes\mathcal S_\Phi^*
\end{equation}
of $M_{2d}=M_2\otimes M_d$. Throughout, $\|\cdot\|_F$ is the Hilbert--Schmidt norm and $\|\cdot\|$ the operator norm, as in
Section~\ref{sec:weighted-proof}; an operator $X$ on the bath $\mathbb C^M$ is identified with $1\otimes X$ when it acts on $\mathbb C^d\otimes\mathbb C^M$,
so that $\|1\otimes X\|_F=\sqrt d\,\|X\|_F$.

\begin{lemma}[Commutator form of the one-round theorem]\label{lem:commutator-form}
Lemma~\ref{lem:weighted-extraction}, Theorem~\ref{thm:weighted} and Corollary~\ref{cor:amplification} remain true if, in the definition
$\nu=\ell+\Delta S$, the entropy production $\Delta S$ of the round $\Psi(X)=\mathrm{Tr}_M\,U(X\otimes\rho)U^*$ is replaced by any $\tilde s\ge0$ with
$\|[U,1\otimes\sqrt\rho\,]\|_F^2\le4d\,\tilde s\ln2$, and if the hypothesis $\ddia(\Psi,\Phi_G)\le1/8$ is replaced by the weaker hypothesis
\begin{equation}\label{eq:coherence}
|\langle z|\Psi(|z\rangle\langle0|)|0\rangle|\le\tfrac34,
\end{equation}
where $0$ and $z$ are the anchor positions of the empty word and of $j$; in Lemma~\ref{lem:weighted-extraction}(iii) the bound then reads
$E(\phi(a)\phi(b))\ge0.66-e$. Conversely, every round can be rewritten, without changing its channel or its leakage, as a round
$(U',\rho)$ with $U'=(1\otimes W)U$ for which $\tilde s=\Delta S$ is admissible.
\end{lemma}

\begin{proof}
The rewriting is the normalization step of the proof of Lemma~\ref{lem:weighted-extraction}, after which~\eqref{eq:block-commutators} is
the inequality $\|[U',1\otimes\sqrt\rho\,]\|_F^2\le4d\,\Delta S\ln2$; its proof uses only that $U'$ is a unitary on $\mathbb C^d\otimes\mathbb C^M$. In
the proof of Lemma~\ref{lem:weighted-extraction}, $\Delta S$ enters only through~\eqref{eq:block-commutators}: afterwards only the bounds
$\epsilon_S\le2\sqrt{\ln2}\,\eta$, $\epsilon_\ell+\epsilon_S\le\sqrt{1+4\ln2}\,\eta$ and $2\epsilon_\ell^2+2\epsilon_S^2\le8\ln2\,\eta^2$ with $\eta=\sqrt{d\nu}$ are used,
together with the unitarity of $U$, its leakage and $\ddia(\Psi,\Phi_G)$. These bounds hold with $\epsilon_S^2\le4d\tilde s\ln2$ and $\nu=\ell+\tilde s$. The diamond hypothesis enters only in the displacement step, through
$|\Tr(\rho B_\varnothing^*B_j)|=|\langle z|\Psi(|z\rangle\langle0|)|0\rangle|\le2u\le\frac14$. Under~\eqref{eq:coherence} instead,
$E(\phi(j))^2\ge2-\frac32-4\epsilon_\ell\ge\frac12-4\eta$, so $E(\phi(a)\phi(b))\ge\sqrt{1/2-4\eta}-e\ge0.66-e$ if $\eta\le\frac1{64}$, while $e>1$ otherwise. In
the proof of Theorem~\ref{thm:weighted} this gives $E(Z_i)\ge0.66-\delta_L$ and $\Delta_i\ge0.66-30\delta_L\ge0.66-b_*>\frac12$ in place of
$\Delta_i\ge1-b_*$, and $\Delta_i\ge\frac12$ is all that Corollary~\ref{cor:sector} needs. Theorem~\ref{thm:weighted} and Corollary~\ref{cor:amplification} use
the round only through Lemma~\ref{lem:weighted-extraction}.
\end{proof}

\begin{lemma}[Weighted intertwining]\label{lem:twisted}
Let $B$ be an operator on $\mathbb C^d\otimes\mathbb C^M$ with $\|B\|\le1$, let $R,R'$ be self-adjoint operators on $\mathbb C^M$, $\xi$ an operator on $\mathbb C^M$, and
$\Sigma=BR-R'B$.
\begin{enumerate}[label=(\roman*),itemsep=1pt]
\item If $h$ is Lipschitz with constant $L_h$ on an interval containing the spectrum of $R$, then $\|[h(R),\xi]\|_F\le L_h\,\|[R,\xi]\|_F$.
\item If $g(x)=\int\hat g(t)e^{itx}dt$ with $M_1=\int|t\,\hat g(t)|\,dt$ and $M_2=\int t^2|\hat g(t)|\,dt$ finite, then
\[
\bigl\|\bigl(Bg(R)-g(R')B\bigr)\xi\bigr\|_F\ \le\ M_1\,\|\Sigma\xi\|_F+\tfrac12M_2\,\|\Sigma\|\,\sqrt d\,\|[R,\xi]\|_F .
\]
\end{enumerate}
\end{lemma}

\begin{proof}
(i) In an eigenbasis of $R$ with eigenvalues $r_k$, the entries of $[h(R),\xi]$ and $[R,\xi]$ are $(h(r_k)-h(r_l))\xi_{kl}$ and $(r_k-r_l)\xi_{kl}$.
(ii) Differentiating $\theta\mapsto e^{i\theta tR'}Be^{i(1-\theta)tR}$ gives
$Be^{itR}-e^{itR'}B=it\int_0^1e^{i\theta tR'}\,\Sigma\,e^{i(1-\theta)tR}\,d\theta$. Write $e^{i(1-\theta)tR}\xi=\xi e^{i(1-\theta)tR}+[e^{i(1-\theta)tR},\xi]$ and use
(i) with $h(x)=e^{i(1-\theta)tx}$; unitaries do not change Hilbert--Schmidt norms, so
$\|(Be^{itR}-e^{itR'}B)\xi\|_F\le|t|\,\|\Sigma\xi\|_F+\frac{t^2}2\|\Sigma\|\sqrt d\,\|[R,\xi]\|_F$. Integrate against $|\hat g(t)|\,dt$.
\end{proof}

We fix the cutoff once and for all. Let $S:\mathbb R\to[0,1]$ be smooth with $S=0$ on $(-\infty,0]$ and $S=1$ on $[1,\infty)$, let $0<s_0\le\frac14$, and put
\[
g(x)=\sin\bigl(\tfrac\pi2S(2x/s_0-1)\bigr)\,S(3-2x),\qquad g_2(x)=\cos\bigl(\tfrac\pi2S(2x/s_0-1)\bigr),\qquad f(x)=\sqrt x\,g(x).
\]
Then $g$ is smooth with compact support, so $M_g:=\int(|t|+t^2)|\hat g(t)|\,dt<\infty$; on $[0,1]$ we have $0\le g\le1$, $g=0$ on $[0,s_0/2]$,
$g=1$ and $g_2=0$ on $[s_0,1]$, and $g^2+g_2^2=1$. Let $L_g$ be a common Lipschitz constant of $g$ and $g_2$ on $[0,1]$, put $\sigma=\sqrt{2/s_0}$ and
\[
K_a=11M_g+1+1.5L_g+2\sigma(4+1.5L_g).
\]
These depend only on $s_0$.

\begin{lemma}[Block extraction]\label{lem:block-extraction}
Let $\Phi$ be a unital channel on $M_d$, and let $\Psi(X)=\mathrm{Tr}_M\,U(X\otimes\rho)U^*$ be a round of $\Theta_+$, with $U$ unitary on
$\mathbb C^2\otimes\mathbb C^d\otimes\mathbb C^M$, $u=\ddia(\Psi,\Theta_+)$ and leakage $\ell$ relative to $\Theta_+$. Put
$\epsilon_\ell=\sqrt{2d\ell}$, $\epsilon_S=\|[U,1\otimes\sqrt\rho\,]\|_F$, $\gamma=\epsilon_\ell+\epsilon_S$ and $\mathcal E=2u+s_0+4K_a\sqrt d\,\gamma$. If $\mathcal E<\frac12$,
there are a unitary $U'$ on $\mathbb C^d\otimes\mathbb C^M$ and a state $\rho'$ on $\mathbb C^M$ such that the round
$\Psi'(X)=\mathrm{Tr}_M\,U'(X\otimes\rho')U'^*$ of $\Phi$, with leakage $\ell_\Phi(\Psi')$ relative to $\Phi$, satisfies
\[
\ddia(\Psi',\Phi)\le\frac{\mathcal E}{\frac12-\mathcal E},\qquad \ell_\Phi(\Psi')\le\frac{4K_a^2\gamma^2}{d(\frac12-\mathcal E)},\qquad
\bigl\|[U',1\otimes\sqrt{\rho'}\,]\bigr\|_F^2\le\frac{8K_a^2\gamma^2}{\frac12-\mathcal E},
\]
and $|\langle x|\Psi'(|x\rangle\langle y|)|y\rangle|\le \mathcal E/(\frac12-\mathcal E)$ for all unit vectors $x,y$ with $\langle x|\Phi(|x\rangle\langle y|)|y\rangle=0$.
\end{lemma}

\begin{proof}
Put $\xi=\sqrt\rho$ and write $U=\sum_{a,b\in\{0,1\}}e_{ab}\otimes U_{ab}$ with $D_0=U_{00}$, $B=U_{01}$, $D_1=U_{11}$, operators on $\mathbb C^d\otimes\mathbb C^M$.

\emph{Residuals.} As in the proof of Lemma~\ref{lem:weighted-extraction}, $\mathcal J(\Psi)$ is the marginal of the unit vector
$(2d)^{-1/2}U\xi\in M_{2d}\otimes M_M$, so $\ell=(2d)^{-1}\|(\Pi^\perp\otimes\mathrm{id})(U\xi)\|_F^2$, with $\Pi$ the projection of $M_{2d}$ onto the span
of~\eqref{eq:lazy-support}. Put $\alpha=(2d)^{-1}(\Tr\otimes\mathrm{id})(U)$, the average of the $2d$ diagonal entries of $U$ in $M_M$, and
$E_a=D_a-\alpha$, $F=(\Pi_\Phi^\perp\otimes\mathrm{id})(B)$. By~\eqref{eq:lazy-support}, $(\Pi^\perp\otimes\mathrm{id})(U)$ has the diagonal blocks $E_0,E_1$ and
the upper right block $F$; since $\Pi^\perp\otimes\mathrm{id}$ commutes with right multiplication by $\xi$, the numbers $e_a=\|E_a\xi\|_F$ and
$e_F=\|F\xi\|_F$ are at most $\epsilon_\ell$. Also $\|\alpha\|\le1$ and $\|E_a\|\le2$.

\emph{Commutators.} Put $k_X=\|[X,\xi]\|_F$ for $X\in\{D_0,B,D_1\}$; these are at most $\epsilon_S$. The diagonal entries of $U$ are entries of $D_0$ and
$D_1$, so by Cauchy--Schwarz $k_\alpha:=\sqrt d\,\|[\alpha,\xi]\|_F\le(k_{D_0}^2+k_{D_1}^2)^{1/2}/\sqrt2\le\epsilon_S/\sqrt2$. With $R=1-\alpha^*\alpha$ and
$R'=1-\alpha\alpha^*$, both between $0$ and $1$, $k_R:=\sqrt d\,\|[R,\xi]\|_F\le2k_\alpha\le\sqrt2\,\epsilon_S$, and likewise $k_{R'}\le\sqrt2\,\epsilon_S$.

\emph{Defects.} Unitarity gives $B^*B=1-D_1^*D_1$ and $BB^*=1-D_0D_0^*$, hence
\begin{gather*}
\Delta:=R-B^*B=\alpha^*E_1+E_1^*D_1,\qquad\Delta':=R'-BB^*=\alpha E_0^*+E_0D_0^*,\\
\Sigma:=BR-R'B=B\Delta-\Delta'B,
\end{gather*}
with $\|\Delta\|,\|\Delta'\|\le1$ and $\|\Sigma\|\le2$. From $\xi E_1=E_1\xi-[D_1,\xi]+[\alpha,\xi]$ and $D_1\xi=\xi D_1+[D_1,\xi]$,
$\|\Delta\xi\|_F\le e_1+\|\xi E_1\|_F+2k_{D_1}\le2e_1+3k_{D_1}+k_\alpha\le4\gamma$; symmetrically $\|\Delta'\xi\|_F\le2e_0+3k_{D_0}+k_\alpha\le4\gamma$;
and $\Sigma\xi=B\Delta\xi-\Delta'\xi B-\Delta'[B,\xi]$ gives $\|\Sigma\xi\|_F\le9\gamma$.

\emph{Two factorizations.} Let $P$ and $P'$ be the spectral projections of $R$ and $R'$ for $(s_0/2,\infty)$. Since $g$ vanishes on $[0,s_0/2]$,
$g(R)=R^{-1/2}Pf(R)$ with $\|R^{-1/2}P\|\le\sigma$, and likewise for $R'$. The operators $B_1=BR^{-1/2}P$ and $B_1'=R'^{-1/2}P'B$ satisfy
\[
B_1^*B_1=P-\tilde\Delta,\quad \tilde\Delta=PR^{-1/2}\Delta R^{-1/2}P,\qquad B_1'B_1'^*=P'-\tilde\Delta',\quad\tilde\Delta'=P'R'^{-1/2}\Delta'R'^{-1/2}P'.
\]
Take polar decompositions $B_1=V_1|B_1|$ and $B_1'=|B_1'^*|V_1'$ with $V_1,V_1'$ unitary. On the range of $1\otimes P$ we have
$|B_1|=1+\varphi(\tilde\Delta)$ with $\varphi(x)=\sqrt{1-x}-1$ and $\tilde\Delta\le1$, and $\varphi(x)^2\le x^2$; so $\|(|B_1|-P)Y\|_F\le\|\tilde\Delta Y\|_F$
for every $Y$ with range in that subspace, and similarly for $B_1'$. Put $Z_0=Bg(R)\xi=B_1f(R)\xi$. Then
$Z_0-V_1f(R)\xi=V_1(|B_1|-P)f(R)\xi$, and $\tilde\Delta f(R)\xi=PR^{-1/2}\Delta g(R)\xi$, while
$\|\Delta g(R)\xi\|_F\le\|\Delta\xi\|_F+\sqrt d\,\|[g(R),\xi]\|_F\le\|\Delta\xi\|_F+L_g\,k_R$ by Lemma~\ref{lem:twisted}(i). Hence
\[
a_1:=\|Z_0-V_1f(R)\xi\|_F\le\sigma(4+1.5L_g)\gamma .
\]
On the other side,
$Z_0=(Bg(R)-g(R')B)\xi+g(R')[B,\xi]+[g(R'),\xi]B+\xi f(R')B_1'$ and $\xi f(R')B_1'=\xi f(R')V_1'+\xi f(R')(|B_1'^*|-P')V_1'$, where
$\|\xi f(R')(|B_1'^*|-P')\|_F\le\|\tilde\Delta'f(R')\xi\|_F\le\sigma(\|\Delta'\xi\|_F+L_g\,k_{R'})$. By Lemma~\ref{lem:twisted}(ii),
$\|(Bg(R)-g(R')B)\xi\|_F\le M_g(\|\Sigma\xi\|_F+k_R)\le11M_g\gamma$. Hence
\[
a_2:=\|Z_0-\xi f(R')V_1'\|_F\le11M_g\gamma+(1+1.5L_g)\gamma+\sigma(4+1.5L_g)\gamma,\qquad a_1+a_2\le K_a\gamma .
\]

\emph{Polar decomposition of $Z_0$.} Put $\zeta_1=|f(R)\xi|=(\xi f(R)^2\xi)^{1/2}$ and $\zeta_2=|(\xi f(R'))^*|=(\xi f(R')^2\xi)^{1/2}$. Then $|V_1f(R)\xi|=\zeta_1$ and $|(\xi f(R')V_1')^*|=\zeta_2$, as operators $1\otimes\zeta_i$. By the
Araki--Yamagami inequality $\||X|-|Y|\|_F\le\sqrt2\,\|X-Y\|_F$~\cite{ArakiYamagami1981}, applied to $Z_0$ and to $Z_0^*$,
$\||Z_0|-\zeta_1\|_F\le\sqrt2\,a_1$ and $\||Z_0^*|-\zeta_2\|_F\le\sqrt2\,a_2$. Let $Z_0=O|Z_0|=|Z_0^*|O$ with $O$ unitary. Then
\[
\|O\zeta_1-\zeta_2O\|_F\le\|O(\zeta_1-|Z_0|)\|_F+\|(|Z_0^*|-\zeta_2)O\|_F\le\sqrt2(a_1+a_2)=:\delta .
\]

\emph{The round.} The Hermitian operators $O(1\otimes\zeta_1)O^*$ and $1\otimes\zeta_2$ are within $\delta$ in $\|\cdot\|_F$, and their spectra are those of $\zeta_1$
and $\zeta_2$, each repeated $d$ times. By the Hoffman--Wielandt inequality~\cite{HoffmanWielandt1953}, the decreasing eigenvalue lists
$(x_k)$ of $\zeta_1$ and $(y_k)$ of $\zeta_2$ satisfy $d\sum_k(x_k-y_k)^2\le\delta^2$; so the unitary $W_h$ of $\mathbb C^M$ that maps an eigenbasis of
$\zeta_1$ to one of $\zeta_2$ in the same order has $\|1\otimes(W_h\zeta_1W_h^*-\zeta_2)\|_F\le\delta$. Put $U'=(1\otimes W_h^*)O$, $c_1=\Tr\zeta_1^2=\Tr\rho f(R)^2$
and $\rho'=\zeta_1^2/c_1$. Then $\|[U',1\otimes\zeta_1]\|_F\le\|O\zeta_1-\zeta_2O\|_F+\|(\zeta_2-W_h\zeta_1W_h^*)O\|_F\le2\delta$, which is the third bound once we
know $c_1\ge\frac12-\mathcal E$. The channel of $(U',\rho')$ is $\Psi'=\Phi_{O\zeta_1}/c_1$, where $\Phi_X(w)=\mathrm{Tr}_M[X(w\otimes1)X^*]$. Its leakage is
$(dc_1)^{-1}\|(\Pi_\Phi^\perp\otimes\mathrm{id})(O\zeta_1)\|_F^2$, and
\[
\|(\Pi_\Phi^\perp\otimes\mathrm{id})(O\zeta_1)\|_F\le\|(\Pi_\Phi^\perp\otimes\mathrm{id})(Z_0)\|_F+\sqrt2\,a_1\le e_F+L_g\,k_R+\sqrt2\,a_1\le2K_a\gamma,
\]
since $(\Pi_\Phi^\perp\otimes\mathrm{id})(Bg(R)\xi)=F\xi\,g(R)+(\Pi_\Phi^\perp\otimes\mathrm{id})(B[g(R),\xi])$. This is the second bound.

\emph{Closeness.} For a completely positive map $\Phi_X$ as above, $\|\Phi_X\|_\diamond=\|\mathrm{Tr}_M(X^*X)\|$, and
$\|\Phi_X-\Phi_Y\|_\diamond\le\|X-Y\|_F\bigl(\|\mathrm{Tr}_M X^*X\|^{1/2}+\|\mathrm{Tr}_M Y^*Y\|^{1/2}\bigr)$, where $\mathrm{Tr}_M$ traces out the input bath.
Let $Y=B\xi$. The upper left block of $\Psi(e_{11}\otimes w)$ is $\Phi_Y(w)$ and that of $\Theta_+(e_{11}\otimes w)$ is $\frac12\Phi(w)$; compressing to
the block, $\|\Phi_Y-\frac12\Phi\|_\diamond\le2u$. Since $g^2+g_2^2=1$ on $[0,1]$, $\Phi_Y-\Phi_{Z_0}-\Phi_{Bg_2(R)\xi}$ is
$w\mapsto\mathrm{Tr}_M[B(w\otimes X)B^*]$ with $X=\sum_{h\in\{g,g_2\}}h(R)[h(R),\rho]$, and $\|X\|_1\le4L_g\|[R,\xi]\|_F$, so its diamond norm is at most
$4L_g\,k_R/\sqrt d$ (split $X$ into positive and negative parts and use $B^*B\le1$). Next,
$\|\Phi_{Bg_2(R)\xi}\|_\diamond\le\Tr\rho\,g_2(R)Rg_2(R)+\|\xi g_2(R)\Delta g_2(R)\xi\|_1\le s_0+\sqrt d(\|\Delta\xi\|_F+L_g\,k_R)$, because $xg_2(x)^2\le s_0$ on
$[0,1]$. Finally $\mathrm{Tr}_M\zeta_1^2=c_1\le1$ and $\mathrm{Tr}_M Z_0^*Z_0\le1$ give $\|\Phi_{O\zeta_1}-\Phi_{Z_0}\|_\diamond\le2\sqrt2\,a_1$. Adding up,
\[
\bigl\|\Phi_{O\zeta_1}-\tfrac12\Phi\bigr\|_\diamond\le2u+s_0+2\sqrt2\,a_1+\sqrt d(4+1.5L_g)\gamma+6L_g\gamma/\sqrt d\le \mathcal E .
\]
Evaluating at $\tau_d$ gives $|c_1-\frac12|\le \mathcal E$, so $c_1\ge\frac12-\mathcal E$, and $\|\Psi'-\Phi\|_\diamond\le \mathcal E/c_1+|1/(2c_1)-1|\le2\mathcal E/c_1$. If
$\langle x|\Phi(|x\rangle\langle y|)|y\rangle=0$, then $|\langle x|\Phi_{O\zeta_1}(|x\rangle\langle y|)|y\rangle|\le\|\Phi_{O\zeta_1}-\frac12\Phi\|_\diamond\le \mathcal E$, and dividing by
$c_1$ gives the last bound.
\end{proof}

\begin{proof}[Proof of (b)]
Let a service of $\Theta_+$ for $n$ rounds have error $\epsilon\le\frac1{16}$ and $P+C\le B$ with $\log n\le B$. Theorem~\ref{thm:extraction},
applied to $\Theta_+$ on $M_{1746}$, gives a round with bath $\mathbb C^M$, $M=2^Q$, $u=\ddia(\Psi_t,\Theta_+)\le\frac1{15}$ and
$\nu_t=\ell_t+\Delta S_t\le t_n\le\frac{16}{15}\cdot\frac{B+1+\log n}n\le\frac{3.2\,B}n$, as in the proof of Theorem~\ref{thm:main-channel}. The normalization
step of the proof of Lemma~\ref{lem:weighted-extraction}, which uses only that $U_t$ is a unitary on $\mathbb C^{2d}\otimes\mathbb C^M$, replaces $U_t$ by
$(1\otimes W)U_t$ without changing $\Psi_t$ or its leakage, after which $\epsilon_S^2=\|[U_t,1\otimes\sqrt{\rho_t}\,]\|_F^2\le4(2d)\Delta S_t\ln2$; so
$\gamma^2\le2(2d\ell_t+8d\ln2\,\Delta S_t)\le12d\,\nu_t$. Apply Lemma~\ref{lem:block-extraction} with $s_0=\frac1{100}$ and $d=873$:
$\mathcal E\le\frac2{15}+\frac1{100}+14K_ad\sqrt{\nu_t}\le\frac3{14}$ as soon as $14K_ad\sqrt{3.2B/n}\le0.07$. Then $c_1\ge\frac12-\mathcal E\ge\frac27$, and the anchors $0$
and $z=z(j)$ satisfy $\langle z|\Phi(|z\rangle\langle0|)|0\rangle=\tau(\lambda([j]))=0$, so the round $\Psi'$ satisfies~\eqref{eq:coherence}. By
Lemma~\ref{lem:commutator-form} it may be used in Corollary~\ref{cor:amplification} with
\[
\nu'=\ell_\Phi(\Psi')+\frac{\|[U',1\otimes\sqrt{\rho'}]\|_F^2}{4d\ln2}\le\Bigl(4+\frac2{\ln2}\Bigr)\frac{K_a^2\gamma^2}{d\,c_1}\le290K_a^2\,\nu_t\le928K_a^2\,\frac Bn .
\]
Let $C_R$ and $y_0$ be the constants of Corollary~\ref{cor:amplification}(a) for $R_\Phi$ and the transport of Section~\ref{sec:lamp-channel}, and put
$y_n=n/(928K_a^2C_RB)$, so that $C_R\nu'y_n\le1$ and $y_n\le n$. With
$c_0=\min\{1/(928K_a^2C_Ry_0),\ (0.07/(14K_ad))^2/3.2\}$ and $B\le c_0n$, the corollary gives
\[
Q=\log M\ \ge\ S(\rho')\ \ge\ \frac{y_n}{32(\log y_n)^{2\kappa}}\ \ge\ \frac1{29696\,K_a^2C_R}\cdot\frac n{B(\log n)^{2\kappa}} .
\]
For the second statement take $B=\max(1,A_0)\log n$.
\end{proof}

\begin{proof}[Proof of (c)]
Let $E$, $K_G$ and a $(1-r^{-1})$-expansive $r^{-1}$-morphism $f:E\to\Sym(\Omega)$ be as in Proposition~\ref{prop:sofic-upper}, with its unitary
$U_f=\sum_ac_aE_a\otimes P_{\pi_a}$ and its set $\Omega_0$ of good points, at which $\mathrm{Tr}_\Omega U_f(w\otimes|\omega\rangle\langle\omega|)U_f^*=\Phi(w)$. Let
$\Omega_1$ be the set of $\omega$ with $(f([x])f([y]))^{-1}(\omega)=f([z])^{-1}(\omega)$ for every triangle $(x,y,z)$ and $f(g)^{-1}(\omega)\ne f(h)^{-1}(\omega)$ for
distinct labels $g,h$. As $U_f^*(\psi\otimes|\omega\rangle)=\sum_a\overline{c_a}E_a^*\psi\otimes|\pi_a^{-1}(\omega)\rangle$, the computation of
Proposition~\ref{prop:sofic-upper} gives $\mathrm{Tr}_\Omega U_f^*(x\otimes|\omega\rangle\langle\omega|)U_f=\Phi^*(x)$ for $\omega\in\Omega_1$. A point outside $\Omega_1$
is the image under the bijection $f([z])$ of a point where $f([x])f([y])$ and $f([z])$ differ, or under $f(g)$ of a point where $f(g)$ and $f(h)$
agree; so $|\Omega\setminus\Omega_1|\le K_GM/r$, and $\Omega_*=\Omega_0\cap\Omega_1$ has $|\Omega_*|\ge(1-2K_G/r)M$.

Let the bath be a coin qubit, a phase qubit carrying $z=\mathrm{diag}(1,i)$, and $\mathbb C^\Omega$, and put
$W_f=|0\rangle\langle0|\otimes1+|1\rangle\langle1|\otimes V_f$ with $V_f=\bigl(\begin{smallmatrix}0&U_f\otimes z\\U_f^*\otimes z^*&0\end{smallmatrix}\bigr)$. For every
state $\beta$ that is a mixture of the $|\omega\rangle\langle\omega|$, $\omega\in\Omega_*$, the bath state $\tau_2\otimes\tau_2\otimes\beta$ makes $W_f$ apply
exactly $\Theta_+$: the coin splits the channel into $\frac12\mathrm{id}$ and $\frac12$ the channel of $V_f$, the phase qubit removes the off-diagonal
blocks since $\tau(z^2)=0$, and the diagonal blocks are $\Phi(w)$ and $\Phi^*(x)$ at the points of $\Omega_*$. Now repeat the exact service of
Proposition~\ref{prop:sofic-upper}(c) with $K_G$ replaced by $2K_G$ and two more qubits, which are imported maximally mixed: with
$\eta=B/(4n)$ and $r=2K_G/\eta$ the imported register is maximally mixed on $\mathrm{span}\{|\omega\rangle:\omega\in\Omega_*\}\otimes\mathbb C^{k_0}$ (and on the
two qubits), its deficit is below $B/n$, every round applies exactly $\Theta_+$, and $Q=\lceil\log\sigma_E(8K_Gn/B)+\log(4n/B)\rceil+2$. The bound
$\log\sigma_E(R)\le C_ER\log R$ of Theorem~\ref{thm:main-group}(b) gives $Q\le K(n/B)\log(2n/B)$. With $\beta$ maximally mixed on all of $\Omega$,
the same unitary gives a finite tracial factorization of $(1-\vartheta)\Theta_++\vartheta\,\Upsilon$ for some channel $\Upsilon$ and some $\vartheta\le2K_G/r$, for every $r$; hence
$\Theta_+\in\cK_{1746}$.
\end{proof}

\begin{remark}[Exact finite factorizations of $\Theta_+$]\label{rem:lazy-exact}
$\Theta_+$ is an exact autonomous value but not an exact finite factor. Indeed, if $\Theta_+(X)=(\mathrm{id}\otimes\tau_m)[W(X\otimes1)W^*]$ for a unitary
$W$ on $\mathbb C^{2d}\otimes\mathbb C^m$, the proof of Lemma~\ref{lem:block-extraction} with $\rho=\tau_m$ has $\ell=0$ and $\epsilon_S=0$: the diagonal blocks of
$W$ are both $1\otimes\alpha$, its upper right block $B$ lies in $\mathcal S_\Phi\otimes M_m$, $B^*B=1\otimes R$ and $BB^*=1\otimes R'$. Then
$B(1\otimes R)=(1\otimes R')B$, so for every eigenvalue $\mu>0$ of $R$ the operator $\mu^{-1/2}B$ maps $\mathbb C^d\otimes\ker(R-\mu)$ unitarily onto
$\mathbb C^d\otimes\ker(R'-\mu)$, and the upper left block of $\Theta_+(e_{11}\otimes w)$ writes $\Phi$ as the convex combination
$\sum_\mu\frac{2\mu m_\mu}m\,\Phi_\mu$ of exact factors with baths of dimension $m_\mu=\dim\ker(R-\mu)\le m$. For the channel of
Theorem~\ref{thm:main-channel} this is impossible: such a combination is a single round with $\ell=\Delta S=0$ (a direct sum of flat baths), and
Theorem~\ref{thm:weighted} with $\nu=0$ would give it entropy at least $N/16$ for every $N$.
\end{remark}

\begin{remark}[A lower exchange]\label{rem:lazy-exchange}
At an anchor $z$ of the gadget, $\Phi(|z\rangle\langle z|)=|z\rangle\langle z|$, so $\Theta_+$ maps $|1\rangle\langle1|\otimes|z\rangle\langle z|$ to the flat state
$\frac12(|0\rangle\langle0|+|1\rangle\langle1|)\otimes|z\rangle\langle z|$ of rank two. Proposition~\ref{prop:purity-budget} with $r=2$ therefore gives $P+C<3Q+6$ for
services of $\Theta_+$ with error at most $1/16$ and exchange $J\le n/2$, and Proposition~\ref{prop:lazy}(b), used as in the proof of
Theorem~\ref{thm:main-channel}(b) with $B=\max(3Q+6,\log n)$ (if $B>c_0n$, then already $Q\ge(c_0n-6)/3$), gives $Q\ge c\sqrt n/(\log n)^\kappa$ for such services, whatever their purity.
\end{remark}

So autonomous dynamics with a tracial bath can produce expensive channels. The reduced dynamics of the probe matters: if the channels
$\Phi_t$ themselves form a semigroup of unital channels, they are mixed unitary for all large $t$~\cite{BhatDevendra2025}, hence cheap, whereas
$\Theta_+^2\ne\Theta_+$. Locality constrains how $\Theta_+$ can be approached: Proposition~\ref{prop:light-cone} and
Proposition~\ref{prop:lazy}(b), combined as in the proof of Corollary~\ref{cor:preparation}, show that a local reservoir with $d=1746$ that
produces a channel within $\delta\le\delta_0$ of $\Theta_+$ at time $T$ has $T\ge c\,\delta^{-1/\nu}(\log(1/\delta))^{-(2\kappa+1)/\nu}$, for constants
$c,\delta_0>0$ that depend only on the parameters of the reservoir. So no local reservoir produces $\Theta_+$ exactly at a finite time, or approaches
it exponentially fast. Whether a local dynamics can produce an expensive channel at all is open (Section~\ref{sec:open}).

%======================================================================
\section{Open problems}\label{sec:open}
%======================================================================

\begin{enumerate}[label=(\arabic*)]
\item \emph{The exact profile.} Theorem~A gives $\exp(r^{1+o(1)})$, with polylogarithmic factors in the exponent between
$c\,r/(\log r)^{\kappa}$, $\kappa=6+\log107$, and $C\,r\log r$. The exponent $\kappa$ comes from the Dehn function of $H_3$, known to be polynomial of
degree at most $\log107+4$ and at least $2$~\cite{DouglasMghirbiA}; a quadratic Dehn function would give $\kappa=4$. What is the correct
polylogarithmic factor; is $\log\sigma_E(r)=\Theta(r\log r)$?
\item \emph{Beyond the counting ceiling.} Does some amenable group have a chunk with profile $\exp(c\,r)$? Does any sofic group have a chunk with
profile $\exp(r^{1+c})$ for some $c>0$? The counting criterion cannot give either; a different lower-bound mechanism would be needed. For the
second question, grids of fixed dimension should not help: the argument sketched in Remark~\ref{rem:grid} would bound the profile of their
chunks by $\exp(O(r\log r))$. By
Remark~\ref{rem:sandwich}, a gadget channel of a sofic group can need superlinear memory at zero purity only if its chunk profile exceeds
$\exp(O(r))$. Every fixed chunk of a sofic group has a computable profile~\cite[Lemma~1]{Ivanov2021} (search the degrees $1,2,\dots$ and the finitely
many maps at each), so no chunk outgrows every computable function; how fast a computable profile can grow is open.
\item \emph{The optimal exchange rate in general.} For $\Phi$ the least memory at the rank rate and the least memory with purity at most the
memory are both $n^{1/2+o(1)}$ (Remark~\ref{rem:sqrt-limit}). For every channel with a flat probe the second is at most three times the first, plus
six. Is the converse true up to polylogarithmic factors for every channel in $\cK_d$ with a flat probe? Theorem~\ref{thm:streaming} bounds the
rank-rate memory by $\log k+n\log(k'/k)+O(\log n)$, so the converse holds whenever some exact rectangular dilation makes this at most a
polylogarithmic factor times the least memory at purity at most memory.
\item \emph{The dyadic group.} Its profile is finite and at least $\exp(c\,r^{1/3})$. In $\mathcal G$ the words $w_0=1$,
$w_{2n}=uw_n$ and $w_{2n+1}=tuw_n$ reach the integer $n$ with length $O(\log n)$, and commutators of copies at points of equal parity reduce
freely to smaller indices. What is the area of the commutators at points of opposite parity?
\item \emph{Thompson's groups.} Is $\mathcal L$ sofic whenever $F$ is? The hyperlinearity transfer of Proposition~\ref{prop:thompson-hyperlinear}
does not obviously give permutations, and not every action of a sofic group is sofic: some coset actions of residually finite Kazhdan groups
are not~\cite{KunThom2026}. We do not know whether the action of $F$ on $\mathbb D$ is sofic when $F$ is. Does $\Gamma$ embed in $V$ or
in $2V$? The embedding of Proposition~\ref{prop:3V} uses one coordinate for each ray.
\item \emph{Local autonomous physics.} Proposition~\ref{prop:lazy} gives an expensive long-time limit of an autonomous dynamics with a tracial
bath, but the dynamics is not local. Can an expensive channel be the long-time limit of a fixed translation-invariant local rule, such as a
quantum cellular automaton or a time-independent finite-range Hamiltonian, with the bath in its tracial state? By Proposition~\ref{prop:light-cone}
such a limit would be approached slowly; for $\Theta_+$ the distance at time $t$ is at least of order $t^{-\nu}$ up to logarithmic factors
(Section~\ref{sec:physics}). By~\cite{BhatDevendra2025} the reduced dynamics of the probe could not be a semigroup.
\end{enumerate}


\section*{Use of AI}
\addcontentsline{toc}{section}{Use of AI}

This work was carried out with extensive use of generative AI systems. They proposed constructions and proof strategies, wrote and
checked proofs, carried out computations, searched the literature, wrote the verification code, and drafted and edited the text. The
authors directed the research, assessed and selected the results, and take full responsibility for the content. The proofs are given in
full in the text; repeated adversarial checking and the accompanying programs support them but do not replace them. The verification
notes state which statements are checked by machine.

\begin{thebibliography}{99}
\bibitem{AlekseevBradford2026} V.~Alekseev and H.~Bradford, Sofic actions, halo products, and metric approximations of groups, \href{https://arxiv.org/abs/2601.18742}{arXiv:2601.18742} (2026).
\bibitem{AD2006} C.~Anantharaman-Delaroche, On ergodic theorems for free group actions on noncommutative spaces, Probab. Theory Related Fields \textbf{135} (2006), 520--546.
\bibitem{ArakiYamagami1981} H.~Araki and S.~Yamagami, An inequality for Hilbert--Schmidt norm, Comm. Math. Phys. \textbf{81} (1981), 89--96.
\bibitem{BhatDevendra2025} B.~V.~R.~Bhat and R.~Devendra, On regions of mixed unitarity for semigroups of unital quantum channels, \href{https://arxiv.org/abs/2512.23598}{arXiv:2512.23598} (2025).
\bibitem{Bradford2022} H.~Bradford, Controlling LEF growth in some group extensions, \href{https://arxiv.org/abs/2201.05052}{arXiv:2201.05052} (2022).
\bibitem{BHV2006} S.~Bravyi, M.~B.~Hastings and F.~Verstraete, Lieb--Robinson bounds and the generation of correlations and topological quantum order, Phys. Rev. Lett. \textbf{97} (2006), 050401.
\bibitem{Bridson2002} M.~R.~Bridson, The geometry of the word problem, in: M.~R.~Bridson and S.~M.~Salamon (eds.), Invitations to Geometry and Topology, Oxf. Grad. Texts Math. \textbf{7}, Oxford University Press, Oxford, 2002, 29--91.
\bibitem{Brin2004} M.~G.~Brin, Higher dimensional Thompson groups, Geom. Dedicata \textbf{108} (2004), 163--192.
\bibitem{Brin2005} M.~G.~Brin, Presentations of higher dimensional Thompson groups, J. Algebra \textbf{284} (2005), 520--558; arXiv:math/0501082.
\bibitem{Brin2010} M.~G.~Brin, On the baker's map and the simplicity of the higher dimensional Thompson groups $nV$, Publ. Mat. \textbf{54} (2010), 433--439; arXiv:0904.2624.
\bibitem{Brown1987} K.~S.~Brown, Finiteness properties of groups, J. Pure Appl. Algebra \textbf{44} (1987), 45--75.
\bibitem{BCMR2016} J.~Burillo, S.~Cleary, A.~Martino and C.~E.~R\"over, Commensurations and metric properties of Houghton's groups, Pacific J. Math. \textbf{285} (2016), 289--301; arXiv:1403.0026.
\bibitem{BCS2001} J.~Burillo, S.~Cleary and M.~I.~Stein, Metrics and embeddings of generalizations of Thompson's group $F$, Trans. Amer. Math. Soc. \textbf{353} (2001), 1677--1689; arXiv:math/9809185.
\bibitem{CFP1996} J.~W.~Cannon, W.~J.~Floyd and W.~R.~Parry, Introductory notes on Richard Thompson's groups, Enseign. Math. (2) \textbf{42} (1996), 215--256.
\bibitem{CapraroLupini2015} V.~Capraro and M.~Lupini, Introduction to Sofic and Hyperlinear Groups and Connes' Embedding Conjecture, Lecture Notes in Math. \textbf{2136}, Springer, Cham, 2015.
\bibitem{CarterWegman1979} J.~L.~Carter and M.~N.~Wegman, Universal classes of hash functions, J. Comput. System Sci. \textbf{18} (1979), 143--154.
\bibitem{Chiribella2009} G.~Chiribella, G.~M.~D'Ariano and P.~Perinotti, Theoretical framework for quantum networks, Phys. Rev. A \textbf{80} (2009), 022339; arXiv:0904.4483.
\bibitem{Chou1980} C.~Chou, Elementary amenable groups, Illinois J. Math. \textbf{24} (1980), 396--407.
\bibitem{Cornulier2006} Y.~de Cornulier, Finitely presented wreath products and double coset decompositions, Geom. Dedicata \textbf{122} (2006), 89--108; arXiv:math/0509090.
\bibitem{Cornulier2013} Y.~Cornulier, Sofic profile and computability of Cremona groups, Michigan Math. J. \textbf{62}(4) (2013), 823--841; \href{https://doi.org/10.1307/mmj/1387226167}{doi:10.1307/mmj/1387226167}. Numbered references use \href{https://arxiv.org/abs/1305.0993v1}{arXiv:1305.0993v1}.
\bibitem{DouglasMghirbiA} S.~Douglas and N.~Mghirbi, Houghton's group $H_3$ has superpolynomial sofic profile, version 1.1.2, Zenodo (2026); \href{https://doi.org/10.5281/zenodo.23070869}{doi:10.5281/zenodo.23070869}. Section numbers refer to version 1.1.2.
\bibitem{DouglasMghirbiB} S.~Douglas and N.~Mghirbi, Causal quantum-channel simulation: memory beyond entropy, version 1.1.2, Zenodo (2026); \href{https://doi.org/10.5281/zenodo.23070870}{doi:10.5281/zenodo.23070870}. Section numbers refer to version 1.1.2.
\bibitem{GKEP2025} D.~Gao, S.~Kunnawalkam Elayavalli and G.~Patchell, Soficity for group actions on sets and applications, Res. Math. Sci. \textbf{12} (2025), 48; arXiv:2401.04945.
\bibitem{Genevois2019} A.~Genevois, Embeddings into Thompson's groups from quasi-median geometry, Groups Geom. Dyn. \textbf{13} (2019), 1457--1510; \href{https://arxiv.org/abs/1709.03888v2}{arXiv:1709.03888v2}. Numbered references use this arXiv version.
\bibitem{GenevoisTessera2024} A.~Genevois and R.~Tessera, Lamplighter-like geometry of groups, \href{https://arxiv.org/abs/2401.13520}{arXiv:2401.13520} (2024).
\bibitem{Guba2006} V.~S.~Guba, The Dehn function of Richard Thompson's group $F$ is quadratic, Invent. Math. \textbf{163} (2006), 313--342.
\bibitem{GutoskiWatrous2007} G.~Gutoski and J.~Watrous, Toward a general theory of quantum games, Proceedings of the 39th Annual ACM Symposium on Theory of Computing (STOC 2007), 565--574; arXiv:quant-ph/0611234.
\bibitem{HaagerupMusat2011} U.~Haagerup and M.~Musat, Factorization and dilation problems for completely positive maps on von Neumann algebras, Comm. Math. Phys. \textbf{303} (2011), 555--594; arXiv:1009.0778.
\bibitem{HaagerupMusat2015} U.~Haagerup and M.~Musat, An asymptotic property of factorizable completely positive maps and the Connes embedding problem, Comm. Math. Phys. \textbf{338} (2015), 721--752; arXiv:1408.6476.
\bibitem{Harper1964} L.~H.~Harper, Optimal assignments of numbers to vertices, J. Soc. Indust. Appl. Math. \textbf{12} (1964), 131--135.
\bibitem{HayesSale2018} B.~Hayes and A.~W.~Sale, Metric approximations of wreath products, Ann. Inst. Fourier \textbf{68} (2018), 423--455; arXiv:1608.02610.
\bibitem{HennigMatucci2012} J.~Hennig and F.~Matucci, Presentations for the higher-dimensional Thompson groups $nV$, Pacific J. Math. \textbf{257} (2012), 53--74; arXiv:1105.3714.
\bibitem{HoffmanWielandt1953} A.~J.~Hoffman and H.~W.~Wielandt, The variation of the spectrum of a normal matrix, Duke Math. J. \textbf{20} (1953), 37--39.
\bibitem{HHO2003} M.~Horodecki, P.~Horodecki and J.~Oppenheim, Reversible transformations from pure to mixed states and the unique measure of information, Phys. Rev. A \textbf{67} (2003), 062104.
\bibitem{Houghton1978} C.~H.~Houghton, The first cohomology of a group with permutation module coefficients, Arch. Math. \textbf{31} (1978), 254--258.
\bibitem{Ivanov2021} A.~Ivanov, Sofic profiles of $S(\omega)$ and computability, Arch. Math. Logic \textbf{60} (2021), 477--494; arXiv:2001.02648.
\bibitem{JNVWY2020} Z.~Ji, A.~Natarajan, T.~Vidick, J.~Wright and H.~Yuen, MIP$^*$ = RE, \href{https://arxiv.org/abs/2001.04383}{arXiv:2001.04383} (2020).
\bibitem{Johnson1999} D.~L.~Johnson, Embedding some recursively presented groups, in: Groups St Andrews 1997 in Bath, II, London Math. Soc. Lecture Note Ser. \textbf{261}, Cambridge Univ. Press, Cambridge, 1999, 410--416.
\bibitem{JungeLeMerdyXu2006} M.~Junge, C.~Le Merdy and Q.~Xu, $H^\infty$ functional calculus and square functions on noncommutative $L^p$-spaces, Ast\'erisque \textbf{305} (2006), vi+138 pp.
\bibitem{Krengel1985} U.~Krengel, Ergodic Theorems, de Gruyter Stud. Math. \textbf{6}, Walter de Gruyter, Berlin, 1985.
\bibitem{KretschmannWerner2005} D.~Kretschmann and R.~F.~Werner, Quantum channels with memory, Phys. Rev. A \textbf{72} (2005), 062323; arXiv:quant-ph/0502106.
\bibitem{Kummerer1985} B.~K\"ummerer, Markov dilations on W$^*$-algebras, J. Funct. Anal. \textbf{63} (1985), 139--177.
\bibitem{KunThom2026} G.~Kun and A.~Thom, Nonsofic wreath products of residually finite groups, \href{https://arxiv.org/abs/2608.06222}{arXiv:2608.06222} (2026).
\bibitem{Lee2012} S.~R.~Lee, Geometry of Houghton's groups, PhD thesis, University of Oklahoma (2012); \href{https://arxiv.org/abs/1212.0257v1}{arXiv:1212.0257v1}. Numbered references use this arXiv version.
\bibitem{LiebRobinson1972} E.~H.~Lieb and D.~W.~Robinson, The finite group velocity of quantum spin systems, Comm. Math. Phys. \textbf{28} (1972), 251--257.
\bibitem{LyndonSchupp1977} R.~C.~Lyndon and P.~E.~Schupp, Combinatorial Group Theory, Ergeb. Math. Grenzgeb. \textbf{89}, Springer, Berlin, 1977.
\bibitem{Milnor1971} J.~Milnor, Introduction to Algebraic $K$-Theory, Ann. of Math. Stud. \textbf{72}, Princeton Univ. Press, Princeton, 1971.
\bibitem{Musat2025} M.~Musat, The Connes--Kirchberg problem and infinite-dimensional phenomena in quantum information theory, \href{https://arxiv.org/abs/2512.00432}{arXiv:2512.00432} (2025).
\bibitem{MusatRordam2020} M.~Musat and M.~R{\o}rdam, Non-closure of quantum correlation matrices and factorizable channels that require infinite dimensional ancilla (with an appendix by N.~Ozawa), Comm. Math. Phys. \textbf{375} (2020), 1761--1776; arXiv:1806.10242.
\bibitem{NSY2019} B.~Nachtergaele, R.~Sims and A.~Young, Quasi-locality bounds for quantum lattice systems. Part I. Lieb--Robinson bounds, quasi-local maps, and spectral flow automorphisms, J. Math. Phys. \textbf{60} (2019), 061101; arXiv:1810.02428.
\bibitem{Rover1999} C.~E.~R\"over, Subgroups of finitely presented simple groups, PhD thesis, University of Oxford (1999).
\bibitem{RybarZiman2008} T.~Ryb\'ar and M.~Ziman, Repeatable quantum memory channels, Phys. Rev. A \textbf{78} (2008), 052114; arXiv:0808.3851.
\bibitem{Slofstra2018} W.~Slofstra, A group with at least subexponential hyperlinear profile, \href{https://arxiv.org/abs/1806.05267v1}{arXiv:1806.05267v1} (2018).
\bibitem{VershikGordon1997} A.~M.~Vershik and E.~I.~Gordon, Groups that are locally embeddable in the class of finite groups, Algebra i Analiz \textbf{9}(1) (1997), 71--97; English transl. St. Petersburg Math. J. \textbf{9}(1) (1998), 49--67.
\end{thebibliography}


\end{document}
